\documentclass[12pt]{article}
\usepackage[utf8]{inputenc}
\usepackage[T2A,T1]{fontenc}
\usepackage{amssymb}
\usepackage{amsmath}
\usepackage{csquotes}

\usepackage[english,russian]{babel}
\usepackage[backend=biber, style=ieee]{biblatex}

\textwidth 14cm
\textheight 21cm

\addbibresource{article_ru.bib}

\title{Несчётность множества континуум, разбор}
\author{Роман Олегович Коротков}

\begin{document}
 
\selectlanguage{russian}

\begin{flushleft}
{\bf УДК 510.223 }\\{\bf MSC 03E99 }
\end{flushleft}

{\let\newpage\relax\maketitle}

\begin{abstract}

Статья посвящена ретроспективному исследованию кризиса оснований математики, связанного с интерпретацией понятия бесконечности. Проведённый в предыдущих работах серии анализ классических доказательств несчётности множества вещественных чисел показал их существенную зависимость от исходных онтологических и методологических представлений о морфологии бесконечного. При допущении универсальной тотальной структуры, обеспечивающей сравнение бесконечных объектов, эти доказательства остаются корректными; при предположении фундаментальной гомогенности всех бесконечных объектов они оказываются внутренне неконсистентными вследствие смешения позиций референта. Детальный разбор существующих аксиоматических теорий множеств (ZFC, NBG и их вариантов) выявил фундаментальные проблемы: ошибки преждевременного обобщения (акциденции и ложной дихотомии) при доказательстве рефлексивности, минимальности $\omega$ и равенства $\bigcup \omega = \omega$ архитектурные дефекты, связанные с неудачной формой аксиомы бесконечности (нестратифицированность, неоднозначность, отсутствие чистоты и независимости); ошибку реификации; и, в конечном итоге, устойчивость условной непротиворечивости, природа которой неясна: либо явные и неявные предположения современной интерпретации действительно не противоречат друг другу, либо аксиомы цензуры выделяют непротиворечивый фрагмент противоречивой теории. Современная интерпретация фактически отрицает исходную теорему Кантора для упорядоченных множеств, распространяя внутреннюю (кардинальную) индукцию на предельные объекты и приводя к равенствам вида $\omega = 2^\omega$, что конфликтует с аксиомой непрерывности и блокирует естественные механизмы усиления предикативности для противостояния семантическим петлям. Авторы выделяют пять ключевых парадигм математики (формалистскую, интуиционистскую, логицистскую, конструктивистскую, структурную) и анализируют их по шести критериям: истина, тождество, выразимость объектов, логическая замкнутость, континуум, метатеория. Этот анализ выявляет существенные различия в допущениях, определяющих границы применимости каждой парадигмы, и объясняет, почему представители разных школ часто говорят на разных языках, даже обсуждая одни и те же теоремы, - но в то же время обнаруживает незаменимость каждой из них для математики.

Вводятся понятия позиций референта, областей регулярности и сингулярности, кардинальной, ординальной и индексальной морфологий, на основе чего формулируется единый концептуальный язык описания бесконечности, уточняющий неоднозначность традиционных понятий (потенциальная/актуальная бесконечность, динамичность/статичность, интенсиональность/экстенсиональность). Через призму разработанного понятийного аппарата проводится анализ истории становления современной интерпретации теории множеств - от оригинальной онтологии Кантора через упрощённые адаптации начала XX века и истоки формирования ZFC/NBG, прорывные работы фон Неймана - до закрепления доминирующей картины, сопровождавшегося постепенной догматизацией и институциональной цензурой альтернативных подходов. Предлагается теоретическая конструкция, объединяющая альтернативные концепции бесконечного с современной интерпретацией в непротиворечивую картину. Вводятся аксиомы стеснения и безопасной индукции, предлагается расслоение отношения принадлежности и аксиомы индексальной и кардинальной бесконечности, позволяющие корректно переключать перспективы и восстановить согласованность с оригинальной онтологией Кантора. Работа демонстрирует, что многие парадоксы самореференции возникают из смешения позиций референта, а не из самой самореференции. В заключении обозначаются ключевые направления дальнейшего развития: построение консервативных расширений (ZFCS, NBGS, ZFCL, NBGL), их синтез с сюрреальными числами, новыми основаниями Куайна, теорией гросс-множеств и нестандартным анализом; исследование диверсификации моделей аннотации континуума и её следствий для теории мер; применение аппарата позиций референта к семантическим теоремам (Гёделя, Тарского, Клини, Тьюринга) и проблеме искусственного интеллекта; разработка формальных критериев выявления цензуры объектов в аксиоматических системах; интеграция когнитивно-математических и историко-философских методов для преодоления институциональной стагнации в основаниях математики. Предложенная объединяющая картина позволяет свести существенное количество разрозненных областей в единую непротиворечивую систему и открывает путь к методологически честной работе с бесконечностью.

{\it Ключевые слова: несчётность вещественных чисел; диагональный аргумент; теорема Кантора; морфология бесконечности; аксиоматическая теория множеств; ZFC; NBG; ординалы; кардиналы; индексалы; позиция референта; потенциальная и актуальная бесконечность; предикативность и импредикативность; парадигмы математики (формализм, интуиционизм, логицизм, конструктивизм, структурализм); аксиома бесконечности; аксиома стеснения; аксиома безопасной индукции; парадоксы самореференции; теоремы Гёделя и Тарского; континуум-гипотеза; сюрреальные числа; новые основания Куайна; гросс-множества; когнитивная математика; основания математики; история теории множеств.}

\end{abstract}

\selectlanguage{english}

\begin{abstract}

The paper is devoted to a retrospective investigation into the crisis in the foundations of mathematics associated with the interpretation of the concept of infinity. An analysis of classical proofs of the uncountability of the set of real numbers, carried out in previous works of the series, revealed their substantial dependence on initial ontological and methodological conceptions of the morphology of the infinite. Under the assumption of a universal total structure that enables the comparison of infinite objects, these proofs remain valid; under the assumption of a fundamental homogeneity of all infinite objects, they prove internally inconsistent due to a conflation of referent positions. A detailed examination of existing axiomatic set theories (ZFC, NBG, and their variants) uncovers fundamental problems: fallacies of hasty generalization (accident and false dichotomy) in proving reflexivity, the minimality of $\omega$, and the equality $\bigcup \omega = \omega$; architectural defects associated with an ill-formed axiom of infinity (lack of stratification, ambiguity, and lack of purity and independence); the fallacy of reification; and, ultimately, the persistence of conditional consistency, whose nature remains unclear—either the explicit and implicit assumptions of the modern interpretation genuinely do not contradict one another, or censorship axioms isolate a consistent fragment of a contradictory theory. The modern interpretation effectively negates Cantor’s original theorem for ordered sets by extending internal (cardinal) induction to limit objects and leading to equalities such as $\omega = 2^\omega$, which conflicts with the axiom of continuity and blocks natural mechanisms for strengthening predicativity to resist semantic loops. The authors identify five key mathematical paradigms—formalist, intuitionist, logicist, constructivist, and structural—and analyze them across six criteria: truth, identity, expressibility of objects, logical closure, continuum, and metatheory. This analysis reveals significant differences in the assumptions defining the boundaries of applicability of each paradigm and explains why representatives of different schools often speak different languages even when discussing the same theorems, while at the same time showing that each of them is indispensable to mathematics.

The notions of referent positions, regions of regularity and singularity, and cardinal, ordinal, and indexal morphologies are introduced, serving as the basis for formulating a unified conceptual language for describing infinity that clarifies the ambiguity of traditional concepts (potential/actual infinity, dynamism/staticism, intensionality/extensionality). Through the lens of this conceptual apparatus, the paper traces the history of how the modern interpretation of set theory took shape—from Cantor’s original ontology, through simplified early-twentieth-century adaptations and the origins of ZFC/NBG, to von Neumann’s breakthrough contributions and the entrenchment of the dominant picture, accompanied by gradual dogmatization and institutional censorship of alternative approaches. A theoretical construction is proposed that unifies alternative conceptions of infinity with the modern interpretation into a consistent picture. Axioms of limitation and safe induction are introduced, along with a stratification of the membership relation and axioms of indexal and cardinal infinity, enabling proper switching of perspectives and restoring alignment with Cantor’s original ontology. The work demonstrates that many paradoxes of self-reference arise from conflating referent positions rather than from self-reference itself. In conclusion, key directions for further research are outlined: constructing conservative extensions (ZFCS, NBGS, ZFCL, NBGL) and synthesizing them with surreal numbers, Quine’s New Foundations, gross-set theory, and non-standard analysis; investigating the diversification of continuum annotation models and its consequences for measure theory; applying the apparatus of referent positions to semantic theorems (Gödel, Tarski, Kleene, Turing) and the problem of artificial intelligence; developing formal criteria for detecting object censorship in axiomatic systems; and integrating cognitive-mathematical and historical-philosophical methods to overcome institutional stagnation in the foundations of mathematics. The proposed unifying picture brings together a substantial number of disparate areas into a single consistent system and opens a path toward methodologically honest work with infinity.

{\it Keywords: uncountability of reals; diagonal argument; Cantor’s theorem; morphology of infinity; axiomatic set theory; ZFC; NBG; ordinals; cardinals; indexals; referent position; potential and actual infinity; predicativity and impredicativity; paradigms of mathematics (formalism, intuitionism, logicism, constructivism, structuralism); axiom of infinity; axiom of limitation; axiom of safe induction; paradoxes of self-reference; Gödel’s and Tarski’s theorems; continuum hypothesis; surreal numbers; Quine’s New Foundations; gross-set theory; cognitive mathematics; foundations of mathematics; history of set theory.}

\end{abstract}

\selectlanguage{russian}

Бесконечность принадлежит к числу наиболее противоречивых и вместе с тем фундаментальных понятий человеческой мысли. Будучи глубоко укоренённой в структуре мышления, она остаётся абстракцией, не имеющей единого, исчерпывающего определения, и в то же время активно и зачастую продуктивно используется в различных областях науки. Любая попытка дать строгое описание или целостную характеристику этого понятия неизбежно сталкивается с его фундаментальной трансцендентностью и концептуальной неуловимостью. Тем примечательнее, что одна из самых строгих дисциплин - математика - не только включила бесконечность в число своих объектов, но и поместила её в основание теории множеств, при этом не фиксируя явно, какое именно содержание вкладывается в данное понятие. Однако, независимо от того, мыслим ли мы бесконечность как сущность, предикат, начало, категорию, монаду или иным образом, речь идёт об одном и том же феномене. Из этого следует, что существенные свойства и характеристики, выявленные в рамках одного подхода, должны, при корректном рассуждении, воспроизводиться и в других. Различные описания бесконечности, таким образом, представляют собой не конкурирующие онтологии, а взаимодополняющие перспективы, применимые к разным областям и задачам. Расхождения в получаемых результатах в таком случае выступают надёжным индикатором ошибки в рассуждении. Отсутствие бесконечности в сфере непосредственного чувственного опыта существенно затрудняет выработку точного и однозначного языка её описания. Тем не менее, ключ к её осмыслению лежит не столько в аксиомах и формальных системах, сколько в самом языке. Именно язык превращает неопределённое явление в концептуальную сущность, задавая рамки возможного мышления о нём.

В работах ''Несчётность множества континуум, шах'', ''Несчётность множества континуум, ещё шах'' и ''Несчётность множества континуум, шах и мат'' было выявлено примечательное явление: доказательства несчётности множества вещественных чисел существенным образом зависят от исходных представлений о морфологии бесконечного. При допущении существования универсальной и тотальной структуры, позволяющей сравнивать бесконечные объекты между собой, данные доказательства остаются корректными. Напротив, при предположении фундаментальной гомогенности всех бесконечных объектов доказательства становятся некорректными вследствие внутренней неконсистентности используемых в них конструкций. Иными словами, ключевое наблюдение Георга Кантора, положившее начало систематическому изучению структуры бесконечного, оказывается критически зависимым от предположения, что бесконечность уже обладает определённой структурой, допускающей унификацию и сравнение. Попытка проанализировать существующие аксиоматические теории множеств с целью объяснить, почему столь сильные ограничения на применимость одного из центральных методов этих теорий остаются в них практически незаметными, привела к выявлению ряда фундаментальных проблем. Во-первых, непротиворечивость этих теорий обеспечивается лишь в контексте довольно сильных метаматематических и философских допущений, в частности - возможности отождествления континуума с множеством и допустимости для теории описывать объекты, невыразимые в её собственном языке. Во-вторых, доказательство ключевых для теории свойств, таких как рефлексивность, минимальность $\omega$, $\bigcup \omega = \omega$ и прочих, опирается на существенные методологические ограничения, связанные с ошибками преждевременного обобщения: фиксированный класс внутренних объектов теории считается единственным классом допускающим интерпретацию функционального отношения внутри теории (ошибка акциденции), свидетельство существования индуктивной морфологии в форме аксиомы бесконечности воспринимается как требование индуктивной морфологии для всех бесконечных множеств (ложная дихотомия). В-третьих, принятая форма аксиомы бесконечности представляется неудачной: она не является ни стратифицированной и конструктивной. Не удовлетворяет формальным требованиям к аксиомам: не является атомарной (влечёт отсутсвие максимальных элементов), оказывается неоднозначной (содержит концептуальную неопределённость относительно свойств сингулярности, отделяющей области регулярности), не является чистой (делает бесконечное замыкание объектом теории и задаёт его структуру), не проявляет независимость (в неочевидной комбинации с другими аксиомами выводит рефлексивность, которую вообще говоря стоит постулировать отдельно). При этом данная аксиома предполагает нечто принципиально неверифицируемое. Перечисленные проблемы сами по себе не приводят к формальным противоречиям внутри теории, однако существенно затрудняют формирование ясной и однозначной картины, переводя её из сферы строгой математики в область герменевтики - искусства интерпретации. Метаматематические и методологические допущения носят характер соглашения. Пока они разделяются большинством представителей научного сообщества, такие соглашения можно считать конструктивными. Тем не менее, они нуждаются в явном и честном формулировании - как для чёткого обозначения границ применимости теории, так и ввиду существования альтернативных подходов, развиваемых другими математиками и философами. Альтернативные взгляды способны внести вклад в познание бесконечного, но размытость интерпретаций требует веры, обеспечивающей смещение взглядов к идеологическим, что блокирует возможность продуктивного синтеза с чётким обозначением линий разграничения. Наконец, указанные проблемы не позволяют должным образом осмыслить принципиальное расхождение между оригинальной картиной Кантора и современной аксиоматической теорией множеств. Последняя фактически отрицает теорему Кантора для упорядоченных множеств, распространяя внутреннюю индукцию на точку разрыва предикативности $\omega$ как на объект. У самого Кантора неравенство $\omega < 2^\omega$ является доказуемой теоремой, выражающей невозможность переупорядочить множество порядкового типа $2^\omega$ в порядковый тип $\omega$. Современные же аксиоматики, допуская применимость внутренней индукции к внешним объектам - предельным порядковым типам, - приходят к тому, что $\omega = 2^\omega$, и далее к более высоким порядковым типам, например $\omega^\omega$ допускает переупорядочение в тип $\omega$. Что в принципе может быть допустимо, поскольку порядковые типы представляют собой синтез внутренней и внешней алгебры, а определение предела через супремум высвечивает именно внутреннюю составляющую, но аксиоматики идут дальше не различия ограничений сделанного выбора и насильно связывают алгебраическую структуру таких теорий не различающих внутренней и внешней перспективы с методологически обоснованной картиной Кантора. Такие допущения приводят к побочным эффектам, в частности, вступают в прямое проитворечие с аксиомой непрерывности в контексте математического анализа, что ограничивает их использование и требует довольно синтетических обоснований. Более того, в конечном итоге, они блокируют естественные механизмы противостояния семантическим петлям в современных теориях, и в более содержательном смысле, ведут к условной непротиворечивости или недостоверности. В частности, если в метатеории, рассматривающей арифметику Пеано, мы предполагаем, что системы объектов, маркируемые арифметической и геометрической прогрессиями, образуют завершённые целостности и могут служить областями определения и значений биекции (взаимно однозначного соответствия в сильной интерпретации), то мы сильно рискуем получить противоречивую теорию. В отличие от этого, теория Кантора уже изначально нивелирует опасность самореференции благодаря теореме, носящей его имя, допускающей обобщение строго разводящее внутреннюю и внешнюю составляющую математического описания. Её объекты образуют упорядоченную структуру, которая в современных аксиоматических системах была адаптирована в виде кумулятивной иерархии фон Неймана, более того они допускают введение более детальной индексной структуры с использованием нормальных операций. Неосторожное толкование учения Кантора в рамках современных аксиоматик приводит к двойному эффекту от парадоксов самореференции (синтаксическому и семантическому), двойному иерархическому расслоению (иерархия фон Неймана и иерархия мета-теорий Тарского) и устойчивости условной непротиворечивости. Где под устойчивостью понимается невозможность блокировать данный феномен, заимствованием механизмов доступных в естественном языке. Такая интерпретация существенно усиливает идеологический компонент подхода мейнстрима и, как следствие, приводит к случаям институциональной дисфункции и научной цензуры альтернативных взглядов на природу бесконечного.

В настоящей работе авторы ставят перед собой задачу ретроспективно исследовать причины возникновения и развития кризиса оснований математики. Этот кризис привёл к тому, что математическое сообщество признало корректной лишь одну картину условно непротиворечивой математической реальности, оставив за её пределами другие, не менее легитимные альтернативные подходы и интерпретации. Мы попытаемся сформулировать несколько математических парадигм, выражающих ключевые метаматематические, философские и методологические допущения относительно процесса математического познания. К сожалению, сложившаяся система научных рейтингов и публикационных механизмов способствует преимущественной поддержке исследований, резонирующих с господствующим мейнстримом, в то время как темы, затрагивающие границы применимости и фундаментальные основания теорий, остаются существенно менее освещёнными. Между тем, по мнению авторов, такие работы имеют не меньшее научное значение, чем наиболее цитируемые прорывные исследования, поскольку именно они ограничивают возможности концептуальной цензуры и злоупотребления доминирующим положением в научном сообществе. Авторы считают целесообразным обратиться к сравнительно новому направлению - когнитивной математике, - чтобы продемонстрировать ещё один слой анализа и дополнительный ракурс на обсуждаемые проблемы. По глубокому убеждению авторов, математика представляет собой взгляд вовнутрь: это попытка постичь суть процесса отражения действительности в сознании. Математика выступает универсальным языком описания этого отражения, а эволюция обеспечивает его тесную связь с реальностью, позволяя математике с высокой точностью описывать и её. Именно поэтому современная математика немыслима без когнитивистики, которая обеспечивает прочный эпистемологический фундамент тем концепциям, что в классической традиции апеллируют исключительно к интуиции и часто опираются на прыжки веры. 

Кроме того, авторы предпримут попытку сформировать единый концептуальный язык описания бесконечности, дополнив традиционную терминологию рядом новых понятий. Мы покажем неоднозначность и контекстуальную зависимость традиционных терминов, таких как потенциальная и актуальная бесконечность, динамичность (незавершённость, процесс, поток) и статичность (целостность, сущность, объект), а также интенсиональность и экстенсиональность. Именно к этой неоднозначности адаптировались представления мейнстрима, что существенно затрудняло адекватное позиционирование и восприятие альтернативных взглядов. С помощью этого концептуального языка предполагается продемонстрировать ограничения и неявные метаматематические и философские допущения современных подходов в теории множеств. Далее планируется предложить теоретическую конструкцию, которая объединяет альтернативные концепции бесконечного с основными положениями современной интерпретации математики в единую непротиворечивую картину, сопровождаемую явными разграничительными утверждениями. Для иллюстрации излагаемых положений мы предлагаем простую структурную модель, интерпретирующую современные аксиоматические системы теории множеств. Данная модель отчётливо демонстрирует ключевые концептуальные развилки и необходимость дополнительных вспомогательных предположений в классической интерпретации. Наконец, на основе разработанного языка будет представлен исторический экскурс в развитие теории множеств и противоречий связанных с ней, в ходе которого все основные неопределённости и концептуальные разрывы становятся наглядными и органично интегрируются в предлагаемую конструкцию. В заключение авторы обозначают ключевые направления дальнейшего исследования, развития и уточнения современных взглядов на основания математики в свете предложенной объединяющей картины, которая позволяет свести существенное количество разрозненных областей в единую непротиворечивую систему.

Работая с математическими текстами, невозможно обойти вниманием тот факт, что метаматематический, методологический и философский контекст исследований играет существенную роль в оценке достоверности тех или иных выводов и результатов. К сожалению, большинство математиков крайне редко проговаривают в явном виде фундаментальные выборы, которые определяют границы применимости их теорий. Именно эти неявные допущения часто становятся причиной взаимного непонимания между различными школами, а также источником неожиданных открытий при переходе из одного концептуального контекста в другой. В настоящее время теория категорий и теория топосов фактически выступают в роли универсального языка современной математики, при этом не укладываясь в рамки классических школ обоснования математики - формализма, интуиционизма и логицизма. Мы наблюдаем взрывной рост алгоритмизации математического знания, автоматического доказательства теорем и конструктивной математики, тесно связанных с программированием и верификацией программ. В контексте квантовых вычислений и искусственного интеллекта вопросы вычислимости, непрерывности и самореференции приобретают новое звучание. Эти процессы вновь возвращают актуальность проблемам оснований математики и необходимости строгой фиксации границ теоретических построений. Авторы выделяют пять наиболее влиятельных парадигм математики:

\begin{enumerate}
    \item {\it Формалистская} (Д. Гильберт, А. Тарский, К. Гёдель) - подход, ставящий синтаксическую форму, правила и процедуры выше семантического содержания. Именно эта парадигма превратила математику в строгую символическую игру и легла в основу всех современных систем автоматического доказательства теорем.
    \item {\it Интуиционистская} (Л. Брауэр, А. Гейтинг) - философски наиболее глубокая школа, подчёркивающая роль разума и саморефлексии по отношению к математическому мышлению. Здесь на первое место ставятся мыслеобразы, возникающие в сознании при взаимодействии с реальностью.
    \item {\it Логическая} (Г. Фреге, Б. Рассел, А. Уайтхед) - классический логицизм, стремящийся свести всю математику к чистой логике, где все математические структуры и формы выступают лишь сложными комбинаторными проявлениями базовых бинарных примитивов.
    \item {\it Конструктивистская} (А. Марков, Э. Бишоп) - прагматичный подход, ориентированный на вычислимость и алгоритмическую верифицируемость; семантика рассматривается здесь как процессуальное свойство.
    \item {\it Структурная} (С. Эйленберг, С. Маклейн, С. Лоувер, Дж. Конвей) - доминирующая парадигма современной чистой математики, ключевым представителем которой является теория категорий. Объекты полностью определяются своими отношениями внутри структур, а семантика возникает как эмерджентное свойство этих структур.
\end{enumerate}

Выбор именно этих пяти парадигм обусловлен тем, что они в совокупности охватывают практически весь спектр классических и современных подходов - от максимально формального и технического до наиболее философского и рефлексивного. Другие направления (например, платонизм или финитизм) либо представляют собой частные случаи перечисленных парадигм, либо уже интегрированы в них и могут рассматриваться как их дальнейшая специализация. Для содержательного сравнения парадигм между собой авторы предлагают шесть ключевых критериев, которые, на наш взгляд, наиболее точно выявляют те неявные допущения, которые обычно остаются за кадром:

\begin{enumerate}
    \item {\it Истина} - что именно считается истинным в рамках данной теории и как возникает это понятие (включая природу его противоположности).
    \item {\it Тождество} - когда два объекта считаются одним и тем же, а когда - различными (критерий, тесно связанный с онтологией).
    \item {\it Выразимость объектов} - как интерпретируется существование математических объектов и возможности их выражения на языке теории (эпистемологический аспект).
    \item {\it Логическая замкнутость} - в какой мере теория допускает тотальную замкнутость и универсальную категоризацию своих объектов (закон исключённого третьего).
    \item {\it Континуум} - природа непрерывности и способ её введения в теорию, связь с концепциями импредикативности и бесконечности.
    \item {\it Метатеория} - как теория относится к самой себе: отделена ли она, иерархична или самореферентна.
\end{enumerate}

Эти шесть критериев позволяют компактно и глубоко раскрыть границы каждой парадигм, показать, где она обладает наибольшей силой, а где вынуждена идти на компромиссы, а также объяснить, почему представители разных математических культур часто говорят на разных языках, даже обсуждая одни и те же теоремы. 

Реформа образования, инициированная Вильгельмом фон Гумбольдтом и развитая далее, приводит к вытеснению философского базиса инженерными практиками. Философия постепенно вытесняется специализированными курсами, утрачивая статус основополагающей дисциплины лежащей в сердце науки, и уступает место технической рациональности. Математика выделяется как самостоятельная дисциплина. В это время растёт количество учебных заведения, что, в частности, приводит к быстрому росту числа математиков, вследствие чего возникает потребность в стандартизации языка и обмена опытом. Ответом на потребность сообщества в стандартизации становится возникновение формализма, который часто связывают с именем Давида Гильберта, очень точно уловившего структуру момента. В основе программе Гильберта лежит простая инженерная идея создания песочниц (формальных систем), в которых можно безопасно работать, не боясь столкнуться с противоречиями. Предполагалось, что каждая математическая теория должна быть переведена в такую формальную систему - строгий язык с четкими аксиомами и правилами вывода. Сама формальная система становится в таком случае математическим объектом, который можно изучать, не обращаясь к внутреннему содержанию (семантике) формул, пользуясь лишь правилами их преобразования. Ключевой шаг замыкающий идею песочницы - доказательство непротиворечивости формальной системы, желательно надёжными финитарными методами, что должно гарантировать достоверность утверждений теории. Давайте рассмотрим формализм в разрезе наших шести критериев.

\begin{enumerate}
    \item {\it Истина} - в классическом формализме Гильберта истина отождествляется с непротиворечивостью формальной системы: если аксиомы и все выводимые из них утверждения не приводят к противоречию, то они считаются истинными (в рамках данной системы), а объекты, ими определяемые, - существующими. Иными словами, истина является атрибутом выводимых утверждений синтаксически непротиворечивой системы; синтаксическая непротиворечивость (консистентность) порождает семантическое содержание, включая понятие истинности, в том смысле, что любая непротиворечивая теория первого порядка имеет модель \cite{completeness-theorem}. Однако современный формализм (в широком смысле) вынужден учитывать платоновский взгляд на истину как на идею, независимую от конкретной формальной системы. Даже радикальный формализм (С. Феферман) и ультрафинитизм (Э. Нельсон, Д. Зеильбергер) признают: существуют утверждения, истинные в стандартной модели арифметики (например, гёделевское предложение G), которые не могут быть ни доказаны, ни опровергнуты в формальной системе (например, в арифметике Пеано, PA) при условии её непротиворечивости. Но понятие стандартная модель само требует дополнительных метатеоретических допущений: существования множества натуральных чисел как объекта (не класса, не совокупности, а именно множества) и его бесконечности по Дедекинду. Формально корректнее было бы говорить не о стандартной модели PA, а о стандартном расширении PA - то есть о теории, включающей PA и дополнительные утверждения (например, о теории, интерпретируемой в ZFC с аксиомой бесконечности), поскольку стандартная модель де-факто сильнее PA. Внутри же самого PA Гёделевское предложение G является независимым: обе теории $PA + G$ и $PA + \neg G$ непротиворечивы (при условии непротиворечивости PA, ведь если бы в PA выводилась противоречивость $PA + \neg G$, то PA бы доказывала $G$). Апелляция к стандартной модели без явного указания на используемые метатеоретические предположения скрывает существенное усиление исходной формальной системы и может вводить в заблуждение, что можно объяснить только платоновским признанием стандартной модели натуральных чисел.
    \item {\it Тождество} - в формализме тождество $=$ определяется исключительно синтаксически: как бинарный предикат, подчиняющийся аксиомам рефлексивности, симметричности, транзитивности и подстановки (принцип Лейбница). Семантически $=$ может интерпретироваться и как строгое тождество элементов модели, и как более слабое отношение эквивалентности, являющееся конгруэнцией относительно всех сигнатурных символов. Формальная система не различает эти интерпретации: любая формула, истинная в одной, истинна и в другой. Эта семантическая слепота делает формализм уязвимым для ошибки реификации - отождествления формальной репрезентации объекта с самим объектом (карты и территории). Классический пример - теорема Банаха–Тарского: в ZFC доказывается, что трёхмерный шар (понятый как множество точек) может быть разбит на конечное число частей, которые после изометрий перегруппировываются в два шара того же объёма. Это не логическое противоречие, а демонстрация того, что формальная карта шара (структура описываемая формальной системой) изоморфна карте двух шаров. Противоречие возникает лишь тогда, когда карту принимают за территорию - то есть полагают, что формальное описание исчерпывает природу шара, включая такие выходящие за рамки формальной системы свойства, как быть одним связным объектом. Интересно, что в отличие от двумерного пространства, формальная структура которого остаётся изоморфной самой себе при масштабировании, трехмерное обладает свойством, позволяющим его формальной структуре описывать не только увеличенную копию, но и две разделённых копии самого себя. Это отражает не парадоксальность шара, а лишь то, что формальное описание (множество точек с группой изометрий) допускает неизмеримые разбиения, которые не сохраняют объём и связность - но эти последние не входят в формальную структуру.
    \item {\it Выразимость объектов} - формализм полностью принимает неконструктивное существование объектов: они не обязаны быть вычислимыми, явно построенными, идентифицируемыми (в смысле наличия разрешимого критерия равенства) или иметь хоть какое-либо выражение на языке теории (например, замкнутый терм). Достаточно, чтобы они удовлетворяли аксиомам системы. Например, в теории меры Лебега доказывается, что мера одноточечного множества равна нулю: $\lambda(\{x\}) = 0$ для каждого $x \in \mathbb{R}$. Это универсальное утверждение доказывается для всех точек континуума одновременно, без возможности индивидуальной верификации. Формальная система позволяет делать такие универсальные утверждения, не требуя перебора или явного построения каждого экземпляра. Однако явно сформулировать в языке теории требование, чтобы существующий объект был выразим термом, мы не можем: для этого требуется квантификация по термам, что выходит за рамки логики первого порядка и требует языка второго порядка или метаязыка. В результате, в формализме нет возможности в рамках формальной системы первого порядка отличить ситуацию, когда объект существует и имеет замкнутый терм, от ситуации, когда он существует, но не имеет замкнутого терма (невыразим). Обе ситуации выглядят одинаково с точки зрения доказуемых формул: теория просто утверждает $\exists x \; \varphi(x)$, ничего не говоря о наличии терма. Но в действительности, из-за мета-допущения, которое мы не фиксируем формула должна выглядеть примерно так: $\exists x \wedge expressible(x) : \varphi(x)$. Как мы видели в статье ''Несчётность множество континуум, шах и мат'' данная проблема способна приводить к формальной независимости утверждений, которые кажутся доказуемыми в логике первого порядка.
    \item {\it Логическая замкнутость} - формализм в своей классической форме использует логику с законом исключённого третьего (LEM): $p \vee \neg p$ является тавтологией (аксиомой или выводимым правилом). Это позволяет доказывать утверждения через сведение к абсурду: если $\neg p$ ведёт к противоречию, то $p$ истинно. В современном контексте допускается также работа с более слабыми формальными системами, в которых LEM не принимается. Однако, формализм оказывается слеп к природе парадокса лжеца (данное утверждение ложно). Он не может его объяснить или разрешить - только избежать через ограничение выразительной силы (запрет самореференции), ослабление логики (отказ от LEM) или введение иерархий. Формализм не способен дать позитивное описание того, почему лжец не является корректным предложением в идеальном языке, или почему он ведёт к противоречию. Природа парадокса остаётся интуитивной, мета-формальной (выходящей за пределы формального анализа). Этот эпистемический дефицит проявляется в фундаментальном компромиссе: в рамках одной формальной системы невозможно одновременно обладать тотальной логической замкнутостью (применимостью всех предикатов ко всем объектам) и универсальной категоризацией (способностью представлять любое разделение объектов в виде объекта системы). NF жертвует логической замкнутостью ради универсальной категоризации (существует множество всех множеств и любые семейства множеств заданные стратифицированными формулами). ZFC, напротив, логически замкнута, но не допускает множества всех ординалов или кардиналов. При этом речь идёт не о принципиальной невозможности, а об отсутствии в формализме аппарата для выработки нового языка: формализм всегда оперирует готовыми языковыми конструкциями, не имея механизма рефлексивного выхода за их пределы.
    \item {\it Континуум} - в формализме континуум предстаёт в двух не сводимых друг к другу ипостасях: как аксиоматическая структура - полное упорядоченное архимедово поле (единственное с точностью до изоморфизма), и как теоретико-множественный объект - множество всех подмножеств натуральных чисел $\mathcal{P}(\mathbb{N})$, мощностью $2^{\aleph_0}$. Обычно считается, что именно несчётная мощность ответственна за полноту поля, хотя аксиома полноты является свойством порядка, не ссылающимся на мощность. Возможность построить неархимедовы модели (например, сюрреальные числа) показывает, что архимедовость - это дополнительное условие, связывающее континуум со счётным натуральным рядом. Без архимедовости классические доказательства о несчётности континуума утрачивают силу, а сам континуум может быть реализован в структурах сколь угодно большой мощности (или даже в собственных классах, как в случае сюрреальных чисел). Архимедовость и отождествление континуума с $\mathcal{P}(\mathbb{N})$ приводит к утрате древнегреческой идеи непрерывности как самоподобия: любая часть континуума подобна целому и делима в тех же пропорциях. Последовательное деление отрезка пополам в формализме превращает связный континуум в множество изолированных точек, которое имеет мощность континуума но утрачивает связность и подобие исходному объекту. Таким образом, континуум в формализме утрачивает свою интуитивную природу, сохраняя лишь формальные атрибуты (полнота).
    \item {\it Метатеория} - классический формализм Гильберта предполагал, что непротиворечивость формальных систем можно доказать финитными средствами, которые считались абсолютно надёжными и не нуждающимися в дальнейшем обосновании. Эволюция должна была идти от простого (финитная метатеория языка) к сложному (богатым объектным теориям реальности). Однако развитие логики в XX веке показало обратное: теорема о полноте логики первого порядка требует для своего доказательства средств, выходящих за пределы самой этой логики. Теоремы Гёделя о неполноте устанавливают, что непротиворечивость достаточно сильной теории не может быть доказана внутри неё самой и требует метатеории с отличной интерпретацией объектов. Иерархия Тарского показывает, что понятие истины для формального языка может быть определено только на метаязыке более высокого уровня, и этот процесс не имеет последнего шага. Следствием является принципиальная незавершимость формализма: любая формальная система $T$ существует в контексте метатеории $T'$, которая, в свою очередь, требует своей метатеории $T''$, и так далее. Уже само определение базовых понятий - истинности и доказуемости - требует бесконечной иерархии метатеорий. Не существует абсолютной метатеории, которая могла бы служить окончательным фундаментом для всех формальных систем. Отношение теории к метатеории является рекурсивным и открытым, что исключает возможность самодостаточного обоснования математики в рамках одного формализма.
\end{enumerate}

С современных позиций мы видим, что формализм по своей сути является инженерной практикой и, как любой стандарт, требует органов по сертификации, надзору и развитию. Можно выделить три слоя аппарата бюрократии, развившегося в контексте распространения формализма:

\begin{enumerate}
    \item {\it Стандартизация публикаций}: журналы начинают требовать всё более строгой, аксиоматической подачи. Рецензент превращается в инспектора по качеству, проверяющего, соблюдены ли правила игры.
    \item {\it Сообщества экспертов}: появляется слой математиков, специализирующихся не на содержательных задачах, а на мета-работе - проверке непротиворечивости систем, формализации доказательств.
    \item {\it Учебный процесс}: формальный метод встраивается в систему образования. Студентов учат не рефлексировать о смысле, а оперировать символами по правилам, что делает формализм педагогическим стандартом, который воспроизводит сам себя.
\end{enumerate}

Стандартной реакцией на технологическое развитие в социуме является неолуддизм, противопоставляющий творчество и инженерию, искусство и производство. В контексте математики начала XX века такой реакцией становится интуиционизм, наиболее ярким представителем которого был Лёйтзен Эгберт Ян Брауэр. Он исходит из противоположной установки: математика не есть игра с символами, а изначальная, не языковая, психическая активность человека, коренящаяся в ''первоинтуиции времени''. Интуиционизм процессуален по своей природе: само наше мышление - это процесс, где факты не существуют сами по себе, а проявляются в актах познания. Если формализм движется от синтаксиса к семантике (непротиворечивость порождает модель), то интуиционизм движется от природной интуиции к языку: язык - лишь бледная и всегда неполная запись внутреннего математического опыта (А. Гейтинг). В разрезе выделенных нами шести критериев интуиционизм представляется следующим образом.

\begin{enumerate}
    \item {\it Истина} - В интуиционизме Брауэра истина не является свойством высказывания самого по себе, существующим независимо от познающего субъекта. Истина представляет собой событие - переживание осуществления доказательства. Высказывание становится истинным лишь в тот момент, когда субъект конструирует для него доказательство или демонстрирует процесс, в ходе которого оно осознаётся как необходимое. Таким образом, истина выступает атрибутом позитивного процесса верификации. Она всегда процессуальна и связана с определённым набором ментальных действий, а не с пассивным соответствием высказывания объективной реальности. Аксиомы в интуиционизме - это не произвольные утверждения, принятые на веру в рамках формальной символической игры. Это утверждения, почерпнутые из непосредственного математического опыта и допускающие естественную верификацию, которая и делает их истинными. Все остальные математические истины обеспечиваются процессом доказательства. Вне процесса верификации истины не существует; не может быть истинных, но недоказуемых утверждений. Такой подход накладывает серьёзное ограничение на возможность постижения истины в платоновском смысле. Хотя подобные платоновские истины, возможно, и существуют, для нас они доступны исключительно через совокупность актов верификации. В этом контексте истина предстаёт как класс эквивалентности актов верификации, каждый из которых высвечивает определённые аспекты данной истины. Отсюда следует принципиальная невозможность совершенного языка для формулировки платоновских истин; возможно лишь итеративное развитие, ведущее к их становлению. Вместе с тем Брауэр уделяет значительно больше внимания чувственному, интуитивному переживанию, связанному с истиной, нежели её строго процессуальному аспекту. В результате ключевая для науки воспроизводимость математических истин в интуиционизме приобретает в значительной степени субъективный характер. Ещё одна существенная трудность, вытекающая из процессуальности истины, состоит в необходимости избежать коллапса интуиционизма в более радикальный конструктивизм - задача, с которой было непросто справиться даже самому Брауэру. Тем не менее, развивая эту линию, можно заключить, что интуиционизм, в отличие от конструктивизма, не требует полного отказа от прыжков веры, но требует для них процессуальной фальсифицируемости. Однако неразвитость формального аппарата, описывающего такой процесс, делает интуиционизм в большей степени искусством, нежели инженерной практикой, пригодной для массового использования.
    \item {\it Тождество} - Брауэр принимает логику равенства, но с важным отличием: тождество должно быть разрешимо. В интуиционизме тождество требует основания - мы должны уметь предоставить акт идентификации. Тождество обеспечивается совпадением мыслеобразов по значимым в данном контексте характеристикам. Два объекта считаются тождественными, если в рамках конкретного интеллектуального построения они не различаются по рассматриваемым свойствам. Например, последовательности $1/n$ и $1/2n$ в контексте предела (становления) определяют одно число 0 - в этом контексте они равны. Однако, их отношение всегда 2 - в этом контексте они не равны. Последовательности $1/n$ и $1/(n + 1/n)$ равны в контексте становления предела отношения, но не равны в интенсиональном контексте (способ задания).  В таком случае корректнее говорить об эквивалентности, оставляя тождество для случаев полного совпадения всех релевантных характеристик, включая способ задания. Доказательство тождества зачастую позитивно, как и большинство утверждений в интуиционизме. Отсутствие отличительных свойств влечёт тождество только в контексте теории, замкнутость которой относительно рассматриваемого набора свойств доказана. Интуиционизм различает изоморфизм и тождество: эквивалентность подразумевает лишь частичную неразличимость в пределах фиксированного контекста и языка, в то время как тождество подразумевает принципиальную невозможность различить объекты средствами теории. К сожалению, интуйционизм не смог дать общего строгого описания контекстности определения тождества, за исключением частного реализованного в конструктивном подходе. Что снова сталкивает нас с фундаментальным ограничением - отсутствием языка для описания самих интуитивных процессов.
    \item {\it Выразимость объектов} - математические объекты - ментальные конструкции, возникающие из базовой интуиции. Они существуют тогда и только тогда, когда существует ведущее к ним рассуждение. Нет платоновского существования вне познания, нет неидентифицируемых объектов, поскольку любое утверждение о таких объектах неверифицируемо, а значит, не имеет смысла. Математика - это попытка описать то, как реальность отражается в сознании; это язык описания мыслительной деятельности. Любая математическая теория - это совокупность взаимоотношений мыслеобразов. Любой объект является мыслеобразом, а значит, выразим в теории. Объекты связаны с актом мышления; они не обязательно являются завершёнными конструкциями - это могут быть процессы и циклы. Однако наши суждения о них могут быть почерпнуты только из непосредственного конструктивного изучения таких объектов; любые свойства и характеристики, не следующие из опыта взаимодействия с ними, отвергаются. Например, последовательности выбора - последовательные процессы выбора чисел, разворачивающиеся в некоторой среде и являющиеся инструментом познания этой среды. Такие последовательности принципиально незавершимы, доступны нам в виде частичных конечных отрезков, тем не менее позволяют экстраполировать опыт, связанный с ними, на среду их формирования. Однако, любое описание способно выразить лишь ограниченную проекцию исходных ментальных конструкций, оставляя смысл интуитивных утверждений расплывчатым. Что накладывает существенные ограничения на продуктивный обмен опытом в рамках чистого интуиционизма.
    \item {\it Логическая замкнутость} - логика в интуиционизме - не универсальная истина, а всего лишь формальная фиксация наиболее ярких закономерностей, почерпнутых из опыта высказываний о конечных объектах. В интуиционизме не требуется, чтобы каждое высказывание было либо истинно, либо ложно. Содержание высказывания процессуально, а процесс не обязательно завершается позитивно или негативно - существует зона неразрешимого. Брауэр отказывается от закона исключённого третьего в силу того, что некоторые вопросы ещё не могут быть разрешены или принципиально неразрешимы, а без разрешения мы просто не можем назвать нечто истинным или ложным. Именно поэтому интуиционизм не впадает в противоречие при обсуждении парадокса лжеца. Парадокс возникает не из-за незаконной самореференции, а из-за: ложного предположения, что об истинности или ложности можно говорить без акта их установления, а также из-за отсутствия явного различия уровней языка и метаязыка. В данном контексте утверждать собственную ложность без доказательства невозможно, и интуиционистски это просто некорректное утверждение. Тем не менее, закон исключённого третьего в некоторых контекстах вполне может оказаться доказуемой истиной. Если мы можем показать логическую замкнутость теории, то закон исключённого третьего доказуем - например, в случае конечных объектов, в противном случае у него нет разрешимого статуса. Это, пожалуй, одно из ключевых наблюдений интуиционизма: классическая логика фиксирует ретроспективные отношения состоявшихся процессов, а не потенциальные возможности текущих.
    \item {\it Континуум} - Брауэр понимает континуум как среду свободного становления, которая никогда не исчерпывается вставкой новых единиц. Континуум дан в интуиции как протяжённость. По сути, Брауэр апеллирует к древнегреческому представлению о непрерывном как о бесконечно делимом субстрате, подобном любой своей части. Таким образом, точки - это не составные элементы континуума, а сечения, которые не способны изменить природы непрерывного. Всякая внутренняя точка делит непрерывный отрезок на две части, каждую из которых мы можем делить в тех же пропорциях, что и другую, и исходный объект. В этом и состоит подобие континуума своим частям. Континуум - это среда определения отношений между объектами. С какими бы наборами объектов мы ни работали, эта среда позволяет охарактеризовать их взаимные связи линейным образом. Континуум - это мыслеобраз единой линейной шкалы для сравнения характеристик произвольных классов объектов. Это область определения величины - фундаментальной линейной характеристики объектов. В интуиционизме величина - это не статичная метка, а процесс сужения интервала. Континуум представляет интерес не сам по себе, а именно как среда: помещая объекты в эту среду, мы получаем возможность сравнивать их через отношение разделённости.
    \item {\it Метатеория} - в интуиционизме математика - это единая ментальная активность, которая не расщепляется на уровни рефлексии. Метатеория - это просто ещё одна порция математического конструирования. Если мы рассуждаем о свойствах системы, мы всё равно занимаемся описанием активности интуиции по отношению к зафиксированной форме. Таким образом, математика обосновывается изнутри самой себя через рефлексию над интуитивными актами. Метатеория - часть мыслеобраза, неотделимая от теории. Теория сама является совокупностью связанных мыслеобразов и допускает исследование. Саморефлексия - неотъемлемая часть процесса познания и рассуждения о природе. Нет иерархии Тарского как жёсткого разделения уровней; есть рефлексивный круг, в котором теория может оборачиваться на себя, потому что и объект, и мета-объект - это мыслеобразы одного порядка. Метатеория не сильнее и не вне теории; она есть самосознание теории, разворачивающееся в том же пространстве интуиции. Вопрос о предикате истины не возникает, поскольку на любом уровне рефлексии истина неотделима от процесса её верификации. Однако, из-за невозможности строгой фиксации мыслеобразов, метатеоретические размышления в интуиционизме приобретают ещё более размытый и отвлечённый характер.
\end{enumerate}

Необходимо отметить, что авторы сознательно отделяют интуиционизм от конструктивизма, которые часто путают. Брауэр был скорее мистиком и экзистенциалистом, чем инженером, поэтому конструктивность возникла как попытка найти альтернативу формальным системам. Конструктивизм стал языком, выражающим интуиционизм в научном сообществе. Тем не менее, сам интуиционизм значительно шире и направлен именно на выработку математического языка, примером чего и служит конструктивизм. Вместе с формализмом как альтернативный ответ на массовость и специализацию математики появляется логицизм. Если формализм предлагал песочницы - изолированные формальные системы, заданные набором непротиворечивых утверждений, принятых на веру, то логицизм ставит своей целью выводимость всех математических принципов из чистой логики. Математика, с этой точки зрения, - не игра с символами, а продолжение логики, её естественное расширение и усложнение. По сути, все математические структуры появляются как результат многократного воспроизводства логических примитивов. Логика представляется как фундаментальная истина о реальном мире, а математика - её сложное дедуктивно развёрнутое следствие. Готлоб Фреге пытался вывести арифметику из законов логики (''Основания арифметики'' и ''Основные законы арифметики''), добавив лишь один сомнительный принцип - ''основной закон V'' - объём понятия, универсализирующий концепцию числа, за что подвергся жёсткой критике со стороны Кантора. Именно этот постулат стал причиной тяжёлого удара по логицизму - парадокса, сформулированного Бертраном Расселом, о невозможности в системе Фреге предиката, который неприложим к самому себе. Рассел и Альфред Уайтхед попытались реабилитировать логицизм и развить программу Фреге, введя теорию типов для избежания парадоксов, но и их труд оказался не лишённым утверждений, принятых на веру, которые невозможно объяснить логически. В современной форме логицизм представлен неологицизмом (Криспин Райт, Боб Хейл, Ричард Кимберли Хек), пытающимся возродить программу Фреге, используя как аналитическое основание арифметики принцип Юма (два числа равны, когда между соответствующими им множествами можно установить взаимно однозначное соответствие), также делающий понятие числа универсальным - вопреки интуиции Кантора.

\begin{enumerate}
    \item {\it Истина} - в классическом логицизме истина и ложь постулируются как примитивные, семантически неделимые объекты. Истина объективна и не зависит от познающего субъекта - это свойство мысли, соответствующее действительности, резонирующее с ней. Согласно Фреге, истинность является не свойством знака, а объективным свойством мысли, указывающим на соответствие положению дел в мире. В отличие от формализма, сводящего истину к синтаксической непротиворечивости, логицизм настаивает на дескриптивном характере логических истин. Центральным инструментом системы Фреге является разделение смысла и значения (Sinn/Bedeutung), где смысл выступает как контекстуальный способ презентации объекта, а значение отождествляется с самим объектом. Однако в стремлении к универсальности Фреге совершает онтологическую редукцию: он выводит истину и ложь за пределы смыслового расслоения, наделяя их статусом абсолютных, самотождественных значений. Авторы полагают, что эта категорическая установка (абсолютизация дихотомии истина/ложь без возможности интенсионального расслоения самих логических примитивов) стала фундаментальной причиной кризиса логицизма и возникновения теоретико-множественных парадоксов. Фактически, Фреге неявно использует механизм дифференциации смыслов лжи (например, при определении нуля через отрицание самотождества), но отказывается обсуждать этот процесс. В этом как раз и состоит основное ограничение логицизма, мы вынуждены описывать уже сформированные сущности нашего мышления, оставляя за кадром сам процесс их конкурентного становления.
    \item {\it Тождество} - в системе Фреге концепт тождества неразрывно связан с категорией значения, а не смысла. Тождество постулируется как объективное отношение между значениями (объектами), независимое от познающего субъекта. Тождество фиксирует совпадение значения при различии способов его презентации, когда различные пути расслоения смысла сходятся в единой точке (например, ''Утренняя звезда'' и ''Вечерняя звезда'' как Венера или $2 + 2$ и $4$). Формально подход логицизма опирается на экстенсиональный принцип Лейбница, согласно которому тождество объектов гарантирует совпадение всех их предикативных свойств. При этом в логицизме проводится строгое разделение между тождеством и эквивалентностью. Тождество рассматривается как отношение, приложимое исключительно к объектам (например, ''Лейбниц'' и ''Фреге'' не тождественны). Эквивалентность выступает как функциональное понятие, фиксирующее общность различных отображений. Так, концепты ''быть Лейбницем'' и ''быть Фреге'', будучи различными по смыслу, являются эквивалентными, поскольку совпадают их объёмы. Переход от функциональной эквивалентности к объектному тождеству осуществляется через принцип абстракции - ''основной закон V''. Совокупность понятий, находящихся в отношении эквивалентности по объёмам, определяет новое логическое значение - число. Таким образом, канонический концепт числа легитимизируется как логический объект, возникающий в результате абстрагирования от содержательного различия смыслов в пользу тождества их количественных характеристик. Однако, логицизм по замыслу Фреге разделявший синтаксическое и семантическое тождество, под влиянием формализма, очень быстро эволюционирует к неразрывности синтаксической формы и семантического содержания (появляется негласное правило: синтаксические конструкции имеют одно и то же значение независимо от контекста), что существенно затрудняет анализ парадоксов. 
    \item {\it Выразимость объектов} - фундаментальная установка логицизма постулирует возможность полной редукции математики к логике, где математические объекты мыслятся как производные логические конструкции, детерминированные исключительно набором логических констант (кванторов, связок, предикатов) и переменных. В этой парадигме объекты (числа, объёмы понятий) обладают статусом объективных сущностей, которые, не будучи данными в чувственном опыте, легитимизируются через процедуру строгого логического определения. Для Фреге логика выступала не просто формальным исчислением, но универсальным языком описания реальности, способным гарантировать существование объектов на основании их непротиворечивой дескрипции. Однако обнаружение парадокса Рассела выявило сущностный разрыв между логической определимостью и онтологическим существованием. Демонстрация невозможности ''объёма понятия, не содержащего самого себя'' в системе Фреге разрушила претензию логицизма на создание бессубъектного, самодостаточного языка реальности. Попытка Рассела и Уайтхеда спасти систему не увенчалась успехом, поскольку не давала логического обоснования введённой иерархии типов, а также ряду принятых ими аксиом. Проблема того, что заданное синтаксическое выражение может иметь различные семантические значения в зависимости от контекста высказываний, способная пролить свет на сущность противоречий так и не была в должной мере осознана, о чем говорит отсутствие явного аксиоматического связывания синтаксических конструкций с семантическим содержанием.
    \item {\it Логическая замкнутость} - в логицизме закон исключённого третьего $p \vee \neg p$ обладает статусом фундаментального, детерминированного природой истины и лжи как исчерпывающей дихотомии значений. Для Фреге он являлся базисом дедуктивной системы, гарантирующим определённость любого логического акта. Однако эта претензия на исчерпывающую универсальность делает логицизм уязвимым перед лицом парадоксов самореференции. Осознавая опасность, Фреге вводит уровни применимости концептов, стремясь исключить прямую самореференцию. Тем не менее, парадокс Рассела продемонстрировал, что ''лжец'' может быть реконструирован опосредованно, а концепция универсального числового объекта - расширения, представляющего предикат - внутренне противоречива. Вместо того чтобы интегрировать теорему о неустранимости противоречий в структурах определённого типа, логицизм ограничился ad hoc решениями - внешним выстраиванием иерархий. Тем не менее, в рамках логицизма может быть сформулировано фундаментальное ограничение: невозможно существование предикатов как объектов, обладающих универсальной применимостью. Логическая замкнутость системы оказывается несовместимой с претензией на универсальную (предикативную) категоризацию - что мы и наблюдаем в рамках формализма. Логицизм пасует перед лицом возникших парадоксов стараясь избегать их, вместо того, чтобы интегрировать и исследовать.
    \item {\it Континуум} - в логицизме континуум вводится как производный логический объект, конструируемый из рациональных чисел посредством логического понятия сечения (как у Дедекинда). Поскольку натуральный ряд полагается выводимым из чистой логики, вещественные числа должны представлять собой результат дальнейшего развёртывания логических истин. В этом контексте континуум определяется как линейно упорядоченное поле, подчинённое требованию полноты порядка: для любого свойства $P$, если все элементы, удовлетворяющие $P$, превосходят все элементы, его не удовлетворяющие, между ними существует точная грань - супремум, устраняющий любые структурные пробелы. Непрерывность в логицизме постулируется не как геометрическая интуиция, а как логическая необходимость. Однако противоречивость универсального понятия числа вынуждает перенести конструкцию континуума в жёсткие рамки теории типов. Что приводит к зависимости анализа от нелогических постулатов: аксиомы бесконечности, гарантирующей наличие необходимого субстрата, и аксиомы сводимости, без которой определение вещественных чисел через сечения оказывается невозможным. В конечном итоге логицизм оказывается неспособным обеспечить обоснование континуума.
    \item {\it Метатеория} - логицизм базируется на принципе логического монизма: отрицании разделения между теорией и метатеорией. Логика понимается как универсальный, самоприменимый язык, не допускающий внешней позиции наблюдения. В этой парадигме любые метатеоретические утверждения (включая вопросы о непротиворечивости или полноте системы) должны быть выражены в терминах самой логики. Однако развитие логицизма привело к вынужденному отказу от этого монизма в теории типов. Введение иерархических уровней организации фактически превратило каждый последующий тип в метатеорию по отношению к предыдущему. Несмотря на этот технологический разрыв, логицизм сохраняет принципиальное отличие от формализма: истина в нём остаётся единой и содержательной на всех уровнях иерархии. В то время как формализм рассматривает металогику лишь как синтаксический инструмент проверки правил игры, логицизм настаивает на онтологической целостности истины. Таким образом, даже признавая поражение в попытке построения универсальной системы, логицизм сохраняет интуицию единого смыслового поля.
\end{enumerate}

Зафиксируем, что Фреге пытается обойти парадокс лжеца, считая его слабостью логики, вместо того чтобы обратить его в силу. В самом деле, мы можем сформулировать следующую теорему:

Теорема 1 (основная теорема логики): {\sl Никакая теория с классической логикой не может быть предикативно замкнутой (иметь предикаты как объекты, применимые ко всему). Усиленная версия: никакая теория с классической логикой не может быть опосредованно предикативно замкнутой (иметь предикаты, применимые к любым объектам, и объекты, однозначно связанные с предикатами).}

Доказательство: 

Предположим противное и рассмотрим в предполагаемой теории предикат $P(x) := \neg x(x)$, но тогда $P(P) \Leftrightarrow \neg P(P)$ - противоречие. Для усиленной версии обозначим факт того, что объект $x$ связан с предикатом $X$, как $x \models X$. Рассмотрим предикат $P(x) := \exists X (\neg X(x) \wedge x \models X)$, откуда $P(p) \Leftrightarrow \neg P(p)$ - противоречие. Теорема доказана.

Отметим, что в современном контексте данную теорему можно сформулировать и в логике первого порядка для случая, когда на языке теории выразимо отношение $R(x, y) \Leftrightarrow y \models Y \wedge Y(x)$. В этом случае отношение $R$ позволяет кодировать применимость предиката к объекту, не требуя квантификации по предикатам напрямую, и теорема сохраняет свою силу. Данное наблюдение открывает интересный взгляд на теории, где отношение $R$ не может быть выражено явно, без объединения синтаксиса и семантики, но допустимо в интерпретирующих эту теорию моделях. Эта простая теорема возвращает логицизму его основополагающую глубину, декларируя, что любые системы, построенные из чистой логики, будут обладать данными ограничениями. Что мы де-факто наблюдаем: в новых основаниях Куайна (NF), ограничивающих применимость предикатов (стратификация); в ZFC и NBG, ограничивающих универсальность предикативной категоризации; в теориях типов, вводящих иерархии применимости предикатов; в $\lambda$-исчислении и других формализмах. Теорема показывает, что парадокс лжеца - не дефект, а позитивный закон, фиксирующий границу любой формальной системы, претендующей на универсальность. Это устраняет необходимость в ad hoc запретах (теория типов, стратификация) и превращает их в следствия фундаментального ограничения, а не во внешние дополнения. Кроме того, любопытно сравнить данный результат с результатом в системе Кантора о невозможности консистентной теории, где выполняется аксиома выбора и существует множество замкнутое относительно взятия степени (Теорема 11, "Несчётность множества континуум, шах и мат"). Удивительно, что обе теоремы фактически об одном и том же логическом ограничении, которое теория Кантора не просто его избегала через иерархию, а умела формулировать, задолго до того, как оно стало причиной кризиса оснований математики.

Конструктивизм зарождается как попытка синтеза альтернативного языка в рамках интуиционизма Брауэра, но быстро обретает самостоятельное значение в условиях растущей потребности в алгоритмической ясности и верифицируемости математических утверждений. Если интуиционизм Брауэра остаётся глубоко философским и почти мистическим, то конструктивизм пытается перевести эту философскую позицию на язык точной, алгоритмически осмысленной математики. Он открыто критикует формализм впитавший в себя метатеоретические допущения вызванные недостаточной выразительностью языка, на которые мы указали ранее. Математика здесь понимается как деятельность по построению объектов и доказательств, обладающих явной вычислимой природой. Андрей Марков выделил три абстракции, лежащие в основании классической математики: абстракцию отождествления, абстракцию потенциальной осуществимости и абстракцию актуальной бесконечности. Конструктивизм принимает первые две, но отвергает третью. Перейдём к обсуждению наших критериев:

\begin{enumerate}
    \item {\it Истина} - истина в конструктивизме отождествляется с наличием явного конструктивного доказательства (или алгоритма). Утверждение считается истинным только тогда, когда мы обладаем конкретным, конечным, в принципе выполнимым процессом, который это утверждение обосновывает. Наиболее естественно считать, что сам термин истина обозначает класс конструктивных построений. Колмогоров предложил интерпретировать интуиционистские формулы как задачи, логические связки - как действия над ними. Например, $p \vee q$ задача, требующая решить хотя бы одну из задач (предъявить решение $p$ или $q$). $\exists x \; \varphi(x)$ задача, решаемая предъявлением конкретного элемента $a$ с доказательством $\varphi(a)$. Здесь нет места истинным но недоказуемым утверждениям. Более того, конструктивисты часто отвергают само понятие платоновской истины как бессодержательное. Это делает конструктивизм наиболее строгим в эпистемологическом смысле (соответствующим критерию Поппера) среди всех пяти парадигм. В конструктивизме признаётся принцип Маркова: для обоснования уже проделанных построений можно пользоваться классической логикой. Это позволяет отделить процесс построения от процесса обоснования, не смешивая их. Стоит обратить внимание на фиксацию логики, как способа ретроспективного анализа. Что приводит к сильному ограничению конструктивизма на формулировку абстрактных определений. Необходимость абстракции признаётся, но повисает в воздухе без аппарата способного зафиксировать свойства новых сущностей. Любые математические рассуждения базируются на абстракциях, в том числе и фундамент конструктивизма, но на практическом уровне мы ограничиваем работу с ними, что создаёт внутреннее напряжение. 
    \item {\it Тождество} - тождество в конструктивизме должно быть конструктивно разрешимым. Два объекта считаются тождественными только тогда, когда существует алгоритм, который за конечное время устанавливает их равенство (или, в более мягких версиях, когда мы обладаем конструктивным доказательством их равенства). Центральная находка, выделенная Марковым - абстракция отождествления - метод формирования абстрактных понятий в рамках конструктивизма. Этот приём позволяет игнорировать различия между равными объектами, не учитывая свойства, которые не релевантны в данном контексте. Здесь также проводится строгое различие тождества и эквивалентности. Тождество фиксируется в контексте всей теории, эквивалентность - на уровне контекстов заданных в рамках теории. Две конструкции могут быть эквивалентны (т.е. давать одинаковые результаты в заданном контексте), но не тождественны по способу построения. Конструктивисты стремятся к тому, чтобы равенство было не только доказуемо, но и вычислимо. Это ещё больше ограничивает возможности конструктивного описания абстракций, которые зачастую связаны с неконструктивными совокупностями и свойствами. Что усиливает напряжение, мы полностью выключаем один из механизмов анализа, который иногда приводит к контринтуитивным результатам, но существенно чаще оказывается обоснованным, не имея возможности ввести для него полноценную замену.
    \item {\it Выразимость объектов} - здесь конструктивизм занимает наиболее радикальную позицию. Математический объект существует только постольку, поскольку мы можем его явно построить (дать алгоритм его построения). Как и в случае истины, сам термин существование естественно считать обозначением соответствующего класс математических конструкций. Неконструктивное существование (существует $x$, такое что $\varphi(x)$, полученное, например, через сведение к абсурду) отвергается как бессмысленное. Последовательности выбора Брауэра (принципиально незавершимые процессы) в конструктивизме заменяются вычислимыми алгоритмическими последовательностями. Это делает конструктивизм более техничным и менее метафизичным, чем интуиционизм. Здесь допустимо отождествление конструктивных объектов со словами в конечных алфавитах, а конструктивных процессов - с алгоритмами. Даже кванторы рассматриваются как конструктивные объекты, так квантор всеобщности интерпретируется как построение функции, сопоставляющей каждому элементу доказательство. Однако, на этом уровне конструктивизм сталкивается с фундаментальным идейным противоречием, синтаксическая формальная игра в символы по заданным правилам оказывается вполне конструктивной, но не может иметь семантического наполнения в рамках конструктивизма. В рамках данной парадигмы не удаётся найти полноценной альтернативы бюрократической целостности формализма без возврата к расплывчатости интуиционизма, что делает отказ от закона исключённого третьего экзистенциальным тезисом конструктивизма. 
    \item {\it Логическая замкнутость} - конструктивизм отказывается от закона исключённого третьего $p \vee \neg p$ по тем же основаниям, что и интуиционизм: для произвольного высказывания может не существовать ни доказательства, ни опровержения. Более того, он часто отказывается и от закона двойного отрицания $\neg\neg p \to p$, по крайней мере, в тех системах, где отрицание понимается как приведение к абсурду. Тогда отрицание $p$ становится алгоритмом, который превращает любое гипотетическое доказательство $p$ в противоречие, а двойным отрицанием - алгоритм, который доказывает, что допущение о противоречивости $p$ само ведет к противоречию, что без заданного доказательства $p$ не очевидно. Конструктивизм не впадает в противоречие при обсуждении парадокса лжеца. Парадокс лжеца в конструктивизме не возникает как серьёзная проблема, потому что самореференция либо блокируется на уровне языка (как в теории типов), либо теряет смысл из-за процессуального понимания утверждений. Утверждение: данное утверждение ложно - не может быть конструктивно обосновано, следовательно, оно просто не считается осмысленным высказыванием в математическом смысле. Полный отказ от закона исключённого третьего очень сильно ограничивает аналитический аппарат теорий, в особенности по отношению к абстрактным понятиям. В результате, в конструктивизме развиваются различные альтернативы. Такие как концепция позитивистской абстракции, позволяющей делать обобщения избегающие отрицания. Или би-интуиционистские подходы стремящиеся компенсировать утраченную часть логики постулируя предположения о дуальных операциях: ко-отрицинии и ко-импликации. Но даже они не снимают напряжённости вызванной отказом от части логического аппарата. В то же время, такой отказ позволяет конструктивизму допускать тотальную предикативную замкнутость, при использовании позитивной или паранепротиворечивых (например, минимальной логики Йохансона) логик.
    \item {\it Континуум} - континуум в конструктивизме - это не множество актуально существующих точек (формализм) и не среда становления (интуиционизм), а конструктивный процесс порождения вещественных чисел. Марков развил конструктивный анализ, в котором: вещественное число - это алгоритм, вычислимая последовательность рациональных приближений с заданным модулем сходимости, а не готовый точечный объект, в то же время, алгоритм как последовательно интерпретируемая программа строго конструктивен и не допускает произвола. Таким образом, последовательности выбора Брауэра в конструктивизме интерпретируются как алгоритмические детерминированные процессы. Например, допустимый в интуицианизме алгоритм вычисления числа, такой что $a_n = 0$, если среди чисел от $0$ до $n$ нет контрпримера для гипотезы Гольдбаха, $a_n = 1/k$ в противном случае, где $k$ - первый контрпример, не является числом в конструктивизме, так как не имеет заданного модуля сходимости. Конструктивисты отказываются от готового несчётного континуума. Непрерывность вводится через процессы и аппроксимации, а не через актуально бесконечные множества. Само понятие непрерывности претерпевает изменения, она понимается как возможность вычисления расстояния между двумя объектами с какой-угодно точностью. А непрерывность функции - как свойство, при котором малые изменения входных данных с любой заданной точностью приводят к малым изменениям выходных данных, при этом данное соответствие вычислимо.
    \item {\it Метатеория} - конструктивизм стремится к максимальной внутренней прозрачности и самодостаточности. Идеал - такая теория, в которой метатеоретические рассуждения могут быть проведены внутри самой теории конструктивными средствами. Поскольку паранепротиворечивые логики допускают тотальную предикативную замкнутость необходимости в мета-уровнях не возникает. Теория может описывать свои истины без противоречия. В то же время, использование логик с принципом взрыва (из противоречия следует всё что угодно) приводит к расслоению, не позволяя полностью избежать иерархии (универсумы типов, иерархия предикативности и импредикативности). Однако конструктивизм относится к этому более осознанно, чем классический формализм: он явно признаёт границы своей выразительной силы и старается работать в рамках явно заданных конструктивных ресурсов. С другой стороны, конструктивизм очень быстро принимает правила формальной игры, не формулируя проблему соответствия синтаксиса и семантики. Отказ от однозначной интерпретации синтаксической формулы семантическим значением представляет собой альтернативу жесткому отрицанию закона исключённого третьего, позволяющим усомниться в единственности семантического содержания синтаксических выражений полученных в рамках формальной конструктивной игры, но данное направление не исследуется.
\end{enumerate}

Каким бы парадоксальным это не казалось, на сегодняшний день конструктивизм - это наиболее инженерная из всех рассмотренных парадигм, несмотря на её происхождение из интуиционизма. Если формализм создаёт "песочницы" (формальные системы), то конструктивизм предлагает инструменты - алгоритмы, программы, автоматические доказательства. Конструктивизм, в отличие от других парадигм, не постулирует неявных допущений о существовании объектов вне конструктивной деятельности. Именно эта установка делает конструктивизм естественным языком для вычислительной математики и программирования.

Структурная парадигма возникает в середине XX века как наиболее естественный и мощный ответ на кризис классического формализма. Пока формализм пытался свести математику к безопасным «песочницам» - изолированным формальным системам с жёстким синтаксисом, - математика продолжала бурно развиваться вглубь, оперируя всё более абстрактными понятиями. Формализм, ориентированный на синтаксическую непротиворечивость и внутреннее строение объектов, всё хуже справлялся с объяснением того, почему одни и те же структуры возникают в совершенно разных областях и почему изоморфные объекты ведут себя похожим образом независимо от своей реализации. Теория категорий, предложенная Эйленбергом и Маклейном, и её дальнейшее развитие в работах Гротендика предложили радикально иной взгляд: математика изучает не объекты сами по себе, а структуры, определяемые исключительно через отношения (морфизмы) между ними. Объект полностью исчерпывается своими связями внутри категории. Семантика перестаёт быть первичной или вторичной по отношению к синтаксису - она становится эмерджентным свойством сети отношений. Структурная парадигма становится доминирующей в современной чистой математике и фактически выступает универсальным языком, способным объединять результаты разных школ. Рассмотрим структурную парадигму в разрезе шести ключевых критериев.

\begin{enumerate}
    \item {\it Истина} - в структурализме истина - семантическое понятие возникающее в контексте логических систем, где зачастую понимается как воспроизводимость структурных паттернов. Истинным считается утверждение, которое сохраняется при всех изоморфизмах (или, шире, при всех эквивалентностях категорий). Категориальные утверждения часто имеют условный, схематический вид: если такая-то диаграмма коммутативна, то выполняется такое-то свойство. Истина может определяться взаимосвязностью морфизмов вывода, обеспечивающей непротиворечивость, потенциальной завершимостью конструктивных морфизмов, или фундаментальной примитивной категорией. Она возникает как инвариантность относительно преобразований, сохраняющих структуру. Это позволяет говорить как об истине в контексте, так и об истине до выбора оснований, что делает структурный подход исключительно гибким и устойчивым к смене фундамента. Однако, во многих случаях структурализм ограничивается упаковыванием старых понятий в новую обёртку. Множества упаковываются в категории, отношения в морфизмы. Мы отказываемся от части аксиомантической семантики, что позволяет делать обобщённые абстрактные наблюдения, над уже установленными фактами, например, при соотнесении супремума и ко-предела, но не затрагивает внутренней логики теории, наполняющей содержанием требование универсальности свойства, которое в результате так и остаётся интуитивной концепцией. Впрочем, это наблюдение не отменяет того факта, что структурализм делает категории классических формальных система и альтернативных теории равноправными обитателями категориального пространство, что само по себе ценно, поскольку снимает напряжённое идеологическое противостояние.
    \item {\it Тождество} - в структурализме тождество радикально переосмыслено. В структурной математике два объекта считаются одинаковыми не тогда, когда они совпадают поэлементно или синтаксически, а когда между ними существует изоморфизм. Классическое равенство $=$ часто заменяется на изоморфизм $\cong$ или даже на более слабые понятия эквивалентности. Однако, это не означает, что структурализм не различает форм эквивалентности, просто в контексте структурализма зачастую они не играют существенной роли. В целом, изоморфные в рамках структуры объекты считаются тождественными, в то время, как эквивалентность распространяется на ограниченные контексты внутри структуры (подструктуры), но подструктура также связана с абстракциями узлов - объектами, для которых эквивалентность становится тождеством. Все объекты в структурализме определяются с точностью до изоморфизма. В пределе (в высшей категории теории или гомотопической теории типов) тождество само становится структурой - пространством морфизмов высшего порядка. Это позволяет избежать многих проблем реификации, характерных для формализма. Например, парадоксы типа Банаха–Тарского перестают восприниматься как странности и видятся как особенности конкретной структурной организации категории множеств. В то же самое время, объект при таком подходе превращается в набор чистых абстракций представляющих его проекции на структуры в которые он вовлечён. Классическая универсальная семантика здесь лишена смысла, теория не позволяет делать утверждения о таких объектах как число или множество, только об их структурных проекциях. 
    \item {\it Выразимость объектов} - объекты в структурализме представляют собой узлы в сети морфизмов, сама возможность идентифицировать объект обусловлена его изоморфизмами. Существование объекта оправдано, если он занимает определённое место в сети морфизмов, то есть обладает универсальным набором свойств. Таким образом, объекты возникают как эмерджентные сущности в структурах, где их морфология полностью обусловлена связями с другими объектами. Категориальная формулировка часто позволяет вообще не указывать, что объект представляет собой. Важно лишь, как он взаимодействует с другими объектами. Это даёт выразительную силу: одна и та же категория может иметь совершенно разные реализации (сетевую, топосную, алгебраическую), но при этом все структурно значимые теоремы переносятся автоматически. Выразимость здесь возникает как часть семантики: объект существует постольку, поскольку он распознаётся через свои структурные характеристики, если семантика позволяет говорить об объекте через граф связей допускающий символьную репрезентацию, то объект выразим в теории. Однако, данный подход не позволяет говорить о существовании объекта, который объединяет ряд структурных проекций. Поэтому вопрос о существовании реальных объектов остаётся в рамках структурного подхода переводится в разряд метатеоретических, более того, метатеоретическим является сам вопрос о выразимости структурной позицией конкретного объекта.
    \item {\it Логическая замкнутость} - структурная парадигма гибко относится к классической логике. В обычных категориях (например, в категории множеств) используется классическая логика с законом исключённого третьего. Однако в более общих топосах логика может быть интуиционистской или любой другой. В частности, как и в конструктивизме существуют категории допускающие тотальную предикативную замкнутость. Теория категорий не навязывает единую логику, а позволяет изучать, как именно логика возникает из структуры топоса. Парадокс лжеца и другие парадоксы самореференции не представляют собой проблемы в рамках структурализма, они рассматриваются как простые структурные паттерны, которые становятся коллизиями лишь в эмерджентной семантике, которая не всегда однозначна. В зависимости от контекста, для работы с классической логикой здесь может применятся как стратификация - ограничение применимости предикатов, так и ограничение универсальности категоризации - возникновение иерархий. Структурный подход не требует тотальной замкнутости в одном языке - он работает с семействами эквивалентных описаний и переходами между ними. Распределив разные подходы по разным контейнерам - категориям, структурный подход зачастую ограничивается структурными обобщениями, без детального анализа совокупности разделённых подходов и их взаимосвязи. Что обеспечивает возможность широкого выбора подходов, в то время как инструментарий по созданию новых остаётся ограниченным.
    \item {\it Континуум} - структурализм с одной стороны не имеет единого понятия континуума. Континуум может пониматься как несчётное множество в рамках категории множеств. Или как поле сюрреальных чисел не имеющее представление на языке теории множеств ни как собственно множество, ни, возможно, даже как собственно класс. С другой стороны, структурализм обладает достаточной выразительностью, чтобы говорить об общей идеи континуума как о среде линейных отношений между любыми объектами. Само понятие непрерывности может пониматься по-разному в зависимости от контекста. Оно может быть классическим - наличие грани разделяющей два множества. Оно может быть аристотелевским - неограниченная делимость промежутков. Или любым другим способом. Это возвращает древнюю идею самоподобия и подобия частей целому на новом, более абстрактном уровне, избегая при этом ловушек актуальной бесконечности в её наивном теоретико-множественном виде. Структурализм оказывается верен себе и в вопросе континуума, предлагая серию различных альтернатив описания континуума, избегая вопрос о его сущностной природе.
    \item {\it Метатеория} - структурализм не имеет выраженной позиции относительно мета-теории. Если рассматривать вопрос в общем, то в структурный подход описывает сам себя, поскольку любую теорию можно рассматривать как структуру. В то же время, в контексте категорий неполных теорий структурализм принимает необходимость иерархии семантики, рассматривая это скорее как ограничение на возможность задать плоскую семантику для определённых структур, нежели как какое-то структурное ограничение само по себе. При этом структурализм сознательно отказывается от претензии на абсолютный фундамент: разные основания (ZFC, топосная теория, гомотопическая теория типов, ETCS и др.) могут служить внутренними языками для одних и тех же структурных утверждений. Интерес представляют лишь инварианты, независимо от того какую семантику они приобретают в тех или иных случаях. К сожалению, структурализм во многом остаётся поверхностным обеспечивая равноправие альтернатив и универсальность структурных суждений, но не предоставляет инструментарий анализа внутренней морфологии объектов и разработки новых подходов. В частности, задача изоморфизма синтаксических и семантических структур внутри одной категории зачастую не ставится.
\end{enumerate}

Структурная парадигма представляет собой наиболее последовательную попытку освободить математику от зависимости от квинтесенции объектов и от выбора логических или онтологических примитивов. Объекты здесь полностью реляционны: они - узлы в сети морфизмов. Семантика возникает как эмерджентное свойство структуры, а не как заданное наперёд содержание или конструктивный процесс. Эта парадигма лучше всех остальных объясняет единство современной математики и её способность переносить результаты между, казалось бы, далёкими областями. Она не отрицает другие подходы, а скорее поглощает их. Однако, к сожалению, инерция традиционных взглядов оказывается чрезвычайно высокой и в теорию категорий просачиваются убеждения вызвавшие кризис формализма. В частности, взгляд на актуальную бесконечность, экстенсиональность, специфическую экспоненциальную структуру, унаследованную из аксиоматических теорий множеств. В результате, структурализм даёт толчок в избавлении проблема реификации, однако, не устраняет концептуальные проблемы формализма. В то же время, как мы видим работы Кантора аппелируют к совершенно другой интерпретации.

Приведённый анализ ограничивается наиболее базовыми конструкциями в рамках парадигм и не затрагивает не менее интересные вопросы связанные со смыслом более сложных объектов, таких как натуральные числа или функции. Тем не менее, он наглядно демонстрирует, насколько фундаментально различаются подходы к изучению математики в зависимости от лежащих в их основе метаматематических, философских и методологических допущений. Эти допущения играют решающую роль, поскольку определяют сам язык описания математических конструкций и характер их взаимодействия. В связи с этим любая серьёзная дискуссия о математических утверждениях должна начинаться с чёткого указания того теоретического и методологического контекста, в рамках которого она ведётся. К сожалению, математики крайне редко учитывают это обстоятельство, критикуя альтернативные подходы и полученные в их рамках результаты. Авторы полагают, что все рассмотренные здесь парадигмы сохраняют свою актуальность в современной математике и не исключают, а взаимно дополняют друг друга. Формализм предоставляет ряд несомненных преимуществ: максимальную точность и строгость контроля над доказательствами, лёгкость стандартизации и алгоритмизации (что делает математику технологичной дисциплиной), минимальные требования к содержательному пониманию (достаточно владеть правилами символической игры), а также возможность работать с очень сильными формальными системами, включая аксиомы, которые мы не в состоянии содержательно обосновать. Вместе с тем формализм обладает существенными ограничениями. Он замкнут в рамках уже выработанного языка и системы понятий и не содержит внутренних механизмов их генерации. Аксиомы в нём могут быть произвольными и слабо связанными с интуицией или физической реальностью. Получаемые результаты зачастую трудно интерпретировать содержательно из-за уязвимости к ошибке реификации. Кроме того, формализм индифферентен к вопросам онтологии и эпистемологии, что открывает дверь для интерпретаций, основанных на вере или идеологических предпочтениях (в частности, платоновского реализма или некритически принимаемых метатеоретических допущений). 

Формализм представляет собой мощный инженерный инструментарий, идеально подходящий для хранения, передачи и формальной верификации знаний. Однако, будучи вырванным из строгого методологического каркаса, он склонен к чрезмерной бюрократизации и злоупотреблению своим положением. Особо следует отметить, что современный формализм в значительной степени утратил первоначальную чистоту, впитав ряд необоснованных допущений. Большинство из них восходит к Фреговской универсализации понятия числа и к отказу от предложенного Кантором принципиального разделения кардинальной и ординальной алгебр как независимых структур. В интерпретации Кантора, в частности, множество $2^\omega$ не может быть переупорядочено в порядковый тип $\omega$. Вследствие этого биекция между Гёделевской нумерацией формул и порядковыми числами становится невозможной. Это обстоятельство снимает напряжение с понятия истина внутри формальной теории, делая её индуктивно (иерархически) допустимой. Существование принципиально невыразимых объектов перестаёт быть неизбежным следствием теории, поскольку отсутствие биекции с Гёделевской нумерацией не приводит к противоречию. Кроме того, ослабевают жёсткие требования к метатеории. Мало кто отмечает, что сама по себе невыразимость истины в достаточно сильной теории порождает тонкость с самореференцией, близкую к парадоксу лжеца. Действительно, если никакая достаточно сильная теория не может выражать собственную истину, то в какой именно теории мы делаем это утверждение? Претендуя на его истинность, мы тем самым предполагаем, что теория, в которой оно формулируется (а она сильнее рассматриваемой, так как выражает истинность в ней), способна выражать свою собственную истинность, что приводит к противоречию: теория утверждает своё собственное несуществование как достаточно сильной теории. Более корректно утверждать, что для всякой достаточно сильной теории $T$ существует более сильная метатеория $T'$, в которой доказуема невыразимость истины в $T$. Само же утверждение о невыразимости истины в достаточно сильной теории вообще не имеет формального доказательства - оно фиксирует структурное свойство иерархии на внетеоретическом философском уровне. Заметим, что такого эффекта не возникает в теореме 1, поскольку в рамках логицизма истина первична и существует как данность. Таким образом, формализация и критический анализ неявных допущений современных аксиоматических теорий множеств позволили бы очистить формализм от идеологической составляющей и вернуть ему первоначальную интеллектуальную мощь и методологическую честность.

Интуиционизм, в свою очередь, обеспечивает максимальную глубину рефлексии и осмысления механизмов математического мышления. Он позволяет формировать понятийный аппарат и язык, способный схватывать то, как реальность отражается в сознании. Интуиционизм не только требует глубокого понимания контекста, но и активно развивает его, избегая необоснованных скачков веры в виде аксиом, не подкреплённых интуитивным обоснованием. В этой парадигме истина возникает как результат понимания и невозможна без него. Вместе с тем интуиционизм страдает высокой степенью расплывчатости и сильной контекстной зависимостью. Его сложно формализовать и превратить в эффективную математическую технологию. Он требует постоянного внимания к состоянию сознания математика и поэтому ближе к искусству, нежели к инженерной дисциплине. Необходимость обосновывать и глубоко понимать каждое утверждение приводит к существенной потере математической силы по сравнению с подходами, допускающими аксиоматические скачки веры. Таким образом, интуиционизм остаётся преимущественно инструментом решения сложных концептуальных задач: определения границ применимости формальных систем и создания новых понятийных каркасов. Преимуществом логицизма является стремление к максимальному единству математики и её радикальной редукции к простейшим логическим примитивам. Он даёт глубокое понимание того, как сложные математические паттерны возникают из элементарной дихотомии истины и лжи, сочетая это с высокой философской строгостью и ясностью оснований. Логицизм легко поддаётся стандартизации и алгоритмизации, что позволяет использовать его в качестве технологической платформы. Однако логицизм оказывается крайне громоздким по сравнению с другими подходами, что ярко демонстрирует ''Principia Mathematica'' Рассела и Уайтхеда. Это делает его применение в повседневной математической практике малоэффективным. Логицизм традиционно считается уязвимым к парадоксам самореференции. Тем не менее, как было показано в Теореме 1, данная уязвимость скорее свидетельствует о его способности явно формулировать принципиальные ограничения, присущие любым достаточно сильным формальным системам. Тем самым логицизм предоставляет логическое обоснование для введения ограничительных аксиом и специальных методов. Авторы убеждены, что именно способность к смысловому расслоению уже на уровне примитивных понятий, наиболее ярко проявляющаяся в парадоксе лжеца, составляет созидательную силу логицизма. В частности, первый арифметический объект - ноль - возникает из абсурдного свойства $x \neq x$: то, что представляет собой абсурд, превращается в идентифицируемый объект с истинным свойством $0 = 0$. В этом смысле логицизм можно рассматривать как своеобразный клей, соединяющий различные формальные системы в единую осмысленную конструкцию и обеспечивающий им необходимое философское и логическое обоснование.

Конструктивизм обеспечивает максимальную достоверность результатов: все утверждения здесь верифицируемы, а конструкции воспроизводимы. В рамках данного подхода не возникают парадоксы и внутренние противоречия, поскольку всё, что сконструировано, по определению является частью существующей математической действительности. Конструктивизм легко поддаётся инженеризации, так как каждое утверждение обладает явным алгоритмическим содержанием. Благодаря этому он оказывается наиболее близким к реальной вычислительной практике среди всех рассмотренных парадигм. Вместе с тем конструктивизм существенно ограничивает выразительную силу математического языка. Жёсткие ограничения на формирование абстракций (в частности, на понятие предела) заметно снижают мощь конструктивных теорий. В этом смысле конструктивизм можно рассматривать как прагматичный минимальный каркас, задающий необходимые требования к любой теории, претендующей на достоверное описание математической реальности. Наконец, структурализм обеспечивает максимальную гибкость, провозглашая семантику эмерджентным свойством математических структур. В рамках этого подхода объекты определяются не через внутреннее содержание, а исключительно через систему отношений, в которые они вступают. Структурализм представляет собой своеобразный гибрид, сочетающий синтаксическую строгость структурного уровня с контекстно-зависимой семантикой. Такой подход позволяет избежать ошибки реификации, прочно связывая математические структуры с картой, а не с территорией. Структурализм допускает активную стандартизацию и инженеризацию в пределах конкретных структурных доменов. Он соединяет высокую абстрактную выразительность, позволяющую формулировать фундаментальные идеи (например, континуум как самоподобие), с мощью теоретических построений, выходящих за рамки описания формальных систем (в частности, сюрреальные числа). Однако, подобно интуиционизму, структурализм требует глубокого понимания абстрактных связей и их контекстной зависимости. Хотя семантика здесь эмерджентна, она далеко не всегда очевидна и зачастую допускает множественные интерпретации. В результате структурные утверждения нередко приобретают форму чрезмерно абстрактных конструкций, оторванных от наблюдаемой действительности. Структурализм, как и логицизм, объединяет различные формальные системы через воспроизводимые структурные паттерны. Однако, в отличие от логицизма, он не предлагает им философского или логического обоснования, а лишь устанавливает связи между разными семантиками. Тем не менее, на практике такой подход оказывается чрезвычайно эффективным благодаря возможности автоматического переноса методов и теорем в изоморфных структурах.

Если рассматривать математическое творчество Георга Кантора, то его с полным основанием можно считать одним из первых структуралистов, сознательно интегрирующих формальный подход для доказательства чистого существования. Кантор последовательно выстраивает две независимые алгебраические структуры - кардинальную и порядковую, неоднократно подчёркивая их принципиальную несводимость друг к другу. Независимость данных алгебр является методологическим условием использования классической логики внутри каждой из них, именно существование одной бесконечности как противоположности конечности гарантирует двузначность пространства альтернатив. Мейнстрим не заметив глубины онтологии Кантора утрачивает методологическую объективность в попытке описать трёхзначное пространство альтернатив: конечное, кардинально бесконечное и ординально бесконечное, в языке предназначенном для описания двухзначного пространства. Кантор утверждает, что числа являются эмерджентными свойствами структуры, ещё задолго до того как это становится мейнстримом. Обычно понятие поля описывается через две основные алгебраические структуры - аддитивную и мультипликативную группы. Однако на эту же структуру можно взглянуть с более фундаментальной точки зрения. Арифметика Пеано напрямую не постулирует указанные алгебраические операции, однако они естественным образом возникают внутри неё как производные структуры. Простое индуктивное построение лежит в основе сложения; сложение, в свою очередь, служит основой для определения умножения. Таким образом, мультипликативная структура является производной по отношению к аддитивной. Это особенно наглядно проявляется при рассмотрении аддитивной структуры в показателе степени: даже нейтральный элемент мультипликативной группы (число 1) выражается через аддитивный нейтральный элемент ($1 = x^0$). Аддитивная структура воспроизводит себя на всё более высоких уровнях - в показателе показателя степени и далее, вплоть до пределов примитивной рекурсии. Данный факт имеет принципиальное значение для анализа. Он препятствует попыткам исправить определение предела путём простого учёта обеих алгебраических структур поля. Казалось бы, если поле полностью определяется аддитивной и мультипликативной составляющими достаточно вместо классического расстояния $|x_n - x|$ рассматривать комбинированную метрику, например, вида $|\ln(x_n/x)| \cdot |x_n - x|$. Однако такой подход остаётся недостаточным: необходимо учитывать все уровни воспроизводства аддитивной структуры, а не только её мультипликативное воплощение. Переходя на язык Брауэра можно сказать, что аддитивно-мультипликативная пара просто задаёт новое направление первоинтуиции времени. Кантор, по-видимому, ясно осознавал эту глубокую структурную особенность при построении порядковой арифметики.

Кантор строит алгебру порядковых типов на основе операции конкатенации - сложения в упорядоченном смысле. В этом контексте предельный объект принципиально отличается от предела, возникающего при неупорядоченных (чисто кардинальных) операциях, и допускает осмысленное внешнее обобщение операций сложения и умножения (через порождение новых объектов $\omega + 2$, $\omega*2$). Онтологический статус образующих алгебру операций играет центральную роль в корректном разграничении внутренних и внешних функциональных отношений. Стремление воспроизвести аддитивную структуру для бесконечных упорядоченных объектов закономерно приводит Кантора к необходимости введения бесконечного показателя степени. Если попытаться конструктивно проследить последовательность $2, 2^2, 2^3, \dots, 2^n, \dots$, то обнаруживается, что она перемежается с последовательностью натуральных чисел, а её предел в наивном смысле оказывается равным $\omega$. Этот результат связан с тем, что все конечные степени, обусловленные внутренней структурой монады $\omega$, попадают в области сингулярности между зонами регулярности аддитивной структуры. Однако равенство $2^\omega = \omega$ явно противоречит теореме Кантора. Кантор находит изящное решение, переходя от структурной парадигмы к формальной. Он доказывает чистое существование функции возведения в степень, не требуя, чтобы область значений этой функции была замкнута в области определения. Благодаря этому аргумент о перемежании последовательностей теряет силу, поскольку последовательности теперь принадлежат разным множествам. Кантор явно требует, чтобы перемежающиеся последовательности принадлежали одному и тому же множеству, и строго соблюдает это условие в своих доказательствах \cite{cantor-alternate}. В частности, при доказательстве равенства $(\alpha + \nu_0)\omega = \alpha \omega$ общее множество фиксировано - это именно результат умножения $(\alpha + \nu_0)\omega$, заданный онтологическими операциями порождающими алгебру. В этом множестве последовательности перемежаются, и не существует подмножества (а следовательно, и порядкового типа), которое полностью содержало бы одну последовательность и исключало другую \cite{cantor-outer-multiply}. Напомним, что в онтологии Кантора, порядковые типы являются аннотациями множеств и в один порядковый тип могут быть включены множества разной длины. Мы сначала комбинаторно перемножаем множества, а потом аннотируем конструкцию последовательностями порядковых типов. На первый взгляд аналогичная ситуация возникает и при возведении в степень: кажется, что $\omega$ должно содержать как последовательность $\nu$, так и последовательность $2^\nu$. Однако здесь требуется функция степени вида $p : \omega \to \{2^\nu \mid 2^\nu \in 2^\omega\}$, а не $p : \omega \to \{2^\nu \mid 2^\nu \in \omega\}$. Мы не можем гарантировать принадлежность обоих последовательностей одному множеству, что не позволяет использовать теорему о перемежающихся последовательностях. Причина заключается в том, что для аддитивного воспроизводства структуры в качестве показателя степени необходим именно объект $\omega$, а не некий замещающий объект, обозначаемый тем же символом. В противном случае пришлось бы заново строить всю арифметику для этих новых объектов, включая пределы и степени. Использование уже имеющегося объекта $\omega$ существенно упрощает построение. В результате, нам требуется бесконечная сумма из правых операций сложения или бесконечное произведение из правых операций умножения, которые по определению являются внешними и способны пересекать предельные границы. Таким образом, у нас практически отсутствуют какие-либо основания считать множества на которых определены данные последовательности совпадающими.

Именно в эту ловушку, которую Кантор тщательно обошёл, попадают современные аксиоматические теории множеств. В них возникает промежуточный объект $2^{\tilde{\omega}}$ - множество всех конечных подмножеств бесконечного множества. Показатели степени ведут себя здесь иначе, чем сами ординалы: если для ординалов $\bigcup \omega = \omega$, то для экспоненциальной конструкции мы не можем записать аналогичное равенство $2^{\bigcup \omega} = 2^{\tilde{\omega}} \neq 2^\omega$, поскольку объединение всех конечных степеней приводило бы к полноценному булеану, биективно не отображаемому на $\omega$. Таким образом, появляется идея промежуточной бесконечности - объекта, который включает все комбинации до любого конечного уровня, но при бесконечном количестве элементов уже не допускает полного комбинаторного составления. По своему смыслу это возвращает нас к кардинальной алгебре в её потенциально-бесконечной форме. Операции над таким объектом (например, $2^{\tilde{\omega} + 1}$) ведут себя скорее кардинально: $2^{\tilde{\omega} + 1} = 2^{\tilde{\omega}}$, поскольку добавление одного элемента ко всем конечным комбинациям оставляет их конечными. Обратим внимание, что когда мы говорим, о всех конечных комбинациях порядковый тип исходного множества перестаёт играть какую-либо роль. В построениях Кантора отсутствует аксиома экстенсиональности в современном виде, а предел последовательности связан с множеством, которому эта последовательность принадлежит. В частности, последовательность $2^\nu$ может принадлежать только множеству $2^\omega$, поскольку $\omega$ выступает как класс эквивалентности (монада), и её внутренняя структура играет существенную роль при помещении в показатель степени - что подтверждается теоремой Кантор. Для иллюстрации можно привести пример $\omega - 1$, который на уровне порядкового типа соответствует $\omega$: $\omega - 1$ = $\omega$. Все операции в онтологии Кантора связаны с множествами, мы лишь постфактум аннотируем их порядковыми типами. Мы строим множество $2^\omega$ как предел - бесконечный ряд, но последовательности $\nu$ и $2^{\nu}$ не могут перемежаться, поскольку они принадлежат множествам, которые невозможно переупорядочить изоморфно друг другу. Поэтому, при аннотации картина выгдлядит непривычно: хотя все члены последовательности $2^\nu$ строго меньше $\omega$, их предел перескакивает сразу к $2^\omega$. В случае онтологии Кантора это не представляет проблемы, так как у нас отсутствует экстенсиональность и мы справедливо можем сказать, что $\{2^\nu \mid 2^\nu \in 2^\omega\}$ и $\{2^\nu \mid 2^\nu \in \omega\}$ представляют собой разные множества, а невозможность выразить отличающие их элементы - всего лишь следствие импредикативности правого конца вполне упорядоченных множеств, что необходимо нам для замыкания алгебраической структуры. Кантора в этом контексте демонстрирует более уверенное владение инструментами предикативности и импредикативности, чем современный мейнстрим. 

В аксиоматической теории множеств реификация множеств и аксиома экстенсиональности меняют картину: подмножества степеней двойки от $2^\omega$ и от $\omega$ совпадают, и как следствие оба являются подмножествами $\omega$, что позволяет говорить о перемежании $\nu$ и $2^{\nu}$. Это не создаёт противоречия при определении пределов через бесконечные ряды, но приводит к трудностям при переходе к определению предела через супремум. В последнем случае последовательность $2^\nu$ не может иметь пределом $2^\omega \neq \omega$, так как все степени двойки строго меньше $\omega$, а значит именно $\omega$ является супремумом для них. Данный эффект довольно естественен, поскольку свойства алгебраических операций напрямую зависят от онтологии образующих операций, Кантор выбрал в качестве такой операции конкатенацию вполне упорядоченных отрезков, где добавление отрезков справа способно изменять структурный паттерн, что противоречит замкнутости сепремума. Иными словами, данный эффект связан с сознательным отказом от исходной онтологии. Однако, такой ход не был в достаточной мере продуман, поскольку при объяснении счётности ординала $\omega^\omega$, который является пределом последовательности $\omega^n$, возникает более серьёзная проблема. Здесь требуется нечто большее, чем любая конечная степень, что заставляет вводить гипотетическую промежуточную бесконечность $\tilde{\omega} < \omega$, которая уже не конечна, но ещё не приводит к несчётности результата. Можно было бы попытаться обосновать такую картину, введя, например, факториал и определив $\omega^\omega = \omega! = \sup \{ n! \mid n < \omega \}$. Тогда результат оставался бы в рамках счётных объектов. Однако подобные попытки, как правило, не предпринимаются; вместо этого новая интерпретация принимается без достаточного критического анализа. Стремление к ней понятно: математическое сознание естественным образом коллапсирует к привычной картине и пытается применить к трансфинитным степеням теоремы о перемежающихся последовательностях вместо теоремы Кантора. С формальной точки зрения одной интуитивной картины недостаточно. Необходимо строго обосновать легитимность именно картины с возникновением новой сущности - промежуточной бесконечности, - а не той, в которой теорема Кантора защищает теорию от парадоксов. На данной развилке отсутствуют убедительные аргументы в пользу промежуточной бесконечности. Следовательно, её принятие требует дополнительных допущений, формализацию которых мы рассмотрим далее. Отказ от теоремы Кантора в данном контексте, как мы увидим, ведёт к теориям условную непротиворечивость, которых сложно устранить, тогда как сохранение подхода Кантора позволяет поддерживать полноценную непротиворечивость.

В статье «Несчётность множества континуума: шах и мат» были введены понятия внутренней и внешней алгебры. Эти понятия связывались с ролью, которую играет внутренняя структура множества при выполнении над ним операций: в кардинальной алгебре результат операций не зависит от структуры множества, тогда как в порядковой алгебре структура множества существенно определяет результат. Вместе с тем данные понятия обладают ясным интуитивным содержанием, которое, по замыслу авторов, помогает прояснить ту концептуальную картину, которой, по-видимому, придерживался Георг Кантор. В настоящей работе мы даём строгое определение внутренних и внешних конструкций. Для этого заметим, что любое математическое утверждение неявно связано с позицией референта - определённым аспектом контекстной информации, относящейся к субъекту высказывания. Рассмотрим два примера. Утверждение ''две различные точки отрезка ограничивают другой отрезок'' фиксирует позицию референта, связанную с внутренней конституцией отрезка. Напротив, утверждение ''трансляция отрезка $[0, 1]$ на величину $l$ даёт отрезок $[l, l + 1]$'' предполагает позицию референта, описывающего отрезок во внешнем контексте пространства. Формально для высказывания $P$ будем говорить о {\it внутренней позиции референта} по отношению к объекту $O$ или {\it внутреннем} высказывании по отношению к объекту $O$, если все связи между сущностями, влияющие на истинность $P$, полностью определяются морфологией (внутренней структурой) самого объекта $O$. Напротив, будем говорить о {\it внешней позиции референта} в высказывании $P$ по отношению к объекту $O$ или {\it внешнем} высказывании по отношению к объекту $O$, если его истинность зависит от связей между морфологией объекта $O$ (или им самим) и другими объектами контекста. Данные концепции носят в значительной степени интуитивный характер и наиболее наглядно иллюстрируются геометрическими примерами. Рассмотрим два произвольных подобных треугольника $A$ и $B$. Очевидно, что любые отношения между внутренними элементами треугольника $A$ воспроизводятся во внутренних элементах треугольника $B$. В частности, для любой точки на стороне треугольника $A$, делящей эту сторону в определённой пропорции, существует соответствующая точка на соответствующей стороне треугольника $B$, делящая её в той же пропорции. Таким образом, масштабное преобразование устанавливает биекцию между треугольниками $A$ и $B$. Такое утверждение является внутренним по отношению к обоим треугольникам, поскольку соответствие между точками полностью определяется связями внутри каждого из них. Сравним теперь эти два треугольника между собой. Для этого необходимо поместить оба треугольника в единое пространство или, как минимум, соотнести их так, чтобы один оказался вписанным в другое. В этом случае мы переходим к внешней позиции референта по отношению к обоим треугольникам, поскольку положение точек одного треугольника зависит от пространства или точек другого. Примечательно, что в данном контексте биекция между треугольниками $A$ и $B$ как множествами точек с внешней позиции референта (изоморфное отображение сохраняющее метрические отношения) невозможна: каким бы образом мы ни разрезали треугольники, мы не сможем совместить их так, чтобы все точки одного полностью совпадали со всеми точками другого. Здесь уместно сослаться на аксиому полноты Гильберта: в континуум нельзя добавить точку, не нарушив при этом остальных аксиом.

К сожалению, в современной математике позиция референта систематически игнорируется. Предполагается существование универсальной, нейтральной позиции, с которой можно однозначно описать всю математическую реальность. Однако уже парадокс лжеца показывает, что такое предположение некорректно. Истинность высказывания синтаксически определяется отношением между его внутренними терминами. Так, высказывание ''сейчас идёт дождь'' выражает отношение между термом ''настоящий момент'' и термом ''дождь''. В случае же высказывания ''данное высказывание ложно'' его внутренними терминами являются само высказывание и понятие ложности. То есть, высказывание внутренне ссылается на связь между собой и ложностью, однако семантически (внешне) оно эту связь отрицает. Парадокс возникает именно вследствие попытки одновременно занимать обе позиции референта - находиться в двух местах одновременно. Коллизия возникает не столько из-за доступа к внешнему контексту, где определена семантика лжи, сколько из-за противопоставления внутренней (синтаксической) и внешней (семантической) позиции. Таким образом, ключевое ограничение, сформулированное в Теореме 1 и воспроизводимое во всех формальных системах с классической логикой, является по сути ограничением самого языка. Позиция референта всегда подразумевается и составляет неотъемлемую часть передаваемого смысла. В рамках одного высказывания невозможно одновременно занимать внутреннюю и внешнюю позицию референта по отношению к одному и тому же объекту. Когда мы говорим о подобии фигур, мы имеем в виду их внутреннее морфологическое сходство, структурную эквивалентность. Напротив, тождество подразумевает эквивалентность фигур во внешнем контексте - по отношению к окружающему пространству или другим фигурам. Не случайно Дедекинд использует термин подобие для описания систем, находящихся в биективном соответствии друг с другом \cite{dedekind-infinity}. Точно так же Кантор говорит о подобии и эквивалентности, подразумевая именно внутреннее морфологическое сходство рассматриваемых объектов \cite{cantor-similarity}, а отнюдь не их количественное тождество. Оба автора заимствуют свойство самоподобия у непрерывных объектов и распространяют его на бесконечность, что позволяет чётко отделить конечное от бесконечного и исследовать последнее как самостоятельный объект. В структурализме мы обнаруживаем, что изоморфные структуры наследуют свойства, связанные с контекстом изоморфизма. В частности, изоморфизм самоподобных структур указывает на то, что некоторые свойства непрерывного мира могут быть присущи и бесконечным объектам. Позиция референта в непрерывном мире играет при этом существенную роль, определяя, идёт ли речь о внутренней морфологии объекта или о его внешних взаимодействиях.

Понятия внутренней и внешней позиции референта неожиданно обогащают язык описания бесконечных объектов. Классических понятий оказывается явно недостаточно для объяснения различных аспектов проявления свойств бесконечности. Со времён Аристотеля бесконечность традиционно разделяли на потенциальную и актуальную. Однако этого различения недостаточно. Так, $\aleph_0$ представляет собой актуальную бесконечность, которая ведёт себя как потенциальная за счёт неограниченной рефлексивности ($\aleph_0 + 1 = \aleph_0$) - именно то, что Аристотель называл ''всё новое и новое''. В то же время $\omega$ - это иная форма актуальной бесконечности, поведение которой отличается как от потенциальной бесконечности, так и от $\aleph_0$. Авторы часто пытаются подчеркнуть различие между потенциальной и актуальной бесконечностью, используя понятия ''процесс'' или ''поток'', апеллируя к динамическому характеру потенциальной бесконечности. Однако и с этой задачей успешно справляется $\aleph_0$: процесс, рассматриваемый как статический объект, полностью развёрнутый в пространстве, естественным образом представляется именно этой кардинальной бесконечностью. Не менее проблематичным оказывается понятие целостности или тотальности, которое обычно используется для подчёркивания завершённости актуально бесконечного объекта. Согласно Евклиду, часть всегда меньше целого. Следовательно, рефлексивный объект, для которого $\aleph_0 + 1 = \aleph_0$, невозможно последовательно описывать как целостный: $\aleph_0$ оказывается частью $\aleph_0 + 1$, которая при этом меньше целого. Аналогичная трудность возникает и с ординалом $\omega$: хотя $\omega + 1 > \omega$, всё равно выполняется $1 + \omega = \omega$, то есть добавление элемента не меняет объект. Близкими по смыслу являются понятия интенсиональности и экстенсиональности, однако и они в конечном итоге не позволяют провести необходимое разграничение. Один и тот же объект - $\aleph_0$ как множество всех конечных кардинальных чисел или $\omega$ как множество всех конечных порядковых чисел - может быть описан как интенсионально (все числа определённой природы), так и экстенсионально (целостное множество заданных объектов). Оба подхода применимы как к кардинальному, так и к порядковому описанию. Понятия предикативности и импредикативности оказываются более удачными, они проясняют ситуацию с потенциально и актуально бесконечными объектами. Потенциальная бесконечность полностью предикативна относительно элементов домена, но импредикативна относительно описания самого домена. Точно таким же свойством обладает кардинальная бесконечность $\aleph_0$, где импредикативность имеет абсолютный характер $1 + \aleph_0 = \aleph_0 + 1$, в то время как, порядковая бесконечность импредикативна слева $1 + \omega = \omega$, но предикативна справа $\omega < \omega + 1$. Но данные термины дают язык для описания конфликта предикативности и импредикативности, ответственного за все парадоксы бесконечности, но не дают ключ разрешению данного конфликта. Это технические термины, которые констатируют ситуацию не проясняя глубинные механизмы, которые провоцируют такие конфликты. Тем не менее все эти термины не дают точного и исчерпывающего разделения типов бесконечных объектов. Ситуация кардинально меняется при введении внутренней и внешней позиции референта. Когда мы описываем бесконечность как внутреннюю ($\aleph_0$) и внешнюю (индексное число $\dot{\omega}$), порядковый тип $\omega$ естественным образом занимает промежуточное положение, в котором левые операции носят внутренний характер, а правые - внешний. Конфликт заключается лишь в попытке занять обе позиции референта одновременно. Кантор ясно видел, что числа являются артефактами структуры: без структуры нет чисел. При этом внутренняя структура и внешняя структура - принципиально разные вещи.

На определяющую роль языка в математике указывал Людвиг Витгенштейн. Недопонимание между ним и Куртом Гёделем было далеко не случайным \cite{wittgenstein-godel}. Вопреки распространённому интуитивному пониманию, полнота формальной системы определяется как доказуемость всех истинных утверждений. Витгенштейн заметил, что такое употребление термина нарушает правила языковой игры. Корректнее, по его мнению, говорить, что непротиворечивость теории требует принятия дополнительных метатеоретических предположений. Более того, Витгенштейн, по-видимому, первым предложил корректное и математически безупречное решение парадокса лжеца. Он утверждал, что смысл слова определяется совокупностью всех контекстов его употребления. Иными словами, семантика любого термина есть не что иное, как класс эквивалентности всех контекстов его использования. В случае утверждения ''данное высказывание ложно'' мы сталкиваемся с отсутствием общего контекста использования со всеми остальными употреблениями слова ''ложь''. Другие контексты не предполагают неопределённости состояния того, что объявляется ложным. Следовательно, мы имеем дело с некорректным использованием языка: термин ''ложь'' помещается в контекст, в котором он лишается смысла. Это равносильно утверждению ''данное высказывание - полдень''. Откуда следует, что ключ к математическим противоречиям следует искать в языке, а не в логике или формальных системах. Интуиционизм позволил избежать некорректного использования языка, однако не смог вернуть утраченную в процессе силу теориям. 

С точки зрения когнитивной науки, внутренняя (алгоритмическая) и внешняя (интегративная) перспективы, лежащие в основе способности к абстракции и, в частности, к осмыслению потенциальной и актуальной бесконечности, имеют не только философские и культурные корни, но и глубокие биологические предпосылки. Данное обстоятельство делает их чёткое различение в математике не только желательным, но и неизбежным, придавая этим понятиям междисциплинарный статус, используемый в дальнейшем изложении. Антропологические данные свидетельствуют, что активное развитие мозга гоминид хронологически совпадало с усложнением технологии обработки каменных орудий. Согласно некоторым современным гипотезам когнитивной археологии, инженерные вызовы этой эпохи могли выступать одним из факторов эволюционного отбора нейронных структур, функциональные принципы которых впоследствии проявились в самых разных сферах когнитивной деятельности \cite{antropology-infinity}. В частности, сформировался двухуровневый паттерн организации психических процессов, классически иллюстрируемый речевыми зонами: зона Вернике ассоциируется с интеграцией семантического контекста и формированием целостного представления об объекте (внешний, смысловой контур), тогда как зона Брока — с последовательной, алгоритмической организацией действий и синтаксических форм (внутренний, процессуальный контур). Хотя современные нейровизуализационные исследования демонстрируют более сложную сетевую организацию мозга, само концептуальное разделение на смысловой и процессуальный уровни сохраняет высокую объяснительную силу. Аналогичная функциональная дихотомия (фиксация глобальной цели и пошаговое построение траектории) наблюдается в задачах навигации и других видах сложной целенаправленной деятельности, выступая универсальным когнитивным шаблоном. Альтернативная и комплементарная гипотеза социального мозга \cite{antropology-socium}, которая связывает эволюционный скачок объема коры гоминид с необходимостью обсчета сложных социативно-иерархических связей, указывает на тот же фундаментальный паттерн. В рамках данной концепции формирование модели самосознания строится на абсолютно том же механизме двойной перспективы. Внутренняя перспектива в данном контексте отвечает за непосредственный субъектный опыт индивида - его первичные импульсы, эгоцентрические интересы и алгоритмическую оценку текущей выгоды (референт - сам индивид). Внешняя же перспектива возникает как мета-когнитивная проекция: способность построить в собственном сознании устойчивую инвариантную модель взгляда со стороны - то, как индивид и его действия воспринимаются и оцениваются социумом (референт - другой индивид). Более того, способность к рефлексии и моделированию внешнего наблюдателя, ставшая залогом выживания сложнейших социальных групп гоминид, опирается на тот же структурный мотив связанный с семантикой и зоной Вернике, в то время как, расчёт индивидуальных интересов и стратегии поведения в стабильных контекстах связан с синтаксическими мотивами и зоной Брока. 

Функциональное взаимодействие этих зон тесно интегрировано в работу фундаментальных нейронных подсетей мозга. Большую часть времени при выполнении конкретных операций мозг функционирует в режиме сети активного решения задач (Task-Positive Network, TPN). Это режим высокой концентрации, строгой фильтрации информации и фокусировки на ограниченном наборе объектов. Синтаксическая (внутренняя) перспектива в данном состоянии сужает когнитивное поле, отсекая внеконтекстные ассоциации и подавляя попутные нейронные сигналы через механизмы локального ГАМК-торможения. Для решения же сложных, нестандартных задач мозг переключается в режим сети пассивной работы (Default Mode Network, DMN), отвечающей за интеграцию опыта и мышление, направленное на семантическое (внешнее) соотнесение и позиционирование понятий. В этом состоянии мозг связывает разрозненные факты в широкий контекст, выявляет скрытые ассоциации и рассматривает ситуацию как единое целое. Если TPN опирается на последовательную алгоритмическую логику, то DMN задействует комбинаторную сложность и одновременную активацию распределённых нейронных ансамблей. Хотя базовое энергетическое потребление в обоих режимах сопоставимо (так как торможение нейронов - это активный процесс), интенсивный ассоциативный поиск в режиме DMN с привлечением нейромедиаторных контуров (глутамат, дофамин) существенно увеличивает нагрузку на системы нейромедиаторной отчистки (включая утилизацию глутамата астроцитами), а также на сети когнитивного контроля, удерживающие смысловую целостность. Накопление метаболитов и истощение ресурсов исполнительного контроля делают длительное пребывание в состоянии высокого творческого синтеза энергозатратным, требуя специальных условий для активации (глубокое сосредоточение, медитативные практики, творческое погружение). Такое соотношение и динамический баланс сетей представляются эволюционно адаптивным механизмом оптимизации когнитивных ресурсов. 

Ключевую роль в координации этих двух режимов играет салиентная сеть во взаимодействии с зоной Вернике. Именно они выступают функциональным переключателем, предотвращающим синаптическое и когнитивное переутомление. Когда DMN формирует целостное решение (образ идеального орудия, музыкальную тему или архитектурный замысел), зона Вернике (семантический комплекс) фиксирует его в виде устойчивого семантического аттрактора (ориентира). После этого салиентная сеть подавляет ресурсоёмкий режим ассоциативного поиска DMN и активирует TPN. Зона Брока (синтаксический комплекс), в свою очередь, запускает последовательный процесс приближения к зафиксированному образу через алгоритмические операции. Этот нейрокогнитивный паттерн воспроизводится при решении большинства творческих и прикладных задач: при создании материального объекта DMN и зона Вернике формируют смысловую концепцию изделия, а зона Брока в режиме TPN задействует моторную кору для ее пошаговой реализации; в искусстве DMN генерирует семантическое облако будущего произведения, а TPN и зона Брока переводят его в последовательность мазков кисти или звуковой ряд. В качестве наглядной демонстрации данного механизма на когнитивном уровне выступает классический эксперимент Питера Уэйсона \cite{wason-246}. Испытуемому даётся задача угадать некоторый закон, которым связаны тройки чисел, для чего он может предлагать свои тройки чисел, получая положительные или отрицательные ответы испытателя, и в качестве примера предлагается тройка чисел $(2, 4, 6)$ заведомо удовлетворяющая закону. Сам закон при этом формулируется просто, каждое следующее число превосходит предыдущее. Однако, большинство испытуемых сразу переходят к алгоритмической проверке в режиме TPN, фиксируя первую очевидную гипотезу (например, увеличение на 2) и проверяя её линейно. Их внутренний взгляд остаётся жёстко ограничен рамками исходного контекста: мозг естественным образом стремится экономить ресурсы исполнительного контроля и придерживаться алгоритмически безопасной линейной стратегии вместо перехода в режим глобального переосмысления. 

Таким образом, внутренняя и внешняя перспективы глубоко встроены в фундаментальную архитектуру человеческого мышления, что находит своё отражение в расслоении языкового восприятия на синтаксическую и семантическую составляющие. Автоматический, зачастую неосознанный характер переключения между ними создаёт иллюзию их независимости или даже взаимной исключительности. В силу этого в научной практике, и особенно в основаниях математики, опора лишь на одну из этих перспектив оказывается методологически недостаточной. Полноценная концептуализация понятийных систем требует явного определения внутренней и внешней позиций референта, а также их соотнесения с понятиями потенциальной и актуальной бесконечности. Вероятно, внешняя перспектива и понятие актуальной бесконечности (как завершённой целостности) соотносится преимущественно с функциями семантического комплекса (зоны Вернике) и сети DMN. В свою очередь, внутренняя позиция референта и потенциальная бесконечность (как незавершимый процесс) связываются с работой синтаксического комплекса (зоны Брока) и сети TPN. В связи с этим точки зрения, полностью отвергающие одну из этих форм бесконечности, следует признать не просто упрощённой, но и методологически несостоятельной, поскольку она приводит к разрыву преемственности научной традиции, игнорированию богатейшего философского и математического наследия, а также к утрате междисциплинарных связей с физикой, когнитивистикой и антропологией. Различение данных форм возникает на ранних этапах когнитивного развития и пронизывает фундаментальные структуры знания в самых разных областях. Давайте формально зафиксируем данные понятия. {\it Потенциальная бесконечность} - свойство, выражающееся в возможности неограниченного внутреннего продолжения математического построения. В свою очередь, {\it актуальная бесконечность} - абстракция, выражающая внешнюю фиксацию объекта идеала такого построения. Давайте взглянем на историю развития идеи бесконечности через призму введённых нами понятий. 

В научном сообществе существует устойчивая тенденция подчёркивать прогресс современного знания за счёт упрощённой и часто искажённой интерпретации идей мыслителей прошлых эпох. Особенно показателен в этом отношении пример Аристотеля, глубина философской позиции которого, при внимательном ознакомлении с первоисточниками, оказывается существенно большей, чем принято рассматривать в современных экскурсах. Широко распространено мнение, что Аристотель признавал лишь потенциальную бесконечность и категорически отвергал бесконечность актуальную. Однако прямое обращение к его текстам обнаруживает значительно более тонкую картину. Для Аристотеля потенциальная и актуальная бесконечность понимались иначе, чем в современной математике и философии. Он действительно отрицает существование актуальной бесконечности в природе, поскольку физика как наука о природе имеет дело с конечными сущностями. В то же время вопрос о бесконечности в математических и умопостигаемых объектах он относит к более широкому кругу исследований, выходящему за рамки ''Физики'' \cite{aristotle-infinity-scope}. Аристотель лишь отмечает, что в современной ему математике можно обойтись без актуальной бесконечности, однако не запрещает её в принципе \cite{aristotle-actual-infinity}. Бесконечность он рассматривает как свойство числа и величины \cite{aristotle-infinity-property}. Во многих местах своих сочинений он свободно оперирует бесконечным в математических рассуждениях и прямо утверждает, например, что время бесконечно \cite{aristotle-infinite-time}. Характерно его описание потенциальной бесконечности: ''если, взявши от конечной величины определённую часть, прибавлять [к ней дальнейшие части, находящиеся друг к другу] в одинаковом отношении, но [только] не прибавлять повторно ту же самую часть целого, то [исходную] конечную величину нельзя будет пройти [до конца]'' \cite{aristotle-modern-potential}. Таким образом, бесконечное существует как незавершимый процесс. При этом Аристотель допускает и определённый вид актуальной бесконечности в возможности: ''бесконечное есть материя для завершённости величины, и целое только в возможности, а не в действительности; оно делимо и при уменьшении и обратном прибавлении, а целым и ограниченным [бесконечное] оказывается не само по себе, а по отношению к другому'' \cite{aristotle-modern-actual}. Именно через категорию "по отношению к другому" Аристотель вводит принципиально важное различение перспектив рассмотрения. То, что внутренне делимо до бесконечности, может выступать как целое по отношению к внешнему объекту. Он предостерегает от смешения этих двух перспектив в рамках одной теории, критикуя пифагорейцев: ''нелепо мнение тех, кто говорит так же, как пифагорейцы: они одновременно делают бесконечное сущностью и расчленяют её на части'' \cite{aristotle-pythagoras}. Невозможно одновременно рассматривать бесконечное и как целое, и как незавершимый процесс, не указывая явно на смену перспективы (позиции референта). Эта проницательность Аристотеля выглядит особенно глубокой на фоне современных теорий множеств, которые, спустя два с половиной тысячелетия, во многом повторяют пифагорейский подход, соединяя оба взгляда в единой концепции. К сожалению, детальные представления пифагорейской школы о бесконечности до нас не дошли, что не позволяет в полной мере подтвердить точность аристотелевской характеристики их позиции, согласно которой идея рефлексивности как свойства самого объекта (а не только структуры) зарождается именно у них.

На эту фундаментальную двойственность бесконечного позднее обратил внимание Галилео Галилей в своём знаменитом рассуждении о сравнении множества всех натуральных чисел и множества их квадратов. С одной стороны, квадратов явно меньше, поскольку существуют числа, не являющиеся квадратами. С другой - их столько же, поскольку каждому числу соответствует его квадрат. Галилей, подобно Аристотелю, приходит к выводу о невозможности объединить эти два представления в рамках единого плоского отношения порядка  \cite{galilei-squares}. Бернард Больцано продвинулся дальше, чётко сформулировав дуальность: количество квадратов меньше количества всех чисел по отношению включения, но совпадает с ним по отношению взаимно-однозначного сопоставления \cite{bolzano-duality}. Эта позиция практически тождественна аристотелевской. Георг Кантор в работе ''Основы общего учения о многообразиях. Математически-философский опыт учения о бесконечном'' высоко оценил интуицию Больцано и предложил связать ''сопоставление'' со строгим понятием мощности, а ''включение'' - с понятием количества, связанным с порядком элементов множества \cite{cantor-bolzano}. Иными словами, Кантор развивает аристотелевскую дуальность, соотнося мощность с внутренней морфологией множеств, а порядковый тип - с внешним определением количества элементов. Однако, порядковые типы как объекты оказываются существенно сложнее, сочетая в себе обе перспективы. Кантор характеризует схоластические опасения связанные с бесконечностью как предрассудки и приводит убедительные доводы о необходимости актуальной бесконечности как концепции \cite{cantor-actual-infinity}. Рихард Дедекинд в ''Что такое числа и для чего они служат?'' предложил критерий бесконечности множества через наличие взаимно-однозначного соответствия между множеством и его собственной частью - подобие, создав замкнутое внутреннее описание бесконечных систем \cite{dedekind-infinity}. Кантор в ''К основанию учения о трансфинитных множествах'' идёт дальше, создавая единую непротиворечивую теорию, в которой кардинальные числа и левые операции порядковых типов выражают внутреннюю структуру, а правые операции - внешнюю. Оба мыслителя подчёркивают внутренний характер сопоставления, заимствуя терминологию из геометрии и представления о непрерывности (подобие, эквивалентность). Кантор последовательно отстаивал мысль, что число не является универсальным понятием. В рецензии на ''Основания арифметики'' Готлоба Фреге он резко критикует попытку обосновать понятие числа и мощности через ''объём понятия'', рассматривая это как извращение истины \cite{cantor-frege}. Для Кантора число связано с порядком, а мощность - с общим поэлементным отображением. Бесконечные объекты обладают эмерджентными свойствами, которые не могут быть полностью редуцированы друг к другу в силу фундаментальных различий в способах их конструктивного построения. Разделение алгебр не является прихотью или выбором, это естественная необходимость вызванная методологическими ограничениями использования основного языка математики - логики, требующими строгой двузначности пространства альтернатив в теории. Однако, при аксиоматизации теории множеств без критического осмысления совершается формальная ошибка ложной дихотомии и адаптируется идея Фреге об универсальности понятия числа. Столкновение аксиом внутренней структуры, например, схем выделения и преобразования и аксиом внешней структуры - аксиома выбора, в рамках единой теории и приводит к контринтуитивным результатам.

Работа Георга Кантора поражает глубиной и филигранной точностью интуиции относительно разделения внутренней и внешней алгебры порядковых чисел. Благодаря вполне упорядоченным множествам становится возможным вводить новые математические конструкции, в частности бесконечные ряды, которые не возможны в кардинальном контексте. Кантор естественно связывает их с понятием внешнего предела: $\lim\limits_{\nu}^{outer} \alpha_{\nu} = \alpha_1 + (\alpha_2 - \alpha_1) + (\alpha_3 - \alpha_2) + \dots$ \cite{cantor-limit}. В частности, $\lim\limits^{outer} \nu =  \sum\limits^{outer} 1 = 1 + 1 + \dots  = \omega$. В данном определении вполне упорядоченность гарантирует существование предела за счёт внешней суммы, где все слагаемые добавляются справа. Это также можно рассматривать как обобщение диагонального построения Кантора, вполне упорядочение поддерживает диагональную конструкцию, благодаря которой отличие от всех элементов класса эквивалентности через внешнее сложение $\cdot + 1$ имеет воплощение в форме внутреннего объекта алгебры. В принципе, аналогичным образом можно определить и внутренний предел: $\lim\limits_{\nu}^{inner} \alpha_{\nu} = \dots + (\alpha_3 - \alpha_2) + (\alpha_2 - \alpha_1) + \alpha_1$, для первого трансфинитного числа: $\lim\limits^{inner} \nu =  \sum\limits^{inner} 1 = \dots + 1 + 1 = \aleph_0$. В порядковой алгебре обратный (левый) порядок не является вполне упорядоченным, поскольку не обладает минимальным элементом. Поэтому без дополнительных определений внутренние операции могут быть осмыслены лишь на уровне кардинальной абстракции. Кантор систематически сочетает внутренние и внешние операции для сложения и умножения, однако сознательно вводит только внешнее возведение в степень. Степень определяется им через правое умножение: $\gamma^{\xi + 1} = \gamma^{\xi}\gamma$ и внешний предел: если $\xi = \lim\limits_{\nu}^{outer} \xi_{\nu}$, то $\gamma^{\xi} = \lim\limits_{\nu}^{outer} \gamma^{\xi_{\nu}}$. С формальной точки зрения ничто не мешает определить и внутреннее возведение в степень: $\gamma^{1 + \xi} = \gamma\gamma^{\xi}$, а в случае предела: если $\xi = \lim\limits_{\nu}^{inner} \xi_{\nu}$, то $\gamma^{\xi} = \lim\limits_{\nu}^{inner} \gamma^{\xi_{\nu}}$. При таком подходе естественно возникает интерпретация, близкая к современной аксиоматической теории множеств: $2^{\lim\limits_{\nu}^{inner} \nu} = \lim\limits_{\nu}^{inner} \nu$, поскольку внутреннее (левое) умножение на 2 не выводит за пределы $\omega$. Это даёт $2^{\aleph_0} = \aleph_0$, что согласуется с равенством $\aleph_0 \cdot 2 = \aleph_0$. Такое определение перекликается с современным понятием супремума, однако внутренний предел в рамках порядковой арифметики не является порядковым числом и отождествляется лишь с кардиналом $\aleph_0$. В этом контексте проясняется и смысл выражения $\omega^{\nu}$ в современных теориях. В построениях Кантора такое выражение не может сохранять свойство изоморфизма с $\omega$ при сохранении порядка. При внутреннем возведении в степень $\omega^{\aleph_0}$ оказывается неопределённым в порядковой арифметике, однако по своему содержательному смыслу соответствует именно тому объекту, который в современной теории множеств понимается как множество всех функций с конечным носителем (или потенциально бесконечная степень). Кантор вводит двойственную (внутреннюю и внешнюю) природу операций лишь для сложения и умножения, но сознательно отказывается делать это для возведения в степень. По-видимому, Кантор понимает, что степень, как покрытие одного множества другим, всегда требует внешней перспективы по отношению к множеству показателю, поскольку множество всех подмножеств должно содержать такое множество как объект. Он интуитивно понимал, что смешение перспектив способно приводит к противоречиям и конфликтам с альтернативными морфологиями бесконечности. Таким образом, ключевое допущение, которое отделяет современные аксиоматические теории от оригинальных построений Кантора, можно зафиксировать в лаконичной форме: $\aleph_0 = \omega$, то есть отождествление внутренних и внешних пределов. Следует понимать данное равенство не просто как переобозначение $\omega$, что дало бы чистую порядковую алгебру Кантора, где $\omega < 2^{\omega}$ и $\aleph_0 < 2^{\aleph_0}$, но как наделение $\omega$ способностью соответствовать своему экспоненциалу через внутреннюю индуктивную структуру. Как мы установили в статье ''Несчётность множества континуум, шах и мат'' данное предположение независимо от ZFC, поскольку единственность введённых функциональных отношений недоказуема здесь. Мы можем говорить о некоторой рефлексивной структуре реализуемой через внутренние функции, но утверждение об отсутствии другой структуры требует дополнительных предположений. 

В качестве иллюстрации рассмотрим простую конструкцию, вдохновлённую теорией сюрреальных чисел Джона Конвея. Определим {\it наборы} (объекты игры) как деревья вида $O = \{L \mid R\}$, где $L$ и $R$ - ранее построенные наборы. При этом $R$ называется правым элементом набора $O$, а $L$ - левым элементом. Правый элемент {\it принадлежит} набору: $R \in O$, и отношение принадлежности транзитивно по левому элементу: если $x \in L$, то $x \in O$. Дочерние элементы обозначаются через $O.L$ и $O.R$ (например, $x.R.L$). Наборы строятся индуктивно, начиная с пустого набора $\{ \mid \}$. Операция построения нового набора называется {\it композицией}. Каждый новый набор использует только наборы, построенные на предыдущих этапах: $\{ \mid \}$, $\{ \{ \mid \} \mid \}$, $\{ \mid \{ \mid \} \}$ и так далее.

Теорема 2 (декомпозиции): {\sl всякий набор может быть декомпозирован по своим правым элементам и правым элементам своих левых поднаборов.}

Доказательство:
Все наборы строятся индуктивно. Для пустого набора утверждение тривиально выполняется. Предположим, что оно верно для всех наборов предыдущего индуктивного уровня. Рассмотрим произвольный набор $z$ следующего уровня. Возможны три случая:
\begin{enumerate}
	\item $z = \{ x \mid \}$ - тогда $z$ декомпозируется в последовательность элементов набора $x$;
	\item $z = \{ \mid x \}$ - тогда $z$ декомпозируется в единственный правый элемент $x$;
	\item $z = \{ x \mid y \}$ - тогда $z$ декомпозируется в правый элемент $y$ и все элементы $x$.
\end{enumerate}
Таким образом, всякий набор $z$ может быть представлен в виде последовательности $z.R, z.L.R, z.L.L.R, \dots$, которую назовём декомпозицией набора. Теорема доказана.

Теорема о декомпозиции позволяет определить отношение тождества наборов через совпадение их декомпозиционных цепочек: два набора равны тогда и только тогда, когда каждая цепочка элементов одного набора имеет соответствующий равный элемент в цепочке другого, и наоборот. Такое определение является рекурсивным и распространяется на все наборы.

Следствие 2.1 (экстенсиональность): {\sl для любых наборов $A$ и $B$ выполняется $A = B \Leftrightarrow \forall x (x \in A \Leftrightarrow x \in B)$.}

Данное следствие непосредственно вытекает из определения тождества через декомпозицию. Таким образом, на языке модели мы оперируем понятиями композиции: $L \oplus R = \{ L \mid R \}$ и декомпозиции: $\{ L \mid R \} \ominus R = L$, а отношение принадлежности $\in$ выражается рекурсивно: $x \in O \Leftrightarrow x = O.R \vee x \in O.L$, где $O.L = O \ominus O.R$.

Теорема 3: {\sl игра $\{ L \mid R \}$ интерпретирует ZFC без аксиомы бесконечности.}

Доказательство:
Покажем, что все аксиомы ZFC, кроме аксиомы бесконечности, выполняются в данной конструкции.
\begin{itemize}
	\item Экстенсиональность. Следствие 2.1.
	\item Аксиома выбора. Для любого семейства непустых наборов функция выбора задаётся явно: $f(x) = x.R$.
	\item Аксиома пары. Для любых наборов $A$ и $B$ существует набор, содержащий только их: $\{ \{ \mid A \} \mid B \}$.
	\item Аксиома объединения. Для любых наборов $A$ и $B$ набор, содержащий все элементы обоих, строится индуктивно по декомпозиции $B$: $\{ A \mid B.R \}, \{ \{ A \mid B.R \} \mid B.L.R \}, \{ \{ \{ A \mid B.R \} \mid B.L.R \} \mid B.L.L.R \}, \dots$.
	\item Аксиома выделения. Для заданного свойства $\varphi(x)$ и всякого набора $A$ набор элементов $A$, удовлетворяющих $\varphi$, строится индуктивно по декомпозиции $A$: если $\varphi(x)$, то $x$ включается композицией в результирующий набор.
	\item Аксиома преобразования. Для заданного свойства $\varphi(x, y)$ и всякого набора $A$, такого что для каждого элемента $x$ существует набор $y$ удовлетворяющий $\varphi(x, y)$, образ строится индуктивно по декомпозиции $A$, для каждого $x \in A$ в набор композицией включается единственный набор $y$ удовлетворяющий $\varphi(x, y)$.
	\item Аксиома степени. Для любого набора $A$ существует набор $\mathcal{P}(A)$, содержащий все поднаборы $A$. Для пустого набора данное утверждение верно. По индукции проверим переходы к следующему уровню: $\{ x \mid \}$, $\{ \mid x \}$, $\{ x \mid y \}$. Первый случай сводится к степени $x$, второй - к $\{ \{ \mid x \} \mid \{ \mid \}\}$. В третьем случае каждый элемент $\mathcal{P}(x)$ комбинируется с $y$ и включается композицией в набор $Y$, после чего $\mathcal{P}(x)$ и $Y$ объединяются.
	\item Аксиома регулярности. Всякий набор содержит элемент, с которым он не имеет общих элементов. Данное утверждение верно для первого уровня $\{ \mid \}, \{ \mid \{ \mid \} \}$. При переходе к следующему индуктивному уровню возможны варианты: $\{ x \mid \}$, $\{ \mid x \}$, $\{ x \mid y \}$. Первый случай сводится к регулярности $x$, второй следует из того, что $x \notin x$. В третьем случае, если такой элемент $z$ для $x$ не содержит $y$, то мы нашли искомый элемент. В противном случае, уровень $y$ меньше уровня $x$, а значит рассматриваемый набор принадлежит уровню $x$ для которого мы уже имеет требуемый факт. 
\end{itemize}
Теорема доказана.

Фактически мы построили модель ZFC без аксиомы бесконечности, которая в силу своей конструктивности является заведомо непротиворечивой. Непротиворечивость легко доказывается индукцией по уровням: любая формула языка с операциями $\{ \oplus, \ominus, \in, = \}$ оценивается рекурсивно по структуре наборов, и переход на новый уровень $\oplus$ не может ввести противоречия, так как всякая формула сводится к предыдущим уровням посредством декомпозиции $\ominus$. Поскольку структура набора жестко кодирует историю его сборки, аксиома выбора в данной модели приобретает строго конструктивный, детерминированный характер через операцию декомпозиции. Примеры соответствия стандартным конструкциям: $\emptyset = \{ \mid \}$, $\{ \emptyset \} = \{ \mid \{ \mid \} \}$, $\{ \{ \emptyset \} \} = \{ \mid \{ \mid \{ \mid \} \} \}$, $\{ \emptyset, \{ \emptyset \} \} = \{ \{ \mid \{ \mid \} \} \mid \{ \mid \{ \mid \} \} \}$. Будем называть {\it ордами} рекурсивные наборы вида $\{ x \mid x \}$, где $x$ - орд предыдущего уровня. Очевидно, что орды находятся в естественном соответствии с ординалами. В отличие от классических конструкций, в нашей модели мы можем двигаться по ордам как вперёд - с помощью композиции $x \oplus x = \{ x \mid x \} = x + 1$, так и назад - с помощью декомпозиции $x \ominus x.R = \{ x.L \mid x.L \} =  \{ x.R \mid x.R \} = x - 1$. Такая двунаправленность даёт существенные преимущества. В частности, для любой разности наборов (в смысле теоретико-множественной операции разности) мы можем однозначно сопоставить орд, который назовём {\it расстоянием} между наборами. Для этого при декомпозиции уменьшаемого набора каждый раз, когда встречается элемент вычитаемого набора, мы добавляем новый орд посредством композиции $x \oplus x = x + 1$. Благодаря этому наборы позволяют моделировать также аксиоматику нумеросити \cite{numerosities}. В данной интерпретации размер множества представляет собой орд, то есть набор специальной формы. Тождество двух множеств следует из равенства соответствующих ордов, построенных по их декомпозициям. Сложение определяется через совместную декомпозицию двух непересекающихся наборов, а единица - как $\emptyset \oplus \emptyset$. Отношение порядка для ордов полностью определено отношением принадлежности $x < y \Leftrightarrow x \in y$. В целом, возможность движения в обоих направлениях существенно снижает роль линейного порядка в анализе наборов. Любой набор можно полностью декомпозировать до нужного элемента и пересобрать заново, поместив этот элемент в произвольную позицию, в том числе в конец последовательности. Теперь введём понятие бесконечности в рамках нашей модели. Наиболее удобным представляется использовать бесконечность по Дедекинду, поскольку она непосредственно выделяет свойство, отличающее конечные объекты от бесконечных. Предположим, что существует бесконечный орд $\dot{\omega}$, структурно подобный собственной части. Ни один орд, полученный конечным числом индуктивных композиций, таким свойством не обладает. Таким образом, мы постулируем существование нового типа объекта и должны исследовать его свойства. Структурное подобие для ордов означает следующее. По определению всякий орд имеет вид $\{ x.R \mid x.R \}$. Следовательно, $\dot{\omega} = \{ \dot{\omega}.R \mid \dot{\omega}.R \}$, $\dot{\omega}.R = \{ \dot{\omega}.R.R \mid \dot{\omega}.R.R \}$ и так далее. По теореме о декомпозиции $\dot{\omega}$ может быть разобран до терминального пустого набора. Таким образом, структурно бесконечный орд представляет собой объект, в котором процессы композиции и декомпозиции неограниченно стремятся навстречу друг другу. Формализуем это наблюдение. Орды будем называть {\it индуктивными}, если расстояние от пустого набора до данного орда не превосходит расстояния от него до $\dot{\omega}$: $x - \emptyset \leq \dot{\omega} - x$ (заметим, что $x - \emptyset = x$). В противном случае орд называется {\it дедуктивным}. Это позволяет определить функцию структурного подобия для бесконечного орда:
\begin{align}
\begin{split}
f(x) = \left \{ 
\begin{array}{rcl} 
x + 1 & x \leq \dot{\omega} - x \\ 
x - 1 & x > \dot{\omega} - x
\end{array}
\right.
\end{split}
\end{align}

Тогда, бесконечный орд $\{0,1,2,\dots,\dot{\omega}.R.R.R,\dot{\omega}.R.R,\dot{\omega}.R\}$ оказывается структурно изоморфным своей части $\{1,2,\dots,\dot{\omega}.R.R.R,\dot{\omega}.R.R\}$. На этом основании аксиому бесконечности можно записать в следующем виде: {\it существует внутренне подобное своей части множество $\dot{\omega}$, такое что}
\begin{equation}
\{ \mid \} \in \dot{\omega} \wedge \forall x ((x \in \dot{\omega} \wedge x \leq \dot{\omega} - x) \Rightarrow x + 1 \in \dot{\omega})
\end{equation}

Предлагаемое построение существенно отличается от теории нумеросити и гросс-объектов. Первая опирается на нестандартный анализ и, с точки зрения теории множеств, даёт большие конечные числа, в то время как вторая, явно постулирует Евклидову структуру, которая не допускает внутреннего подобия. То есть, обе теории оперируют Евклидовыми сущностями, для которых целое всегда больше своих собственных частей. Это связано с центральной ролью вполне упорядоченных структур в современных теориях, которая не допускает равноправного рассмотрения альтернативных морфологий. В то же время наши наборы оказываются изоморфными омнифическим целым положительным числам, которые также являются внутренне Дедекиндово-бесконечными (допускают изоморфное со своей частью индуктивное построение). Подчеркнём, что морфология определяет внутреннее строение множества, которое вовсе не обязательно является тотально предикативным. Морфология сохраняет возможность подобия структуры, при этом допуская усиление предикативности и выход во внешний контур, где структуры становятся строго Евклидовыми. Таким образом, мы явно фиксируем разницу между морфологическими свойствами, отражающими особенности внутренней структуры (аксиома индуктивно-дедуктивной бесконечности), и онтологическими свойствами, определяющими различие между объектами и их реифицированным описанием (аксиома Евклида: целое больше своих частей). Данные утверждения (о индуктивно-дедуктивной морфологии и Евклидовости) являются независимыми и должны постулироваться отдельно. Операция возведения в степень $x^y$ естественным образом вводится через функциональное отношение $f : y \to x$, то есть через покрытие набора $y$ элементами набора $x$. В рамках игры функция $f$ строится посредством декомпозиции набора $y$: каждому элементу $y$ поочерёдно сопоставляются все элементы декомпозиции $x$, после чего из полученных пар формируется набор всех возможных комбинаций. Это позволяет доказать теорему Кантора в предложенной модели.

Теорема 4 (Кантора): {\it не существует функции, обеспечивающей изоморфизм между набором и набором всех его поднаборов: $\forall X (X < 2^X)$.}

Доказательство:
Предположим, что для некоторого набора $X$ существует изоморфизм (биекция) $f : X \to \mathcal{P}(X)$. В терминах нашей детерминированной игры это означает наличие строго определённого правила композиции/декомпозиции, сопоставляющего каждому элементу $x \in X$ уникальный поднабор $f(x) \in \mathcal{P}(X)$. Построим диагональный набор $D$ с использованием только разрешённых операций. Выполним полную декомпозицию набора $X$, последовательно извлекая его элементы $x$. Для каждого $x$ рассмотрим соответствующий поднабор $f(x)$. Посредством декомпозиции $f(x)$ проверим, принадлежит ли $x$ к $f(x)$. Если $x \notin f(x)$, включаем $x$ в $D$ посредством композиции; в противном случае пропускаем. Поскольку набор $D$ состоит только из элементов $X$, он по определению является поднабором $X$ и обязан находиться в $\mathcal{P}(X)$. В силу предположения об изоморфизме существует элемент $d \in X$ такой, что: $f(d) = D$. Рассмотрим, принадлежит ли $d$ набору $D$:
 \begin{itemize}
	\item Если $d \in D$, то по правилу построения $D$ это возможно только при условии $d \notin f(d)$. Но поскольку $f(d) = D$, получаем $d \notin D$ - противоречие.
	\item Если $d \notin D$, то по правилу построения $D$ элемент $d$ должен был быть включен, то есть $d \in f(d)$, откуда $d \in D$ - противоречие.
\end{itemize}

Таким образом, предположение о существовании внутреннего изоморфизма приводит к противоречию. Теорема доказана.

Данное элементарное построение наглядно демонстрирует глубину и силу оригинальной идеи Кантора. С формальной точки зрения корректнее говорить о классах эквивалентности ордов по внутреннему изоморфизму (типах ордов), как это делали Кантор и Дедекинд. Особенно это существенно при работе с бесконечными ордами. Кантор подчёркивал геометрическую аналогию, называя вполне упорядоченные фрагменты отрезками, под которыми обычно подразумевается именно тип морфологии, а не фиксированная система точек. При необходимости мы можем воспользоваться техникой осознанной реификации Джона фон Неймана и работать с каноническими представителями этих типов. В любом случае внутренние операции получают чёткое процедурное воплощение. Для внутреннего сложения достаточно полностью декомпозировать бесконечный набор, добавить новый элемент посредством композиции и объединить результат с декомпозированной частью. Внутреннее умножение удлиняет встречные композиционную и декомпозиционную цепочки, в то время как внешнее умножение копирует структуру многократно. Однако ни орды, ни их типы не могут быть изоморфны набору (или типу набора) всех своих поднаборов. Теорема Кантора гарантирует принципиальное доминирование внешнего взгляда над внутренним и, соответственно, фундаментальную ограниченность рефлексивности. Именно это мы наблюдаем в оригинальной теории Кантора: внешняя составляющая ординалов доминирует над внутренней, а внутреннее возведение в степень как операция отсутствует, ведь степень предполагает интерпретацию показателя как фиксированного объекта. Для достижения полной рефлексивности, соответствующей классическому представлению о натуральном ряде, Кантор выделил отдельно арифметику кардинальных чисел, которая согласуется с рефлексивными системами Дедекинда. В таких системах теорема Кантора становится недоказуемой, поскольку отсутствует структура, способная корректно выразить отношение самопринадлежности. Поскольку универсум наборов является ограниченным поддеревом сюрреальных чисел Джона Конвея, предложенная игра остаётся непротиворечивой. Типы ордов служат лишь инструментом внутреннего структурного анализа и не влияют на этот факт.

Теорема 5: {\it утверждение $\aleph_0 = \omega$ независимо от ZFC.}

Доказательство:
Определим ординалы как индуктивные орды. Индуктивные уровни построения наборов можно пронумеровать ордами. Назовём {\it индуктивными} те наборы, которые находятся на уровне, соответствующем индуктивному орду. Остальные наборы будем называть {\it дедуктивными}. Теория, включающая только индуктивные наборы, в том числе ординалы, принципиально не описывает все дедуктивные наборы. Аксиома бесконечности в этом фрагменте принимает свою классическую форму. Остальные аксиомы ZFC также сохраняют силу при полном исключении дедуктивных наборов. Введение внутренних операций сводится к конкатенации типов вполне упорядоченности. Таким образом, мы получаем полную теорию ZFC, моделью которой служит индуктивный фрагмент игры $\{ L \mid R \}$. В этой модели утверждение $\aleph_0 = \omega$ опровергается теоремой Кантора. Следовательно, утверждение о тождестве внутренних и внешних пределов $\aleph_0 = \omega$ независимо от классической аксиоматики ZFC, поскольку существует непротиворечивая модель, в которой это тождество опровергается. Теорема доказана.

Таким образом, мы полностью подтвердили результат, установленный в статье "Несчётность множества континуума: шах и мат". Доказательство получилось весьма компактным и элегантным: разделение на индуктивные и дедуктивные наборы вводится в самой аксиоме бесконечности. При исключении из описания теории дедуктивных множеств система самостоятельно воспроизводит всю индуктивную иерархию ординалов на модели, где дедукция является неотъемлемой частью. В частности, $\dot{\omega} - 1$ оказывается определённым в модели, однако в теории обладает индуктивной морфологией, идентичной $\omega$, в полном соответствии с интуицией Кантора: $\omega - 1 = \omega$. Это красиво демонстрирует подводные камни ограничительного подхода к аксиоматике: аксиомы цензуры не позволяют увидеть полный контекст теории. В ZFC просто отсутствуют аксиомы, которые могли бы выявить скрытые дедуктивные объекты. При описании множеств невозможно обойтись без явной фиксации позиции референта. Кантор чётко осознал это обстоятельство в контексте возникновения эмерджентных свойств, связанных со структурной организацией объектов. Кардинальная арифметика соответствует чисто внутреннему взгляду и выражает классическое представление о натуральном ряде как о замкнутом относительно комбинаторных операций объекте. Однако такая арифметика даёт лишь наиболее общую характеристику бесконечного - неограниченную рефлексивность. Лишь дополнительные структурные ограничения позволяют сформировать из неё конкретный образ трансфинитного числа. Кантор исследовал структуру вполне упорядоченных множеств и нашёл в ней подтверждение своим интуициям. Порядковая арифметика явно сочетает внутренний и внешний взгляды через левые и правые операции соответственно. Довести эту идею до логического завершения удалось Джону Конвею, который предложил элегантную структуру, описывающую исключительно внешний взгляд на объекты. Тем самым Конвей завершил программу Кантора, сделав обе перспективы - внутреннюю и внешнюю - полноценными и равноправными инструментами анализа. Необходимость учёта позиции референта легко продемонстрировать на классическом парадоксе Гранд-отеля, предложенном Давидом Гильбертом. Рассмотрим отель с бесконечным числом комнат, в каждой из которых уже проживает постоялец. Обычно утверждается, что, несмотря на полную занятость, можно разместить нового гостя: достаточно переселить жильца из номера $n$ в номер $n + 1$, освободив тем самым первую комнату. Современная математика часто интерпретирует этот парадокс в актуальном и статическом смысле - через существование функции взаимно-однозначного соответствия $f(n) = n + 1$. Однако само существование такой функции как законченного объекта (множества упорядоченных пар) не является тривиальным фактом и требует обоснования через индуктивное построение или его эквивалент. Поэтому целесообразно рассмотреть задачу с чистого листа. Существуют два равноправных индуктивных рассуждения:
\begin{itemize}
	\item Гостя из первого номера можно переселить во второй. Если определено размещение гостей из первых $ n - 1$ номеров, то гостя из номера $n$ можно переместить в комнату $n + 1$, которая ещё не использовалась. Следовательно, для каждого гостя существует номер для переселения.
	\item Мы перемещаем гостя из первого номера во второй - задача остаётся нерешённой, поскольку один гость всё равно остаётся без номера. Если мы переселили гостей из $n - 1$ комнат, то после переселения оставшегося гостя в номер $n$ задача по-прежнему не решена - один постоялец остаётся без размещения. Следовательно, для всех комнат один гость остаётся не заселённым.
\end{itemize}

Эти два рассуждения раскрывают проблему с неожиданной глубиной. Не имеет принципиального значения, рассматриваем ли мы процесс или статическую конструкцию - оба подхода правомерны как в динамическом, так и в статическом контексте. В первом случае мы не можем говорить о переселении как о полностью завершённом факте: даже заявив, что гость заселён в номер $n + 1$, мы обязаны обеспечить возможность дальнейшего переселения в номер $n + 2$. Это приводит к простой смене обозначений и удерживает нас в рамках потенциальной возможности. Данный факт делает невозможной корректную переформулировку задачи для ''ответственного администратора'', который предпочитает не переселять уже проживающих постояльцев, а сразу разместить вновь прибывшего гостя в следующий номер. В этом случае, пройдя по всем комнатам, новый постоялец неизбежно останется без номера. Во втором рассуждении, напротив, мы вынуждены рассматривать переселение как завершённый акт. Здесь отель воспринимается как актуальная целостность, к которой применим принцип Дирихле: в какой-то комнате обязательно должно оказаться два постояльца. В первом подходе для нас существенно лишь сохранение структурного подобия отеля до и после операции, тогда как во втором - сохранение идентичности исходной и результирующей целостности. Первый случай развивает внутреннее стремление натурального ряда к собственному воспроизводству, второй фиксирует внешние ограничения для выхода из имеющейся формы. Таким образом, современная теория множеств, формально провозглашая работу с актуальной бесконечностью, на практике преимущественно оперирует существованием в возможности - то есть потенциальной бесконечностью. Ещё более примечателен следующий вывод. Интуиционизм и конструктивизм рассматривают классическую логику как ретроспективный инструмент, подходящий преимущественно для описания уже установленных связей, но малопригодный для анализа текущих или будущих процессов. Вероятно, именно поэтому применение классической логики к потенциальным событиям в рамках аксиоматических теорий множеств порождает столь значительное количество контринтуитивных следствий.

Позиция референта представляет собой одно из фундаментальных понятий человеческого мышления. Её следы обнаруживаются задолго до обращения к таким сложным объектам, как бесконечность. Она играет основополагающую роль в формировании ключевых математических понятий - множества, числа и преобразования. Сознание человека способно объективизировать коллекции, то есть переходить с внутренней позиции референта (рассматривающей коллекцию как систему связей между элементами) к внешней позиции, при которой сама коллекция становится самостоятельным объектом. Этот механизм широко используется в повседневном мышлении: пара глаз, вязанка дров, коллекция марок, корзина фруктов, тысяча рублей и т. п. Для такого перехода достаточно ввести модификатор множества - реальный (корзина, мешок) или абстрактный (коллекция, множество). Как только собрание элементов формируется как объект, становится возможным его описание через свойства, в частности, через количество элементов. То есть, количество всегда связано с внешней перспективой и выступает базовым числовым свойством коллекции как целостного объекта. Многие модификаторы множества задают не просто совокупность, а структурированную систему (пара, ряд, последовательность и т. д.). В таких объектах возникают дополнительные числовые свойства, прежде всего номера - порядковые числа. Уже на этом уровне находит подтверждение интуиция Георга Кантора: числа возникают как эмерджентные свойства структур объективизированных коллекций. Для описания направленных связей между объектами требуется абстракция преобразования. Преобразование, в свою очередь, также предполагает внешнюю позицию референта по отношению к связываемым объектам. Элементы множества могут быть связаны преобразованиями (например, отношением следования). Чтобы описать преобразование одного множества (например, натуральных чисел) в другое, необходимо занять позицию, внешнюю по отношению к обоим множествам. Именно здесь возникает характерная путаница при описании бесконечных коллекций, которые допускают потенциально неограниченное продолжение, но актуально обладают фиксированной структурой. В этом контексте полезно ввести различение {\it области регулярности} и {\it области сингулярности}. Под областью регулярности понимается индуктивно конструируемая связная система элементов, в которой между любыми двумя элементами существует потенциально конструируемый путь преобразований. Область регулярности соответствует сфере классической математики с её потенциальной бесконечностью - это система, доступная непосредственному описанию. Область сингулярности, напротив, представляет собой переход между областями регулярности - то, что принципиально не поддаётся полному описанию в терминах внутренней логики предыдущей области. Мы можем населять область сингулярности нестандартными бесконечно большими объектами, однако при попытке перейти от них к актуально бесконечному вновь сталкиваемся с принципиальной невозможностью описать сам переход. Когда сингулярность непосредственно предшествует выделенному элементу из области регулярности, например, как в случае предельные ординалов, возникает разрыв непрерывности предикативного описания. Такие выделенные элементы называются {\it точками разрыва предикативности}, но существуют морфологии позволяющие избежать такого разрыва. Здесь уместно вернутся к терминологии сечений Дедекинда - при разделении элементов связанных отношением порядка на два класса $A$ и $B$ могут реализоваться три возможности: оба класса $A$ и $B$ содержит граничные элементы - {\it скачок}; ни один класс не содержит граничных элементов - {\it пробел} и только один из классов $A$ и $B$ содержит граничный элемент - {\it предел}. Таким образом, можно связать точки разрыва предикативного описания со структурами реализующими предельное разделение.

Полезным оказывается ввести понятие {\it общезначности теории}: теория $T$ называется общезначной относительно своего предикативного домена $D$, если для любых двух моделей $M_1, M_2 \models T$ элементарные теории их выделенных базовых ядер совпадают: $\text{Th}(D(M_1)) = \text{Th}(D(M_2))$. Иными словами, теория общезначна тогда и только тогда, когда многообразие её моделей затрагивает только внешние надстройки, но не способно изменить внутренний логический статус базовых элементов. Если базовое утверждение теории сформулировано через квантор всеобщности $\forall x \in D \; \phi(x)$, оно работает как безусловное требование. Элемент $x$ не может попасть в домен $D$, не подчиняясь $\phi(x)$. Внешняя модель может быть сколь угодно нестандартной или содержать новые слои, но локальная семантика объектов $D$ строго фиксирована. Если фундаментальное понятие теории вводится через квантор существования $\exists X \; \psi(X)$ (как аксиома бесконечности в ZFC и NBG), то объект $X$ фиксирует лишь одну из альтернатив, реализация которых возможна в модели. В результате фундаментальные утверждения о свойствах элементов домена $D$ теряют общезначимость. Например, все базовые аксиомы $PA$ - это универсальные суждения ($\forall$). Хотя глобально существуют неизоморфные нестандартные модели $M_1 \neq M_2$, базовое ядро арифметики общезначно: локальные свойства каждого натурального числа одинаковы во всех моделях. В то же время базовые утверждения ZFC о фундаментальном объекте $\omega$ зависят от глобальных параметров модели. Существование нестандартных элементов, например, дедуктивных ордов, разрушающих индуктивную структуру объектов, или зависимость свойств континуума от форсинга означает, что $\text{Th}(\omega(M_1)) \neq \text{Th}(\omega(M_2))$. Теории ZFC и NBG не являются общезначными относительно собственного ядра. Так, например, если в модели реализуются альтернативные морфологии бесконечности - скажем, соответствующие положительным омнифическим целым числам, - то существование $\omega$ как минимального индуктивного замыкания становится справедливым исключительно для индуктивного фрагмента модели, связанного с аксиомой бесконечности, а не в общем случае. Аналогичным образом, утверждение о рефлексивности (бесконечности по Дедекинду) всякого бесконечного множества остаётся справедливым для фрагмента связанного с внутренними функциями, но не для всякого бесконечного множества вообще, поскольку допустимо вводить альтернативные функциональные отношения выразимые объектами теории. Это становится критически важным когда мы пытаемся рассуждать об общепринятом понятии натуральных чисел в контексте современных аксиоматик, объекте который отсутствует в общезначном смысле в самих теориях, что становится причиной недопонимания между математиками. Далее мы рассмотрим некоторые артефакты отсутствия общезначности более подробно. От них легко избавиться, постулировав, например, что всякое множество в ZFC и NBG является индуктивным. Именно с отсутствием общезначности связаны существенные трудности, возникающие при работе со стандартными аксиоматиками теории множеств. Данное наблюдение позволяет нам сформулировать ключевое требование для построения достоверных аксиоматических систем: всякая теория должна оставаться открытой для расширения (дополнения), как $PA$, но закрытой для изменения, как в случае ZFC и NBG, допускающих расслоение смысла своих базовых примитивов, в частности отношения принадлежности $\in$. 

Отметим также, что общезначность можно попытаться определить в рамках современного понятийного аппарата, например как {\it консервативную категоричность}: теория называется консервативно категоричной, если она допускает исключительно консервативные расширения. Однако такое определение наследует относительность понятий категоричности и консервативности, которые как раз связаны с отсутствием общезначности в метатеории. Например, расширение арифметики $PA + \text{Con}(PA)$ не является консервативным синтаксически; тем не менее, никак не влияя на элементы домена, оно позволяет теории формализовать утверждения о своём собственном домене, которые, вообще говоря, следует относить к метатеории. Иными словами, консервативность расширения в современном понимании подразумевает отсутствие метатеоретических расширений, которые не влияют на исходные элементы домена. Фактически это означает, что понятие консервативности, как и категоричности, несёт в себе знание о внешней интерпретации теории. А поскольку не существует единой глобальной теории, такие понятия становятся размытыми философскими мета-абстракциями, с которыми крайне сложно работать. Лучше всего воздерживаться от понятий смешивающих позиции референта, до тех пор пока мы не в состоянии указать на необходимость такого смешения и очертить его границы. На отсутствие общезначности при оперировании с бесконечностью очень точно указал Брауэр, выявив одно из ключевых следствий данного феномена - нарушение закона исключённого третьего. Поскольку различные альтернативы выражаемые через единый синтаксис равноправны, они не исключают друг друга, а сосуществуют в едином многозначном пространстве возможностей, что мы и интерпретируем как нарушение данного логического закона. Дело в том, что сама по себе логика первого порядка FOL полна, что гарантирует истинность всякого выводимого утверждения в любой модели. Действительно, если мы корректно применяем FOL как инструмент к области, где любое пространство альтернатив строго ограничено двумя значениями, то мы имеем тотальный изоморфизм синтаксиса и семантики в такой теории. Однако, FOL + ZFC или FOL + NBG не имеют никаких методологических оснований считать все пространства альтернатив данных теорий двухзначными. Действительно, аксиома бесконечность здесь разрушает ключевое условие применимости FOL - она утверждает существование объекта в пространстве с неизвестным количеством альтернатив. Это означает, что любые выводы связанные с данной аксиомой, любую ветвь в дереве вывода, которая начинается в данном корне, мы обязаны описывать в минимальной логике Йохансона, которая не предполагает никакого априорного знания о количестве альтернатив. Более того, о том что синтаксис и семантика в нашей теории разъехались свидетельствует теорема Гёделя о неполноте, некоторые семантические истины заведомо синтаксически невыводимы в теории. Об этом же свидетельствует теорема Сколема, которая говорит, что отношение принадлежности вообще говоря семантически неоднозначно. Формально данный эффект абсолютно понятен, если мы добавим аксиому индуктивно-дедуктивной бесконечности в ZFC или NBG, в стиле сюрреальных чисел или гросс-объектов мы получим синтаксическое противоречие: для индуктивных множеств $\bigcup \omega \notin \omega$, в то время как для индуктивно-дедуктивных множеств $\bigcup \omega \in \omega$. Но семантически данное противоречие не выводимо в теории, когда в корне дерева вывода индуктивно-дедуктивная бесконечность выводимо $\bigcup \omega \in \omega$, а когда индуктивная - $\bigcup \omega \notin \omega$, однако, обе альтернативы не могут входить в одно и то же дерево, поскольку они взаимоисключающие. Закон исключённого третьего здесь принимает обобщённую форму: множество либо конечно, либо индуктивно бесконечно, либо индуктивно-дедуктивно бесконечно, делая классическую логику неприменимой. Это не вопрос принятия или непринятия аксиом логики, поскольку их принятие не гарантирует изоморфизма синтаксиса и семантики. Приняв аксиомы ZFC или NBG и аксиомы FOL мы просто получаем теорию, которая может оказаться синтаксически противоречивой, но семантически непротиворечивой. Такая теория, может иметь модель, например, нестандартную теория, которая различает индуктивный и индуктивно-дедуктивный синтаксис, в результате чего в модели данное противоречие становится синтаксически невыводимым. Однако, теория объединяет разные синтаксические конструкции модели в единое отношение принадлежности, что и выглядит как синтаксическая противоречивость. Как видим формальная ошибка ложной дихотомии, несмотря на то, что она является методологической, а не внутренней ошибкой теории, имеет чрезвычайно масштабные последствия на интерпретацию теоретических результатов.

Рассмотрим взаимно однозначное соответствие между натуральными числами и конечными начальными отрезками натурального ряда, определяя при этом натуральное число как орд (или как омнифическое целое сюрреальное число), а множества как индуктивные наборы, чтобы избежать терминологической неоднозначности:

\begin{align}
\begin{split}
0 &\Leftrightarrow N_0 = \{\}\\
1 &\Leftrightarrow N_1 = \{1\}\\
2 &\Leftrightarrow N_2 = \{1, 2\}\\
\dots &\Leftrightarrow \dots\\
i &\Leftrightarrow N_i = \{1, 2, 3, \dots, i\}\\
\dots &\Leftrightarrow \dots
\end{split}
\end{align}

Данная конструкция наглядно иллюстрирует принципиальное различие между подходами, опирающимися на потенциальную и актуальную бесконечность. В рамках потенциально-бесконечного подхода мы можем сколь угодно долго продолжать процедуру установления взаимно однозначного соответствия. Однако этот процесс никогда не приводит к числу, которое соответствовало бы всему множеству натуральных чисел. Предположение, что множество $\mathbb{N}$ присутствует где-то в данной последовательности (скажем, на месте $n$), немедленно приводит к противоречию: за элементом $n$ следует $n + 1$, который по определению также должен принадлежать $\mathbb{N}$, но не может содержаться в $N_n$. Таким образом, не существует конечной процедуры, позволяющей построить множество всех натуральных чисел. В рамках потенциальной бесконечности такого актуального объекта просто не существует. Тем не менее, указанная система образует область регулярности, обладающую внутренней связностью: все комбинаторные преобразования остаются замкнутыми внутри неё, формируя алгебраическую структуру кольца. Однако наше мышление позволяет обобщить потенциально бесконечный опыт и говорить о справедливости утверждений для всех натуральных чисел, более того, мы даже можем говорить о существовании такого объекта как натуральный ряд. Иными словами, мы всё равно способны мыслить потенциально бесконечный процесс как целостный объект, но для того, чтобы сохранить его замкнутость, нам необходимо ограничить его внешнюю предикативность, не позволяя диагональной конструкции получить элемент $\aleph_0$ отличный от всех с помощью алгеброобразующей операции сложения. Для этого мы принимаем предположение замкнутости алгебры, которое оставляет полученный диагональным построением элемент частью замкнутой структуры, добавление этого элемента в исходное множество оставляет множество неизменным $\aleph_0 + 1 = \aleph_0$. Таким образом, мы получаем абсолютно бесконечный объект, который импредикативен внешне, к нему невозможно добавить новый отличных всех элемент с помощью внутренней операции.

В то же время, мы можем отказаться от идеи замкнутости и волевым усилием определить множество всех натуральных чисел как результат бесконечного продолжения вышеописанной процедуры или как её предел: $\infty = \lim_{i = 0, 1, 2, \dots} i \Leftrightarrow \mathbb{N} = N_{\infty} = \lim_{i = 0, 1, 2, \dots} N_i$. Цена такого определения заключается в том, что $\infty \in \mathbb{N}$. Иными словами, существование бесконечного множества натуральных чисел как объекта эквивалентно существованию бесконечного элемента, которая сама включена в это множество и характеризует его упорядоченную структуру. Для большей наглядности можно каждой паре $(i, N_i)$ поставить в соответствие вещественное число, имеющее ровно $i$ значащих цифр после запятой и совпадающее с начальными цифрами разложения $\pi$. Поскольку число $\pi$ существует, мы должны признать существование бесконечного множества значащих позиций в его десятичной записи, а также натурального числа $\infty$, соответствующего всей этой последовательности. Такую бесконечность естественно назвать натуральной бесконечностью, поскольку она выражает идею совершенной завершённости результата бесконечной счётной процедуры - как в случае позиций $\pi$, так и в случае элементов $\mathbb{N}$. Подобная абстракция идеала (достигнутой завершённой цели) позволяет строго определить понятие все элементы множества, одним из которых оказывается бесконечность. Это следует из того, что операция сложения является абстракцией объединения: если объединение объектов порождает бесконечное множество, то ему должно соответствовать натуральное число, выражающее мощность данного множества, - что является прямым следствием определения числа. Диагональная конструкция тут же подхватывает такую ментальную конструкцию позволяя получить элемент $\infty + 1$. Мы получаем внешнее предикативное описание бесконечных объектов. Абстрагирование множества - это уже внешняя операция, поэтому нам проще мыслить данный процесс последовательно, сначала мы выделяем множество, а потом понимаем, что множество само является объектом. Что позволяет скрыть бесконечный элемент внутри бесконечного множества, однако сути это не меняет. Предикативность бесконечного множества как объекта обеспечивается внешними функциональными отношениями, для которых диагональная конструкция является интуитивной онтологической поддержкой. Кантор чётко улавливает эту разницу и вводит в порядковой алгебре два независимых типа операций: левые (внутренние, кардинальные) и правые (внешние, индексные). Левые связаны с потенциальной бесконечностью, то есть внешней импредикативностью объекта. В то время как правые связаны с актуальной бесконечностью, когда объект становится предикативным. Это позволяет сформулировать диагональный аргумент для внешних операций, который как мы видели в статье "Несчётность множества континуум, шах и мат" допускает обобщение. Так при унарном кодировании диагональная конструкция позволяет определить элемент отличный от всех, тем самым обеспечивая упорядоченную структуру трансфинитных порядковых чисел. Диагональная конструкция таким образом однозначно ассоциируется с внешним контекстом. 

Стоит отметить, что введение понятия бесконечного позволяет строго демаркировать понятие конечного. Само по себе понятие конечности не может быть определено без порочного круга. Именно поэтому Дедекинд и Кантор определяют бесконечное множество через свойство подобия собственной части, а конечное - как отрицание этого свойства. Бесконечность, таким образом, возникает в результате сознательного постулата возможности внешней позиции референта по отношению ко всей совокупности конечных объектов (по сути - позиции Бога в классическом смысле). Такой подход оставляет область, недоступную полному математическому описанию - область сингулярности, отделяющую конечное от новых мнимых объектов. Примечательно, что концепция мнимых объектов, успешно применённая для определения отрицательных чисел ($x + (-1)y$) и комплексных чисел ($x + iy$), при переходе к бесконечности ($x + \dot{\omega} y$) сталкивается со значительными трудностями. Кантор последовательно сохраняет принцип работы с мнимыми объектами, считая разность легитимной числовой операцией. В частности, он принимает $\omega - 1 = \omega$, что позволяет записать $\omega = \{0, 1, 2, \dots, \dot{\omega} - 1 \}$, где порядковый тип множества $\dot{\omega} - 1$ есть $\omega - 1 = \omega$, подчёркивая тем самым определение порядкового типа $\omega$ как неподвижной точки внешнего упорядочения. В рамках ZFC такая онтология в значительной степени утрачивается. Для порядковых типов мы не можем говорить о максимальном порядковом числе внутри $\omega$, поскольку именно это число и есть $\omega$ (в смысле порядкового типа). Кантор неслучайно проводит строгое различие между множеством и типом множества (классом эквивалентности). Это различие позволяет рассматривать $\omega$ как минимальный объект аннотации (тип), а не как собственно множество. Множество $\dot{\omega} - 1$ в порядковом смысле строго меньше $\dot{\omega}$, хотя их индуктивная структура полностью совпадает и описывается одним и тем же порядковым типом $\omega$.

Однако при переходе к аксиоматическому подходу это представление становится проблематичным. Мы заменяем порядковые типы их каноническими представителями - ординалами. Хотя сам Джон фон Нейман подчёркивал, что корректно говорить именно о типах множеств, а описание в терминах ординалов осмысленно лишь в контексте уже заданной порядковой арифметики (для которой и вводятся канонические представители) \cite{von-neumann-ordinals}. Для полного воспроизведения порядковой арифметики на ординалах необходимо сначала определить классы эквивалентности ординалов по индуктивной структуре, а затем сопоставить каждому такому классу канонический элемент. Орды $\dot{\omega} - 1$ и $\dot{\omega}$ представляют различные объекты на уровне модели (где ординал - это конкретное множество), но соответствуют одной и той же индуктивной структуре на уровне теории, обозначаемой как $\omega$ (где ординалы понимаются как порядковые типы). К сожалению, эта сознательная реификация порядковых типов не была в должной мере осмыслена математическим сообществом. В результате $\omega$ стало восприниматься именно как минимальное множество, а не как тип. Это, с одной стороны, блокирует исходную онтологию Кантора, затрудняя понимание, а с другой - формализует абстракцию порядковых типов как самостоятельных сущностей, отсекая всё лишнее. Тем не менее, данный сдвиг ещё не приводит к отождествлению $\aleph_0 = \omega$ в сильном смысле и к утверждению, что свойства ординалов полностью определяются их индуктивной структурой, замкнутость которой означает $\omega = 2^{\omega}$ как в кардинальной арифметике $\aleph_0 = 2^{\aleph_0}$.

Развитая Кантором порядковая арифметика с введением натуральных операций Хессенберга позволяет стандартным образом построить кольцо частных: $\frac{\alpha_1}{\beta_1} = \frac{\alpha_2}{\beta_2} \Leftrightarrow \alpha_1\otimes\beta_2 = \beta_1\otimes\alpha_2$ c операциями сложения и умножения: $\frac{\alpha_1}{\beta_1} + \frac{\alpha_2}{\beta_2} = \frac{\alpha_1\beta_2 \oplus \beta_1\alpha_2}{\beta_1\beta_2}$, $\frac{\alpha_1}{\beta_1} * \frac{\alpha_2}{\beta_2} = \frac{\alpha_1\otimes\alpha_2}{\beta_1\otimes\beta_2}$. В частности, это позволяет рассматривать такие частные, как $\omega/2$. На уровне модели соответствующее множество представляет собой половину $\omega$, однако на уровне аннотации (порядковых типов) оно сохраняет тип $\omega$. Таким образом, область сингулярности - это не просто переходная зона, а полноценная реальность, населённая на уровне модели объектами типа $\dot{\omega}/2$, $\dot{\omega} - 5$, $\dot{\omega}/7$ и так далее. При аксиоматизации теории множеств мы корректно фиксируем внутреннюю часть - индуктивные отношения между элементами - посредством схем аксиом выделения и замещения. Это гарантирует существование внутреннего описания: в теории появляются объекты, представляющие системы связей между элементами. В терминах таких индуктивных отношений множества действительно самоподобны. Например, отношение следования $s(n) = n + 1$ переводит $\omega$ в $\omega \setminus \{\emptyset\}$ во внутреннем смысле. На уровне модели набор пар $f = \{(0, 1), (1 ,2), \dots, (\dot{\omega} - 2, \dot{\omega} - 1) \}$ задаёт функцию $f : \dot{\omega} - 1 \to \dot{\omega} \setminus \{\emptyset\}$, которая на уровне индуктивной аннотации воспринимается как $f : \omega \to \omega \setminus \{\emptyset\}$. Аксиоматика теории множеств при этом не запрещает существование других объектов, описывающих внешнее отношение, при котором $\omega$ отображается в $\omega + 1$, фактически мы должны лишь каким-то образом зафиксировать выделенную пару $(\omega, \omega)$. Иными словами, она допускает как внутреннее индуктивное представление, так и внешнее целостное представление - при соответствующем расширении аксиом. Тем не менее, многие построения, основанные на сопоставлении чисел и начальных сегментов натурального ряда, часто не находят удовлетворительного объяснения в господствующих интерпретациях и подвергаются критике. Авторы таких рассуждений, как правило, не являются специалистами в теории множеств и нередко смешивают рассуждения на уровне модели с рассуждениями в рамках теории. Однако они и не обязаны владеть всей строгостью аппарата: они формулируют общественный запрос на полную и непротиворечивую картину бесконечности. Математическое сообщество должно было формализовать и перевести этот запрос на строгий математический язык, однако вместо этого зачастую игнорирует его.

Но мы идём ещё дальше. Мы стремимся к тому, чтобы порядковые типы, которые Кантор мыслил как дуальные объекты, сочетающие внутреннее индуктивное описание с внешним представлением, описывались исключительно через внутренние индуктивные отношения. Мы отказываемся от внешних пределов Кантора, задаваемых через ряды, и вводим пределы как супремумы внутренних индуктивных конструкций. При этом возникновение промежуточных объектов (таких, как множество всех функций с конечным носителем) никак не аргументируется. Формально требуется либо доказать их существование, либо явно постулировать соответствующую аксиому. Если объединение всех конечных ординалов даёт $\bigcup_{n < \omega} n = \omega$, то объединение всех конечных носителей даёт нечто иное, несмотря на то, что последовательность конечных носителей, в которой каждый элемент пробегает все возможные значения, эквивалентна последовательности конечных ординалов. Данный шаг представляется по меньшей мере неочевидным. Кантор целенаправленно вынес полностью индуктивное описание в кардинальную алгебру, где внутреннее описание не смешивается с внешней фиксацией структуры. Для аксиоматического описания кардинальной алгебры могут быть использованы новые основания Куайна, позволяющие непротиворечиво замкнуть универсум относительно любых индуктивных преобразований. В порядковой алгебре, несмотря на использование внутреннего описания для сложения и умножения (необходимого для обеспечения самоподобия объектов и их отличия от конечных), Кантор сохраняет исключительно внешнее описание для операции возведения в степень. Степень определяется через внешнее умножение и внешнюю сумму ряда. Именно благодаря этой внешней составляющей в упорядоченном универсуме работает теорема Кантора, позволяющая изолировать области, между которыми не существует биекции сохраняющей порядок. Это даёт возможность выстроить унифицированную иерархию бесконечных объектов. Как мы увидим далее, ключ к построениям Кантора заключается не столько в самой структуре, сколько в способности этой структуры поддерживать унификацию, которая вводит эталон измерения бесконечности - фактически, линейку для бесконечных объектов. Тем не менее, мы отказываемся от аппарата Кантора, вводим альтернативный формализм и определяем кардиналы через ординалы, требуя, чтобы они вели себя изоморфно кардиналам Кантора. Однако минимальный ординал, с которым множество находится во взаимно однозначном соответствии, и класс эквивалентности множеств, отображаемых друг на друга взаимно однозначно (без каких-либо дополнительных структурных требований к отображению), - это принципиально разные объекты. Таким образом, мы без достаточных оснований пытаемся распространить внутреннее индуктивное поведение кардиналов на ординалы, которые как порядковые типы обладают существенной внешней составляющей. Лишь аксиома выбора в современных аксиоматических системах в определённой мере сохраняет внешнюю интуицию Кантора. Кантор действительно полагал, что любое множество можно вполне упорядочить, но при этом считал, что мы создаём новый объект - упорядоченную аннотацию, описываемую в терминах порядковой алгебры, в то время как исходное множество само по себе не является объектом этой алгебры. Иными словами, аксиома выбора выступает мостиком между кардинальной и порядковой алгебрами: для любого произвольного множества мы можем сконструировать структурированный объект, морфология которого выявляет новые свойства в рамках нового алгебраического контекста.

Именно безосновательное включение в аксиоматики тождества $\aleph_0 = \omega$ (где под $\aleph_0$ понимается абстракция внутренней индуктивной морфологии, то есть принципиальная возможность взаимно однозначного сопоставления, не ограниченная рамками упорядоченного универсума) стало одной из главных причин критики аксиоматических теорий множеств. В частности, Анри Пуанкаре прямо писал: ''всякая теорема, относящаяся к бесконечным числам или вообще к тому, что называется бесконечным множеством, или трансфинитным количественным, или трансфинитным порядковым и т. д., не может быть чем-либо иным, как сокращённым способом формулирования предложений, относящихся к конечным числам. Если дело обстоит иначе, то эта теорема недоказуема, а если она недоказуема, то она не будет иметь смысла'' \cite{poincare-infinity}. Пуанкаре ни в коем случае не отрицал саму идею бесконечности. Он лишь настаивал, что изучение её свойств и закономерностей следует начинать с области конечных величин и тщательно подготовить понятийный аппарат, прежде чем вводить такую сущность в математику. Удивляет математическая интуиция Пуанкаре, который, несмотря на то, что не смог точно указать на фундаментальные проблемы оснований аксиоматических теорий множеств, чрезвычайно точно связал эти трудности с апориями Зенона: ''случилось так, что натолкнулись на некоторые парадоксы, которые привели бы в восторг Зенона Элейского и мегарскую школу'' \cite{poincare-zeno}. При этом Пуанкаре длительное время был сторонником теории множеств и даже участвовал в переводе на французский язык ключевых работ Кантора. К сожалению, Пуанкаре не всегда был достаточно сдержан и называл патологией направление аксиоматизации теории множеств, которое он именовал ''Канторизмом''. Судя по всему, он намекал и на личную трагедию Кантора. Между тем, сам Кантор, по всей вероятности, был бы во многом солидарен с Пуанкаре, в частности, в том, что попытка слить внутреннюю и внешнюю перспективы в едином понятии числа является ошибочной и именно она порождает парадоксы, столь напоминающие апории Зенона.

Пусть задана фундаментальная последовательность $a_n$ с пределом $a$. Если последовательность строго монотонна, то она не может содержать свой предел в качестве элемента: противное предположение приводит к противоречию, поскольку невозможно указать порядковый номер такого элемента. Однако с позиций теории множеств ничто не препятствует объединению множества элементов последовательности $\{a_n\}$ с одноэлементным множеством $\{a\}$. В аксиоматике Цермело–Френкеля \cite{fraenkel-axioms} для этого достаточно аксиомы объединения, аналогично в аксиоматике фон Неймана - Бернайса - Гёделя \cite{von-neumann-axioms} и в любой разумной аксиоматике теории множеств. Полученное множество $\{a_n\} \cup \{a\}$ должно быть счётным, поскольку добавление одного элемента не изменяет мощность бесконечного множества. Таким образом, оно представимо в виде последовательности. Нюанс с порядком элементов не спасает ситуацию. Если мы не можем потребовать конкретного порядка, то не можем и запретить его наличие. Кроме того, любое счётное множество можно упорядочить с помощью счётного числа операций, а сам Кантор считал, что порядок при установлении взаимно однозначного соответствия в кардинальной арифметике несущественен \cite{cantor-relation}. Следовательно, теория множеств не содержит средств, запрещающих существование строго монотонной последовательности, содержащей свой предел. Номер такого предельного элемента невозможно указать, не перейдя к внешней позиции референта, позволяющей оперировать бесконечными числами. Данный факт подчёркивает решающую роль структуры в теоремах Кантора о несчётности. Мы говорим об упорядоченном множестве, изоморфном множеству конечных порядковых чисел. Приписывая множеству порядковый тип $\omega$, мы указываем, что неважно, о каком именно представителе в модели идёт речь - $\dot{\omega}-5$, $\dot{\omega}^2$ или любом другом. Каждое из них может быть описано как $\omega$, и ни одно из них при таком описании не может быть поставлено во взаимно однозначное соответствие со множеством всех своих подмножеств. Дело не в конечности индексов, а в самой линейной структуре, которая не допускает этого, если мы способны унифицировано описывать все множества данной совокупности через единый эталон.

Как ни парадоксально, именно об этом говорил Аристотель: ''оно (бесконечно) делимо и при уменьшении и обратном прибавлении, а целым и ограниченным [бесконечное] оказывается не само по себе, а по отношению к другому'' \cite{aristotle-limit}. Две тысячи лет назад он уже различал внешнюю перспективу, необходимую для сравнения, и внутренний взгляд, определяющий морфологию связей между элементами. Именно поэтому для $\omega$ критически важна внешняя составляющая: она обеспечивает унификацию всех множеств в одной области длин. Все порядковые типы индуктивно подобны в смысле вполне упорядоченной структуры, однако именно внешний аспект позволяет описывать группы таких типов через единый эталон. Без предположения о внешней перспективе (возможности сравнения) мы могли бы аннотировать множества $\dot{\omega}$ и $2^{\dot{\omega}}$ одним и тем же порядковым типом $\omega$, что сделало бы теорему Кантора недоказуемой (аналогично неявному принятию $\aleph_0 = \omega$). То есть, в онтологии Кантора неявно присутствует постулат внешнего описания позволяющий сравнивать упорядоченные отрезки. Интересно, что даже математики, принявшие основные положения теории множеств, несмотря на их контринтуитивность, такие как Гильберт и Бернайс, ясно видели этот артефакт, хотя, по-видимому, не вполне осознавали его природу: ''парадокс, заключающийся в том, что некая бесконечная последовательность следующих друг за другом событий, последовательность, завершаемость которой мы не можем себе даже представить (не только фактически, но хотя бы в принципе), на самом деле все-таки должна завершиться'' \cite{gilbert-paradox}. Отсюда до формального различения внутренней и внешней перспектив, ответственных за потенциальную и актуальную бесконечность, остаётся один шаг. Ситуацию можно сравнить с корпускулярно-волновым дуализмом в физике: оба описания корректны в своих рамках, но сознание в каждый момент времени способно удерживать лишь одно из них. Работая с трансфинитными последовательностями, мы преимущественно опираемся на внутренний индуктивный взгляд, связанный с потенциальной бесконечностью и позволяющий продолжать построение неограниченно. Актуальная же бесконечность необходима для сравнения упорядоченных множеств между собой и для их совместного описания через единый эталон.

Эти размышления приводят нас к ещё одному наблюдению, связывающему апории Зенона с трудностями аксиоматических теорий множеств: и те, и другие коренятся в конфликте внутренней и внешней позиций референта. Аристотель элегантно разрешает первую апорию Зенона - дихотомию: ''И вот, бесконечного в количественном отношении нельзя коснуться в конечное время, а бесконечного в отношении деления - можно, так как само время бесконечно именно в таком смысле. Таким образом, бесконечное удается пройти в бесконечное, а не в конечное время и коснуться бесконечного [множества можно] бесконечным, а не конечным [множеством]'' \cite{aristotle-dichotomy}. Он рассматривает проблему с внешней перспективы: расстояние описывается рядом $l/2 + l/4 + \dots + l/2^n + \dots = l$, а время - аналогичным рядом $t/2 + t/4 + \dots + t/2^n + \dots = t$. Бесконечного множества расстояний можно ''коснуться'' бесконечным множеством временных промежутков. Такая интерпретация очень близка современным представлениям, где ''касание'' одного бесконечного множества другим мы называем взаимно однозначным соответствием. По сути, Аристотель, выражаясь современным языком, строит изоморфизм двух упорядоченных множеств - подход, который удалось переоткрыть лишь два тысячелетия спустя. Тем не менее, понимание Аристотеля глубже того, что обычно доступно в современной трактовке. У него ''касание'' множеств происходит не за счёт внутренней организации самих последовательностей, а благодаря тому, что обе последовательности возникают из единого описания фиксированного физического процесса. Требование общей системы координат - пространства, в которое с внешней позиции вписываются сравниваемые множества, - является ключевым для корректной работы с множествами как целостными сущностями. Собственно, гениальность Кантора состояла в том, что он предпринял первую систематическую попытку построить такую систему координат для бесконечных упорядоченных множеств. К сожалению, этот эпистемологический прорыв остался в значительной мере невостребованным в процессе аксиоматизации теории множеств.

Аристотель не останавливается на этом и точно указывает источник возникающего ощущения противоречия: ''Если же их сделать действительными, то [движение] не будет непрерывным, но будет останавливаться, что вполне очевидно произойдет с тем, кто считает половины; ведь ему тогда необходимо одну точку считать за две: одна будет концом одной половины, другая - началом другой, если считать непрерывную [линию] не как одну, а как две половинные''. Когда мы начинаем мыслить половины, мы разбиваем непрерывное расстояние на дискретный ряд отрезков, которые перебираем пошагово. С внутренней точки зрения данной процедуры завершение процесса невозможно - он принципиально неограничен. Аристотель заключает: ''Таким образом, на вопрос, можно ли пройти бесконечное число [частей] во времени или по длине, следует ответить, что в одном отношении можно, в другом нет. Если они будут существовать в действительности - нельзя, если в возможности - можно, так как [предмет], движущийся непрерывно, прошел бесконечное множество по совпадению, а не прямо, ибо наличие бесконечного числа половин в линии есть для нее побочное обстоятельство, а сущность ее и бытие иные'' \cite{aristotle-perspectives}. Иными словами, акт рефлексии позволяет взглянуть на процесс собственного зацикливания внутри бесконечной последовательности со стороны и вписать этот алгоритм как целостность в более широкое наблюдаемое пространство. Аристотель выражает это категорией ''в возможности'': мы сохраняем возможность неограниченного продолжения процесса и одновременно помещаем его в общий контекст, где он становится лишь побочным обстоятельством, а не сущностью. Таким образом, внешняя позиция референта позволяет наблюдать внутренне бесконечный процесс как законченную целостность в рамках системы координат, связанной с описываемым явлением. Эта система координат задаёт внешнюю биекцию между отрезками расстояния и интервалами времени, где количественная тождественность гарантируется общностью исходной конструкции, а не внутренней морфологией последовательностей. Аристотель не просто разрешает апории Зенона - он задаёт метатеоретическую рамку, разделяющую ключевые режимы работы человеческого мышления. Он показывает, что большинство затруднений возникает именно из-за неспособности одновременно удерживать оба режима. Более того, он предлагает аналитический аппарат переключения между перспективами, позволяющий получить унифицированное описание как внутреннего представления ''в возможности'', так и формируемой им целостности в исследуемом контексте.

Оглядываясь назад и анализируя доказательства несчётности множества всех подмножеств счётных множеств, трудно поверить, что математическое сообщество поддалось апориям бесконечности и не разглядело дуализм позиций референта, о котором уже писал Аристотель. Более того, человек, которого принято считать основоположником теории множеств, - Георг Кантор - в своих зрелых работах чётко разделяет общее понятие мощности и тип порядка, явно избегая отождествления теоремы о невозможности биекции между множеством и множеством его подмножеств с кардинальной арифметикой. Он не совершает формальной ошибки ложной дихотомии, в каждой из независимых алгебр пространство всех альтернатив остаётся двухзначным обеспечивая применимость ключевого инструмента математического описания - логики. Судя по всему, Кантор уже осознавал, что одна из его первых теорем - о счётности множества алгебраических чисел, сформулированная в рамках кардинальной алгебры ($|\omega| = |\omega^\omega| = \aleph_0$), - соседствует с теоремой о несчётности множества вещественных чисел, принадлежащей порядковой алгебре ($\omega < 2^\omega$). Для второго доказательства требуется общая система координат, обеспечивающая возможность внешней биекции, на необходимость которой указывал Аристотель. Кантор тонко чувствовал это ограничение, поэтому в своём итоговом труде осознанно разделил кардинальную и порядковую арифметику. В таком подходе функция в алгебре кардиналов $f_c(x)$ и функция в алгебре порядковых чисел $f_o(x)$ представляют собой принципиально разные объекты. $f_c(x)$ - это неупорядоченное множество пар, тогда как порядок $f_o(x)$ индуцируется областью определения и областью значений. В частности, в упорядоченном универсуме $f_o(x)$ является множеством со строго заданным рангом, в то время как $f_c(x)$ может быть классом, не обладающим структурой множества. Смешение этих концепций наиболее корректно связывать с работами Герхарда Хессенберга, который одним из первых обобщил доказательство Кантора на произвольные (не обязательно структурированные) множества \cite{hessenberg}. Хессенберг не акцентировал внимание на различии между внутренней природой $f_c(x)$ и смешанной (внутренне-внешней) природой $f_o(x)$. В результате внешняя составляющая $f_o(x)$, обеспечивавшая корректность оригинального доказательства Кантора, была распространена как требование на все множества. Это положение было унаследовано аксиоматической теорией множеств и стало одной из глубинных причин кризиса оснований математики. История недоразумений, связанных с проблемой несчётности, напоминает детектив с запутанным сюжетом: длинная цепочка неочевидных компромиссов, допущенных разными авторами, привела к созданию дефектного фундамента современной математики. Но раз уж мы взялись за это расследование, следует довести его до конца.

В своей первой работе ''Об одном свойстве совокупности всех действительных алгебраических чисел'' \cite{cantor-algebraic} Кантор доказывает, что множество алгебраических чисел можно поставить во взаимно однозначное соответствие с натуральными числами, тем самым допуская существование бесконечной последовательности, перечисляющей все алгебраические числа. В той же статье он приводит пример множества, которое невозможно взаимно однозначно отобразить на натуральные числа, - множество вещественных чисел. Работа ещё не содержит понятий мощности, кардинальности или порядка. Оба доказательства корректны при аккуратной формулировке исходных предположений, что создавало серьёзный вызов для самого Кантора, осознававшего важность сделанного открытия. С современных позиций находка Кантора особенно примечательна: она иллюстрирует контринтуитивный факт, что одно множество с экспоненциальной комбинаторной сложностью может быть представлено в виде последовательности, а другое - нет. При этом многие модели множества вещественных чисел, в том числе современные, связаны с экспоненциальнымы конструкциями, такими как булеан счётного множества позиций, которые однозначно сопоставимы с подмножеством множества всех многочленов с коэффициентами 0 и 1 (в том числе, при рассмотрении с внешней позиции, многочлены бесконечных степеней). Математическая интуиция учителей Кантора - Карла Вейерштрасса и Леопольда Кронекера - сигнализировала об опасности: они высказывали серьёзные сомнения в корректности формулировок, однако не смогли точно указать источник диссонанса \cite{weierstrass-kronecker}. В глаза бросается противопоставление потенциальной и актуальной бесконечности. Множество алгебраических чисел определяется как потенциально бесконечное (конечные многочлены с конечными коэффициентами), тогда как множество вещественных чисел содержит числа с актуально бесконечной позиционной записью. Именно это поверхностное наблюдение, по-видимому, сыграло свою роль в формировании современного представления о промежуточных множествах всех функций с конечным носителем, хотя такие объекты отсутствуют в работах самого Кантора. Сегодня мы утверждаем, что, объединяя все комбинации одного, двух, трёх и так далее объектов, невозможно получить несчётное множество, несмотря на то, что само множество объектов в таком построении становится бесконечным. Любопытно, что это утверждение до сих пор не имеет строгого обоснования. Оно далеко не очевидно, однако принимается на веру без сколько-нибудь убедительных доказательств.

Действительно, говоря о потенциальной бесконечности, мы подразумеваем лишь отсутствие ограничения на продолжение процедуры. Для натуральных чисел не существует максимального числа; точно так же не существует максимального количества счётных множеств, после которого их объединение становится несчётным, или максимальной степени многочленов, определяющих алгебраические числа. Этот принцип лежит в основе математической индукции и позволяет делать утверждения, справедливые для любого натурального числа, а вслед за Кантором - и для всего счётного множества натуральных чисел. Поэтому разделение на любое конечное, сколь угодно большое и актуально бесконечное в значительной мере условно: через первое мы по определению вводим второе. Начиная объединять множества многочленов, полученные для разных высот, мы не можем указать конечную высоту, на которой это объединение становится несчётным. Следовательно, объединение по всему счётному множеству высот остаётся счётным и содержит многочлены произвольной (неограниченной) длины. Это легко показать: взяв любую неограниченную последовательность целых чисел и пошагово исключая многочлены, отличающиеся от неё коэффициентами, мы никогда не исключим все многочлены за конечное число шагов. Сама невозможность установить различие за конечное число шагов и есть признак бесконечности. При этом от нашего внимания ускользает то, о чём говорил Аристотель: для сравнения бесконечных множеств требуется их ''касание'' по совпадению - биекция, учитывающая общую пространственную структуру, в которую вписываются объекты. Метод, открытый Кантором, вовсе не выявляет несуществующую границу между потенциально и актуально бесконечным; он описывает объекты с внешней позиции, помещая их в общую систему координат, что делает возможным их корректное сравнение. Именно поэтому современная аргументация различия между случаями потенциально бесконечных многочленов и актуально бесконечных вещественных чисел остаётся такой неубедительной, потенциальная и актуальная бесконечность - это всего лишь два способа видеть один и тот же феномен с разных позиций референта, мы в буквальном смысле пытаемся утверждать, что результат эксперимента зависит от нашего описания. Данный факт иллюстрирует всю нелепость утверждения, что в аксиоматических теориях множеств отсутствует разделение на потенциальную и актуальную бесконечность. В действительности, это разделение играет экзистенциальную роль в современной интерпретации данных теорий. Таким образом, первое доказательство Кантора работает в рамках кардинальной арифметики (и распространяется в том числе на множество вещественных чисел), тогда как второе доказательство требует порядковой арифметики, где унифицирована длина стандартных последовательностей.

Впервые понятие мощности появляется в работе Кантора ''К учению о многообразиях'', где оно определяется как возможность однозначного и полного поэлементного сопоставления двух совокупностей \cite{cantor-power}. На этом этапе мощность служит для Кантора синонимом эквивалентности (изоморфизма) без разделения на неупорядоченные и упорядоченные структуры. Это позволило ему изящно сформулировать главный результат работы - равномощность многомерного пространства одномерному. Кантор продолжает развитие теории множеств в серии работ ''О бесконечных линейных точечных многообразиях'', где впервые вводится понятие выделенной счётной мощности - наименьшей бесконечной мощности, присущей множеству натуральных чисел \cite{cantor-countability}. Тем не менее, мощность ещё не формализуется как числоподобный объект: она остаётся абстрактным обозначением возможности установить взаимно однозначное соответствие. В частности, счётность означает возможность соответствия с последовательностью. Таким образом, в этих работах мощность ещё не обладает заземлением во внешнем унифицированном пространстве. Она представляет собой свойство самого математического построения - допускает ли оно ''касание'' по совпадению или нет. Мощность выступает лишь маркером сравнения: в рамках одного построения мощности могут совпадать, тогда как в других построениях такое совпадение приводит к противоречию. Без внешнего заземления счётность в одном контексте и счётность в другом - это принципиально разные понятия. Кульминацией работы Кантора становится идея возникновения бесконечных паттернов из структуры вполне упорядочения \cite{cantor-ordinality}. Для построения трансфинитных порядковых чисел он вводит три фундаментальных логических принципа:

\begin{enumerate}
    \item {\it Первый принцип порождения} - принцип присоединения единицы к некоему имеющемуся уже образованному числу $i + 1$.
    \item {\it Второй принцип порождения} - если задана определённая последовательность введённых целых реальных чисел, среди которых нет наибольшего, то создаётся новое число, которое мыслится как предел этих чисел, то есть как первое число, большее всех них.
    \item {\it Третий принцип стеснения (замыкания)} - всякая последовательность множеств, мощность каждого из которых не превосходит фиксированную мощность, и сама последовательность имеет такую же мощность, порождает множество, мощность которого совпадает с мощностью последовательности.
\end{enumerate}

Эти принципы требуют высокого уровня рефлексии по отношению к собственному мышлению. Особенно второй принцип, который вплотную подходит к способности сознания фиксировать абстрактную цель направленного процесса. Логически строго обосновать этот принцип как во времена Кантора, так и в рамках современных теорий, представляется весьма непростой задачей. При рассмотрении через призму дуализма внутренней и внешней перспектив принцип оказывается уязвимым с точки зрения воспроизводимости. Действительно, если один исследователь получает $\omega_1$ как предел натурального ряда, а другой - $\omega_2$, то без дополнительного критерия у нас отсутствует возможность сопоставить эти объекты. Необходим внешний критерий унификации, позволяющий вписать оба результата в единое пространство. Здесь и вступает в действие третий принцип (стеснения), фиксирующий области числовых классов. Какими бы ни были $\omega_1$ и $\omega_2$, если они имеют одну и ту же мощность, то они определяют один и тот же структурный объект с точностью до изоморфизма. Переход к следующему числовому классу обеспечивается всей совокупностью производных элементов. Важно подчеркнуть, что на этом этапе мощность ещё играет роль специфического для данного построения способа сравнения - детектора коллизий при попытке вписать конструкцию во внешний контекст. Ещё не введено взаимно однозначное соответствие с сохранением порядка. Таким образом, требование сравнимости по мощности выступает критерием внешней унификации, формируя конструкцию, которая образует иерархию мощностей. При этом данная иерархия никак не связана с различием мощностей вещественных и натуральных чисел из его прошлого доказательства. Также следует отметить, что поскольку с самого начала Кантор мыслил порядковые числа как классы эквивалентности определённые с точностью до изоморфизма, его система естественным образом обосновывает почему значение разности $\omega - \omega$ не определено. Он не видит возможности полностью унифицировать множество $\omega$, поскольку в таком случае потерял бы свойство самоподобия отделяющее конечное от бесконечного, ему приходится совмещать внутреннее представление с внешним. В современных теориях $\omega$ считается строго определённым множеством, для которого сохранения неопределённости $\omega - \omega$ никак не аргументируется.

Кроме того, Кантор формулирует в этой работе собственный взгляд на континуум. Он предлагает рассматривать его не просто как замыкание, а как многократное (в том числе трансфинитное) замыкание точечных множеств. Кантор не соглашается с Дедекиндом в том, что для анализа достаточно представлять вещественные числа как сечения множества рациональных \cite{cantor-objection-to-dedekind}. Для него сечения - это всего лишь пределы последовательностей первого порядка. Существуют также пределы последовательностей второго, третьего и вообще любых порядков, в том числе трансфинитных. Пусть они и не дают новых объектов в аддитивной структуре вещественных чисел, но позволяют более просто описывать многообразие аналитических отношений. Такая конструкция подобна конструкции второго числового класса: в первом порядке мы получаем бесконечные числа, использование которых в следующих порядках порождает новые трансфинитные пределы, например, $\lim \omega^i = \omega^\omega$. Поэтому мощность аналитического континуума Кантора оказывается связанной с иерархией мощностей, возникающей в конструкции порядковых чисел, что и приводит Кантора к утверждениям, впоследствии названным континуум-гипотезой. Как уже отмечалось, проблема состоит в том, что аддитивная группа воспроизводит себя не только в мультипликативной (где $x^0$ становится нейтральным элементом), но и последовательно на всех уровнях степенной башни и далее сквозь все счётные паттерны порядковых чисел. Классический предел замыкает лишь аддитивную группу, тогда как даже мультипликативная группа остаётся незамкнутой. Например, $a = \lim 1/n$, $b = a + a = \lim (1/n + 1/n)$, но $a = 2b$, хотя для всех конечных $n$ такого равенства нет. Это приводит к необходимости введения идеалов. Кантор же предлагает изящное решение, позволяющее замкнуть все уровни воспроизводства аддитивной группы в едином аналитическом фреймворке. Так или иначе, с этой работой понятие мощности континуума приобретает двойственность: как для континуума Дедекинда, так и для аналитического континуума Кантора, существует математическое построение, демонстрирующее невозможность внешней биекции. При этом континуум Кантора гораздо шире и включает в себя континуум Дедекинда. Уже в этой работе Кантор чётко разделяет общее понятие мощности как маркера сравнения и понятие структурного количества (которым он называл порядковые числа), противопоставляя свой непротиворечивый подход идеям Больцано, где эти концепции находятся в зачаточном состоянии \cite{cantor-bolzano}. Эту мысль он последовательно излагает также в критике работы Фреге \cite{cantor-frege}. Стоит также отметить, что Кантор, вопреки расхожему мнению, не критикует саму концепцию бесконечно малых, он лишь указывает на проблему совместимости линейных бесконечно малых величин с трансфинитными числами. Он явно позиционирует свои новые числа как к ключ к ответу на вопрос о возможности определения бесконечно малых величин. Но его вывод вполне однозначен: "по моему убеждению, область определимых величин не исчерпывается конечными величинами и границы нашего познания можно соответственно расширить..." \cite{cantor-infinitesimal}. Суть конфликта линейных бесконечно малых величин он раскрывает в своей философской работе "К учению о трансфинитном" \cite{cantor-infinitesimal-problem}. Дело в том, что Кантор видит порочный круг в интуитивном разграничении конечного и бесконечного, то что мы называем конечным имеет смысл только в противопоставлении предполагаемому бесконечному, а потому он видит позиционирование бесконечного через выделенное свойство подобия существенно более строгим. Но если существует бесконечно малая линейная величина, то она обладает аддитивностью, а значит может быть взята трансфинитное число раз. Однако, если бы $\zeta \times \omega$ давало конечную линейную величину $x$ то, взятие ещё $n$ таких величин приводило бы к тому же результату: $\zeta \times (n + \omega) = \zeta \times n + \zeta \times \omega = \zeta \times n + x = x$, что нарушало бы линейность, оставляя лишь вариант, что такое произведение остаётся бесконечно малым. Но поскольку это справедливо для любого трансфинитного порядкового числа, то такая величина лишена смысла, ведь мы никоим образом не можем получить из неё результат доступный нашему восприятию. Фактически Кантор показывает, что линейная величина не совместима с алгеброй порядковых чисел, полученной строгим определением бесконечности через свойство, а значит она не может существовать поскольку Кантору не было известно иных способов введения бесконечности минующих порочный круг. Это действительно проблема, если мы постулируем аксиому Евклида о превосходстве целого над собственной частью, то у нас отсутствует какое-либо свойство позволяющее разделить в теории конечное и бесконечное, кроме интуитивной исходной маркировки. С современных позиций мы можем решить этот вопрос относительно легко, структурная морфология индуктивно-дедуктивных наборов обладает свойством подобия, но при этом с внешней перспективы наборы являются Евклидовыми. Таким образом, бесконечное внешне неразличимо с конечным, но внутренне обладает выделенным свойством.

Взгляды Кантора продолжают эволюционировать, он поликует работу ''Об одном элементарном вопросе учения о многообразиях'', в которой формулирует известный диагональный метод \cite{cantor-diagonal}. Несмотря на то, что данная статья филигранно фиксирует структурные требования критичные для корректной интерпретации того, что означает детектируемое различие мощностей, в каком именно смысле следует понимать невозможность взаимно однозначного соответствия, от которого явно требуется сохранение структурного паттерна. Метод предложенный Кантором очень быстро некритически обобщается, на случаи отсутствия внешнего контекста, обеспечивающего возможность структурного сравнения. Кантор пишет об ''элементах, зависящих от бесконечно многих координат'', он сознательно фиксирует свойство универсальности позиционного пространства, при этом координаты представимы в виде последовательности, что обеспечивает свойство тотальной упорядоченности структуры. Классического определения последовательности ему явно не достаточно, поскольку оно подразумевает потенциальную бесконечность не обладающую достаточной степенью унифицированности, для однозначного покоординатного сопоставления всех элементов $E = (x_1, x_2, \dots, x_{\nu}, \dots$ многообразия $M$. Множество координат должно быть единым для всех элементов, требуется не произвольная структурная топология множества, а топология имеющая внешнее ограничение. Кантор также явно обозначает связь между диагональным методом и доказательством непредставимости вещественных чисел в виде последовательности. Он пишет: ''... содержится первое доказательство теоремы, что существуют бесконечные многообразия, которые нельзя взаимно однозначно отобразить на совокупность всех конечных целых чисел... например, совокупность всех действительных чисел произвольного интервала $(\alpha, \beta)$ нельзя представить в форме последовательности... Можно, однако, получить значительно более простое доказательство этой теоремы, не зависящее от рассмотрения иррациональных чисел.'' Иными словами, он говорит о теореме о существовании неперечеслимых множеств, где действительные числа выступают лишь примером, фактически, что новое и предыдущее рассуждения демонстрируют расслоение классов эквивалентности по мощности на различные типы при определённых предположениях о структуре. Он отнюдь не утверждает, что доказывает теорему о вещественных числах или что вещественные числа допускают описание посредством универсальной тотально упорядоченной системы координат, наоборот он склонен описывать их кардинальными (потенциально бесконечными последовательностями), чему мы находим подтверждение в его следующей работе (будет рассмотрено далее). Математический анализ понимает под вещественными числами объекты с описываемые одним набором допущений, в то время как в данной работе, объекты рассматриваются в разрезе иных допущений, которые возможно являются минимальными. Сведение вещественных чисел к точкам универсального тотально упорядоченного пространства представляет собой сильное упрощение позднейших аксиоматических подходов, а не собственную позицию Кантора. Именно такое упрощение приводит к расхождению между континуумом Кантора и континуумом в современном понимании, которое в свою очередь тесно переплетается с различием в интерпретации кардинальности Кантором и тем её аналогом в аксиоматических теориях. Таким образом, Кантор воспринимает континуум существенно шире, множество точек в универсальном упорядоченном пространстве координат в диагональном доказательстве выступает лишь как подсистема континуума. Данное хитросплетение существенно затрудняет анализ и усиливает идеологическую и догматическую составляющую современных теорий. Кантор обобщает диагональное доказательства на линейный континуум, где аналогичным образом, требует строгого порядка заданного на этом множестве. Если мы фиксируем отрезок $[0, 1]$ некоторого заданного пространства, то мощность функций $f : [0, 1] \to \{ 0, 1 \}$ в этом пространстве превосходит мощность отрезка $[0, 1]$. Без фиксации единого координатного пространства, или пространства определения линейного континуума, доказательство теряет силу. В частности, мы не можем просто так распространить доказательство на современный контекст. Если мы понимаем под тождеством вещественных чисел простую невычислимость различия \cite{compare-undecidable}, то тут же теряем возможность говорить о несчётности таких чисел, любая бесконечность невычислима и любое множество можно вполне упорядочить, отсутствие ограничения на бесконечность делает доказательство некорректным. Но Кантору важно другое, он обращает внимание на то, что мощность как абстрактный маркер структурных коллизий, допускает обобщение математического построения, где воспроизведение коллизии формирует иерархию. А значит можно попытаться построить единое унифицированное пространство позволяющее сравнивать бесконечных конструкции. 

Именно эту задачу Кантор решает в фундаментальном труде ''К обоснованию учения о трансфинитных множествах'' \cite{cantor-transfinite}. Он отходит от понимания мощности как локального маркера невозможности сопоставления и переопределяет её глобально через независимую алгебру кардинальных чисел, основанную исключительно на взаимно однозначном соответствии, без учёта какой-либо структурной (морфологической) информации о множествах. Кантор фиксирует соответствующую терминологию: числа называются кардинальными и соответствуют мощности множества, а взаимно однозначное соответствие без учёта морфологии - эквивалентностью. Мощность теперь применима к любой структуре. Какой бы морфологией ни обладало множество, его внутреннее строение может порождать структурные коллизии при попытке вписать конструкцию в единое пространство, топологически поддерживающее данную морфологию. Однако такие коллизии носят внутренний характер и относятся к ограничениям конкретной алгебры, связанной с топологией пространства, а не к мощности как таковой. Мощность же объединяет фундаментальную невозможность сопоставления, которая воспроизводится в любых возможных структурах. Счётность теперь определяется глобально и допускает внутренние структурные коллизии, обусловленные конкретными морфологиями. Иными словами, счётные множества в рамках заданных топологических пространств могут не допускать взаимно однозначного соответствия, что не влияет на их счётность в кардинальном смысле. Кантор утверждает, что мощность линейного континуума может быть выражена как $2^{\aleph_0}$, но использует это исключительно для воспроизведения результата предыдущей работы: многомерный континуум и континуум счётной размерности имеют ту же мощность, что и одномерный \cite{cantor-multicontinuum}. Таким образом, он интегрирует этот результат в систему переопределённых понятий. Именно здесь он говорит о вещественных числах как последовательностях внутри кардинальной алгебры, разделяя упорядоченные объекты рассматриваемые в рамках диагонального аргумента и вещественные числа. Иными словами, он явно связывает вещественное число с произвольной счётной последовательностью, которая может быть представимы любым порядковым типом, а не только $\omega$. Однако он отказывается от аналогичной интеграции утверждения о несчётности континуума ($\aleph_0 < 2^{\aleph_0}$). Кантор говорит, что кардинальные числа образуют иерархию, однозначно определяемую через порядковые числа, но нигде прямо не утверждает, что множество всех подмножеств также позволяет разграничивать уровни этой иерархии \cite{cantor-cardinal-hierarchy}. Он рассматривает порядковый тип линейного континуума и показывает, что для его характеристики достаточно требовать совершенства (включения всех пределов) и наличия вложенного определённым образом счётного множества (аналогичного множеству рациональных чисел), при этом никаких требований к мощности не возникает \cite{cantor-continuum-order}. Таким образом, в своём зрелом понимании Кантор полностью отходит от мощности, привязанной к конкретному построению, превратив это понятие в универсальное обобщение не замкнутой в каком-либо структурированном универсуме.

Кантор исследует различные морфологии порядковых типов, среди которых особо выделяет порядковые типы вполне упорядоченных множеств. Вполне упорядоченные типы формируют собственную алгебру, где соответствие, сохраняющее порядок, называется подобием, а абстракции свойств конкатенации - порядковыми числами. Объектами, с которыми работает Кантор, являются классы эквивалентности: порядковые типы и типы мощности, которые аннотируют реальные множества, но сами множествами не являются. Нигде в его работах мы не встречаем аксиомы экстенсиональности в современном смысле: совпадение элементов двух множеств влечёт лишь их эквивалентность или подобие, но не обязательно равенство самих множеств. Кантор называет кардинальные числа (как свойства типов мощности) и порядковые числа (как свойства порядковых типов) ''числоподобными'' объектами, утверждая, что вместе они охватывают всё, что поддаётся строгому математическому определению \cite{cantor-numberlike}. Здесь он не вполне точен, поскольку возможно определить классы эквивалентности бесконечных множеств, арифметика которых совпадает с арифметикой конечных чисел (индексные числа). К ним относятся, например, орды, гросс-объекты и целые омнифические положительные числа. Кантор прав в том, что такие объекты можно описывать как вполне упорядоченные, однако их полное описание требует дополнительных морфологических свойств, обеспечивающих возможность единичной декомпозиции, и, как следствие, - однозначной метрики $x - y$, которая в общем случае отсутствует у вполне упорядоченных множеств. Более того, Кантор настолько последователен в разделении алгебр, что считает множества, различающиеся порядком элементов, различными множествами \cite{cantor-order-difference}. Это математически корректно: тождество означает отсутствие любых различающих свойств, в то время как упорядоченные множества различаются порядком элементов, порядковыми типами и порядковыми числами, даже если содержат одни и те же элементы. Кантор был воспитанником классической системы образования, где огромное внимание уделялось философии, онтологическому и эпистомологическом обоснованию математических построений. Именно поэтому, он хорошо знал софизмы сформулированные Аристотелем, среди которых присутствует формальная ошибка ложной дихотомии - необоснованное предположения двузначности пространства альтернатив в некоторой области, позволяющую использовать для описания данной области классическую логику противоположностей. Кантор разделяет кардинальную и порядковую алгебры, чтобы иметь возможность пользоваться логикой внутри обоих. Он неоднократно указывает, что смешение различных представлений о бесконечном числе недопустимо, поскольку приводит к невозможности однозначного логического описания. 

Кантор подчёркивает универсальность кардинальной абстракции и её применимость к любым типам морфологии, при этом всякое равенство между порядковыми типами, получающееся из двух элементарных операций, остаётся справедливым после замены всех типов на соответствующие кардинальные числа \cite{cantor-cardinal-universality}. Определив функцию степени через элементарные операции (включая бесконечное произведение и бесконечный ряд, которые сами по себе определимы благодаря вполне упорядочению), мы естественным образом должны были бы получить $\omega = 2^\omega \implies \aleph_0 = 2^{\aleph_0}$ как результат общего правила сформулированного Кантором. При этом порядковое отношение $\omega < 2^\omega$ прекрасно согласуется с кардинальным $\aleph_0 = 2^{\aleph_0}$. Но в современных интерпретациях мы идём на открытую конфронтацию с Кантором допуская одновременную истинность как $\omega = 2^\omega$, так и $\aleph_0 < 2^{\aleph_0}$. Кантор описывает пределы через бесконечные ряды: $\lim\limits_{\nu}^{outer} \alpha_{\nu} = \alpha_1 + (\alpha_2 - \alpha_1) + (\alpha_3 - \alpha_2) + \dots$ \cite{cantor-limit}, что позволяет ему явно зафиксировать внешнюю составляющую порядковых чисел. Именно следствие $\omega < 2^\omega$ из теоремы Кантора гарантирует однозначность определения порядковых типов внутри экспоненциальных областей с точностью до изоморфизма. Конструкция Кантора получает многоуровневую поддержку: через изоморфизм внутри экспоненциальных областей, через изоморфизм самих этих областей внутри мощности и через иерархию мощностей. Для внешней унификации достаточно предположить, что множество конечных порядковых чисел всегда попадает в фиксированную экспоненциальную область. Это обеспечивает как свойство самоподобия, так и многоуровневую унификацию бесконечного пространства. Иными словами, Кантор аккуратно выстраивает измерительный аппарат для всего бесконечного - универсальную линейку, обеспечивающую сравнимость любых математических конструкций. Интересно, что экспоненциальные области воспроизводства порядковых структур при таком построении находятся в точном соответствии с уровнями воспроизводства аддитивной группы в арифметике натуральных чисел. Выстроенная Кантором система настолько связна и элегантна, что остаётся совершенно неясным, каким образом математическое сообщество могло проигнорировать эту глубину и постулировать обобщение внутренней кардинальной структуры на порядковые типы в виде отождествления $\aleph_0 = \omega$.

Несмотря на то, что множество деталей указывает на то, что Кантор как минимум сомневался в применимости теоремы, носящей его имя, к кардинальной алгебре, а попытка напрямую связать неравенство $2^{\aleph_0} > \aleph_0$ с этой теоремой может объясняться недоучётом эволюции его понятийного аппарата, настоящее наблюдение остаётся интерпретацией. Необходимо отметить, что единственное место, где Кантор явно использует данное неравенство, - это его письмо Дедекинду из санатория Халленклее от 31 августа 1899 года. Однако в этот период Кантор переживал серьёзные профессиональные и личные трудности, которые в конечном итоге привели к тяжёлому нервному срыву \cite{dauben-cantor-1899, grattan-guinness-1899}. Данный факт в сочетании с довольно поверхностным (для Кантора) содержанием письма не позволяет рассматривать это высказывание как выражение его зрелой позиции. Вероятно, систему координат (вместо последовательностей) в диагональном доказательстве, обеспечивающую его математическую корректность, а также отсутствие теоремы Кантора в кардинальной алгебре можно объяснить редакторскими правками. Однако в таком случае трудно объяснить непримиримую позицию самого Кантора по отношению к возможности слияния кардинальной и порядковой алгебр. Более того, морфология вполне упорядоченных множеств естественным образом разделяет внутренние (левые) и внешние (правые) операции. Благодаря этому в чистой порядковой алгебре можно определить только внешние пределы и внешнее возведение в степень, поскольку данные операции связаны с вполне упорядоченными последовательностями и рядами. Внутренние операции в духе современных аксиоматик требуют выхода за пределы вполне упорядоченной морфологии. Кантор свободно оперирует как левыми, так и правыми бесконечными последовательностями при рассмотрении различных упорядоченных структур. Поэтому предположение, что он неосознанно оставляет степень исключительно внешней операцией, выглядит натянутым. Внутренние операции при этом остаются возможными, но они не входят в состав порядковой алгебры, в рамках которой действует теорема Кантора. Как бы то ни было, математическое наследие Кантора фиксирует тщательно структурированную онтологическую картину, в которой числа возникают как эмерджентные свойства структур, а не как универсальные инварианты мышления. Кантор не отрицает своих предшественников, а продолжает их исследование, облекая концептуальные прорывы в строгую математическую форму. Аккуратность, с которой он это делает, позволяет ему избежать догматических коллапсов и оставляет его теорию и программу совместимыми со всеми концепциями, ставшими альтернативой мейнстриму, - от новых оснований Куайна до гросс-объектов.

С современных позиций легко увидеть напряжение между внутренней и внешней перспективами, которое вскрыл Кантор уже в своей первой работе. В одном контексте экспоненциальное множество можно сопоставить с линейным, в другом - нельзя. Сама по себе структура не гарантирует возможности сравнения множеств: как алгебраические, так и вещественные числа можно вполне упорядочить. Однако структура позволяет вообразить пространство с общей дискретной топологией, в которое мы можем вписать разные объекты для их сравнения. Чтобы соотнести бесконечные объекты между собой нам необходима воспроизводимость результата бесконечного построения, формирующего законченную целостность (универсальность в терминах теории категорий). Отсутствие воспроизводимости означает, что в одном исследовании ряд натуральных чисел может быть продолжен до $\dot{\omega}$, а в другом - до $\dot{\omega}^{\dot{\omega}}$ или даже до $\dot{\omega} \uparrow\uparrow \dot{\omega}$. При этом индуктивная морфология этих рядов будет идентичной, позволяя говорить об их подобии через взаимно однозначное соответствие. Такой подход не различает экспоненциальных уровней и сливает бесконечность в единую сущность, обеспечивающую индуктивную замкнутость универсума. Кантор связывал это с отсутствием структуры и кардинальной алгеброй, но разделение внутренней и внешней перспектив позволяет легко увидеть, что дело в отсутствии механизма внешней унификации. Именно это мы фиксируем в новых основаниях Куайна: замкнутость без возможности говорить о каком-либо унифицированном ограничении, которое мы связываем с аксиомой выбора. Если структура задана, то существенно легче вообразить общее пространство с дискретной топологией, которое и обеспечивает необходимую унификацию. В случае вещественных чисел коллизия не говорит о том, что их нельзя расположить в виде последовательности. Она указывает, что представление вещественных чисел унифицированной по длине позиционной записью приводит к тому, что упорядоченная совокупность всех таких чисел не может вместиться в рамки этого эталона длины. Именно это Кантор фиксирует в порядковых типах: операция возведения в степень определяется чисто внешне, обеспечивая унификацию экспоненциальных областей, в то время как внутренняя составляющая, действующая на полиномиальном уровне, гарантирует свойство подобия объектов своим частям, необходимое для того, чтобы считать эти объекты бесконечными. Благодаря морфологии вполне упорядочения, которая не допускает введения внутренних пределов и внутреннего возведения в степень, Кантору удаётся объяснить внешнюю унификацию чисто структурно и сохранить свойство ассоциируемое с бесконечностью. Именно этот механизм гарантирует ограничение воспроизводимости бесконечных построений через теорему Кантора, обеспечивающую $\omega < 2^\omega$. В свою очередь, это гарантирует определение $\omega$ внутри экспоненциальной области, и таким образом вся порядковая арифметика становится строго определённой с точностью до изоморфизма. В одном исследовании мы можем получить натуральный ряд как $\dot{\omega}$, а в другом - как $\dot{\omega}^3$, но для нас это несущественно, поскольку внутренняя составляющая обеспечивает полиномиальную замкнутость и полную воспроизводимость структуры внутри зафиксированных внешне экспоненциальных областей, которые соответствуют уровням воспроизводимости аддитивной структуры в степенных башнях. Ничего подобного мы не наблюдаем в современных аксиоматиках, где внутренняя и внешняя перспективы хаотично смешаны. Вера в унификацию посредством аксиомы выбора и общего функционального аппарата при одновременной рефлексивности бесконечных множеств повисает в воздухе, не находя структурной поддержки. Это приводит к отрицанию теоремы Кантора для ординалов и, как следствие, к устойчивой условной непротиворечивости теории. Аксиомы требуют ритуалов правильной интерпретации, а указание на данный феномен воспринимается как экзистенциальная атака. 

В 1906 году Герхард Хессенберг публикует фундаментальный труд ''Основные понятия теории множеств'', в котором делает ряд смелых обобщений. Он формализует вербальное описание взаимно однозначного соответствия, которым, по всей видимости, умышленно пользовался Кантор, и превращает его в полностью функциональное описание. Однако функция $f_c(x)$ в кардинальной арифметике и функция $f_o(x)$ в порядковой арифметике представляют собой принципиально разные объекты. В современной терминологии $f_c(x)$ можно описать как класс пар, который не обязательно является множеством в упорядоченном универсуме (например, множество в новых основаниях Куайна). Напротив, $f_o(x)$ заведомо обладает порядком, индуцированным областью определения и областью значений, и представляет собой объект с фиксированным рангом в упорядоченном универсуме. Кардинальная функция - это объект чисто внутреннего индуктивного описания; она не подразумевает внешней унификации. По сути, именно она лучше всего описывает индуктивные функции в арифметике Пеано, где законченная целостность всех арифметических объектов не определена. Здесь путь открыт как для минимальных расширений до кардинальной алгебры, так и до расширений с унификацией - порядковой алгебры. Ключевая разница состоит в том, что в кардинальной алгебре бесконечные объекты принципиально не сравнимы между собой: они все подобны друг другу и обладают свойством подобия своим частям, но не предполагают возможности сопоставления в унифицированном пространстве. Поэтому в кардинальной алгебре не работает диагональный аргумент: морфологическое подобие индуктивных структур не гарантирует, что для каждого элемента сингулярности одного множества сингулярность другого содержит соответствующий элемент. Порядковая функция, в отличие от кардинальной, допускает внешний унифицированный контекст. Все вполне упорядоченные морфологии (многочлены) в рамках одной экспоненциальной области могут быть представлены через одну. Иными словами, всякое дискретное пространство конечной размерности может быть представлено одномерным пространством. Именно выразимость внутри области через общую морфологию обеспечивает унификацию понятия последовательность, которая, в свою очередь, гарантирует корректность диагонального аргумента. Однако Хессенберг обобщает диагональный аргумент в функциональной нотации на произвольные множества, смешивая кардинальную функцию $f_c(x)$ и порядковую функцию $f_o(x)$ в одну \cite{hessenberg}. Само по себе это обстоятельство не является фатальной ошибкой: мы просто отказываемся от чисто внутреннего представления и постулируем, что все объекты теории допускают внешнее сравнение. Фактически мы переходим в полностью упорядоченный универсум, где доминирует порядковая арифметика.

Далее Хессенберг утверждает счётность множества $\omega^{\omega} = \omega + \omega^2 + \dots + \omega^n + \dots$, опираясь на то, что формула для всякого элемента $\omega^n$ конечна (состоит из $n$ чисел), а это, по теореме о конечной записи, влечёт счётность $\omega^\omega$  \cite{hessenberg-countable-exponent}. Иными словами, он ссылается на теорему из первой работы Кантора о счётности алгебраических чисел, доказанную через счётность многочленов c конечной записью. Эта теорема по сути интерпретируется как утверждение о счётности любого счётного объединения счётных множеств. Однако в работе ''К обоснованию учения о трансфинитных множествах'' Кантор относит такую теорему строго к кардинальной арифметике ($\aleph_0 \cdot \aleph_0 = \aleph_0$), тогда как порядковая арифметика работает несколько иначе: $\omega \cdot \omega = \omega^2$. Таким образом, Хессенберг не замечает эволюцию взглядов самого Кантора и смешивает кардинальную алгебру с порядковой. Впрочем, в разделе о парадоксах он честно описывает уязвимость концепции конечного представления, приходя к выводу, что само понятие конечной выразимости не может быть разрешимым свойством объектов \cite{hessenberg-paradoxes}. С теоремой о конечной записи оказываются связаны многочисленные парадоксы, такие как парадокс Ришара (Число называется Ришаровым, если оно не удовлетворяет тому математическому свойству, которое стоит под его же порядковым номером в алфавитном списке всех определений свойств чисел. Но поскольку само определение ''быть Ришаровым числом'' тоже находится в этом списке под некоторым номером $n$, возникает вопрос: является ли число $n$ ришаровым \cite{richard-paradox}) и парадокс Берри (наименьшее натуральное число, которое нельзя задать фразой, содержащей менее ста букв \cite{berry-paradox}). Эти парадоксы в конечном итоге привели к существенному ограничению математического языка - переходу к логике первого порядка, сильно ограничивающей возможность теории описывать собственный синтаксис. То есть, Хессенберг хоть и воспроизводит результат Кантора в доказательстве теоремы о конечной записи, явно указывает на потенциальные проблемы подобной интерпретации в контексте возникающих парадоксов. Данный факт мог бы привлечь внимание к доказательству Хессенберга и привести к формулировке ключевого вопроса: если мы отказываемся от чисто внутреннего описания, какие теоремы кардинальной алгебры можно безопасно транслировать в упорядоченный универсум? 

Вместо этого математическое сообщество интегрировало доказательство счётности любого счётного объединения счётных множеств в аксиоматическую теорию множеств. Для этого используется два подхода восходящих к кардинальным рассуждениям Кантора: рассматривается множество функций с конечным носителем, в котором ограничивается сумма участвующих элементов (высота многочлена), рассматривается бесконечная счётная по столбцам и строкам матрица, где индекс элемента выражается через индекс строк и колонок. Однако в упорядоченном универсуме, где существует механизм внешней унификации оба подхода оказываются затруднительными и ведут к порочному кругу. В первом случае, количество комбинаций, где сумма элементов не превышает определённое число, растёт экспоненциально, и, таким образом, само формирует экспоненциальную последовательность порядковых чисел, счётность предела которой требуется доказать. Аналогичным образом в случае матрицы в контексте упорядоченного универсума нам необходимо гарантировать унификацию всех строк и представимость их одним порядковым типом. Но если супремум всех строк описывается порядковым типом следующего экспоненциального уровня, то есть не отображается на порядковый тип самого множества строк (объединение порядковых типов строк не изоморфно строкам, например, как в случае $\bigcup_{n < \omega} n = \omega$), то индекс столбцов уже должен быть определён на множестве следующего экспоненциального уровня равно как и его арифметическая комбинация с индексом строк. Сама по себе аксиоматика не позволяет однозначно интерпретировать поведение двойного предела (по строкам и столбцам). Теорема Кантора позволяет раскрыть эту неопределённость внешне, так как в нашей интерпретации она применима ко всем множествам теории без исключения. Тогда у нас имеется простые контрпримеры: $2 + 2^2 + \dots$ или $\omega + \omega^2 + \dots$, к утверждению счётности любого счётного объединения. В то же время, если мы можем гарантировать изоморфность супремума всех строк множеству строк, теорема остаётся корректной, сохраняя изоморфизм полиномиальной области. В какой-то момент для обоснования применимости кардинальных рассуждений в упорядоченном универсуме возникает идея промежуточных объектов (множеств функций с конечным носителем), восходящим к интуиции о потенциальной бесконечности твёрдо закрепившейся в математике благодаря средневековым схоластам. В ZFC множество функций с конечным носителем существует и определяется стандартными аксиомами. Однако в контексте упорядоченного универсума возникает более тонкая проблема: при переходе к пределу (объединению по всем конечным носителям) этот объект становится неотличимым от множества всех функций. Конечный носитель в пределе даёт бесконечный носитель. Таким образом, использование промежуточных объектов для обоснования счётности $\bigcup_{n < \omega} X_n$ (где каждое $X_n$ счётно) просто добавляет новое звено к порочному кругу: чтобы оправдать счётность предела, мы апеллируем к объектам, чьё поведение в пределе уже предполагает ту самую счётность, которую требуется доказать.

Хессенберг выводит из нормальной формы порядковых чисел третью независимую алгебру - арифметику натуральных операций \cite{hessenberg-sum}. Эта алгебра фиксирует чисто внешнюю составляющую порядковых типов и уже позволяет ввести понятие индексных чисел как абстракций нормальной формы. Коммутативность натурального сложения и умножения делает введение операций разности и деления чисто техническим вопросом - они возникают как естественные эмерджентные свойства полученной алгебраической структуры. Иными словами, Хессенберг доводит теорию Кантора до логического завершения: три алгебры объединяются в единую непротиворечивую картину. Кардинальная становится классической индуктивной арифметикой, замкнутой относительно всех комбинаторных операций; индексная - исключительно внешним описанием, обеспечивающим максимальную точность классификации бесконечных объектов; порядковая - гибридом, сохраняющим внутреннее представление в полиномиальной области, но внешне высвечивающим степенные уровни воспроизводства аддитивной структуры. Все три алгебры оказываются неразрывно связаны, поскольку описывают одни и те же объекты - множества. Связь между кардинальной и порядковой алгебрами демонстрирует Кантор, а Хессенберг находит путь к индексной алгебре. Таким образом, несмотря на то, что Хессенберг не замечает опасности унифицированной функциональной нотации и использует кардинальную логику в контексте, который становится внешним, он одновременно даёт инструмент, позволяющий обнаружить эту погрешность и исправить её. Формально функции $f_i(x)$ в индексной алгебре как объекты отличны от функций $f_o(x)$ в порядковой и $f_c(x)$ в кардинальной. Индексные числа обладают уникальной морфологией, которая позволяет очень точно описывать механику предельных переходов. Благодаря объектам типа $\dot{\omega} - 1$ внешнее представление предела $2^\omega$ через бесконечный ряд приобретает вид $2 + 2^2 + \dots + 2^{\dot{\omega} - 2} + 2^{\dot{\omega} - 1}$. Фактически все элементы $\dot{\omega}$ от $0$ до $\dot{\omega} - 1$ участвуют в комбинациях множества всех подмножеств. Объединение всех конечных носителей функций бесконечно ровно в том же смысле, в каком бесконечно объединение всех конечных ординалов. Ординал как максимальный носитель комбинации включает в себя все предшествующие. Отсюда явственно следует, что предположение о существовании промежуточных объектов (вроде множества функций с конечным носителем) требует дополнительных обоснований. Тем не менее, исключительно значимая находка Хессенберга остаётся невостребованной в аксиоматических теориях множеств до тех пор, пока Джон Конвей не возвращается к ней, чтобы вписать всю ординальную структуру в более общую морфологию древовидных конструкций.

Следующим значительным шагом становится опубликованная в 1908 году Эрнстом Цермело первая аксиоматика теории множеств, ставшая ответом на парадоксы Бурали-Форти и Рассела \cite{zermelo-axiomatization}. Парадокс Бурали-Форти \cite{burali-forti-theorem}, как было продемонстрировано в статье ''Несчётность множества континуум, шах и мат'', указывал на то, что порядковая арифметика Кантора не может быть замкнута. Это простое следствие второго принципа порождения, если всякая последовательность порядковых типов, не имеющая максимального элемента, определяет новый порядковый тип, непосредственно следующие за типами последовательности, то наивное обобщение - совокупность всех порядковых типов - не может быть порядковым типом (множеством) в этой системе. Парадокс Рассела \cite{russel-paradox}, в свою очередь, компрометирует именно кардинальную арифметику, где множество всех подмножеств оказывается изоморфным исходному множеству. Именно это свойство кардинальной арифметики позволяет считать счётными все формальные определения теории, что приводит к ряду синтаксических парадоксов, таких как парадокс Ришара или парадокс наименьшего неопределимого ординала. В порядковой арифметике теорема Кантора позволяет избежать таких противоречий, здесь тип множества всех конечных определений зависит от комбинаторной сложности синтаксиса, в частности, для языка с конечным количеством слов множество определений чисел не изоморфно множеству чисел, что и обуславливает избыточность определений. Даже собственно парадокс Рассела в порядковой алгебре не так страшен, ведь порядковый тип множества всех множеств не содержащих себя больше любого индекса внутренних множеств, что позволяет ограничить самоприменимость. Кардинальная алгебра в наивной предикативной форме оказывается крайне уязвимой для таких парадоксов, здесь требуется стратификация формул (как в новых основаниях Куайна) или существенное ослабление логики (например, наивная теория множеств с позитивной логикой). 

Цермело недавно доказавший, с использованием аксиомы выбора, теорему о том, что любое множество можно вполне упорядочить \cite{zermelo-well-order}, судя по всему, хорошо осознаёт возможность противостоять парадоксам в рамках структурированного универсума, для чего его подход позволяющий вполне упорядочить любое множество кажется идеальным. Он предлагает простую аксиоматику из семи аксиом: аксиома экстенсиональности, аксиомы элементарных множеств (пустого, синглтона и пары), схема аксиом выделения, аксиома степени, аксиома объединения, аксиома выбора и аксиома бесконечности. Цермело вводит понятие достоверной пропозиции, которое предполагает доказательство разрешимости данной пропозиции для всех объектов области. Такое определения требует постоянной проверки всех используемых пропозиций, вследствие чего подвергается критике и в конечном итоге разрешается в логике первого порядка. Проблему которую пытается решить Цермело была продемонстрирована в статье ''Несчётность множества континуум, шах и мат'', дело в том, что парадоксы самореференции вовсе не требуют больших множеств и легко воспроизводятся на конечных, например, множество всех множеств определённых в данном параграфе не содержащих себя в качестве элемента. Формально такую конструкцию можно записать, например, следующим образом $A = \{ \emptyset, \{ x | x \in A \wedge x \notin x \} \}$, пропозицию $x \notin x$ для выделения элементов множества $A$ Цермело предлагает считать недостоверной. Таким образом, исходная аксиоматика Цермело была существенно более мощной, допуская произвольные пропозиции, в том числе из логики произвольного порядка, чем версия закреплённая мейнстримом. Аксиома выбора в рамках данной системы играет ключевую роль, именно она обеспечивает однозначную сравнимость множеств между собой, благодаря возможности вполне упорядочить каждое из них. В определённой степени данное поведение также является интерпретацией, требующей дополнительного предположения воспроизводимости выбора, чтобы гарантировать однозначность вполне упорядочения, но в данном случае несложно видеть из использования аксиомы, что подразумевается именно воспроизводимый выбор. В этой системе аксиом можно также рассмотреть логические принципы порождения порядковых чисел Кантора: аксиома объединения становится первым принципом, а аксиома бесконечности - вторым. Третий принцип порождения, связывающие кардинальную алгебру с порядковой отсутствует. Иными словами, данная система представляет собой аксиоматизацию исключительно порядковой алгебры, что даёт прочный каркас для преимущественно внешнего непротиворечивого описания множеств. 

Цермело вводит понятие эквивалентности и доказывает ряд базовых теорем, оставляя вполне упорядоченность и подобие (эквивалентность сохраняющую её) для второй части, которая так и не была опубликована. Крайне примечательно, что Цермело доказывает теорему Кантора для всех объектов полученного домена, но при это не доказывает теорему о счётности счётного объединения счётных множеств, которая критична для доказательства счётности $\omega^{\omega}$. При это Цермело определяет кардинальность через эквивалентность в универсуме своей аксиоматической системы. Эта кардинальность уже существенно отличается от кардинальности Кантора, поскольку всякая функция (соответствие) в этом универсуме рассматривается с внешней перспективы. Таким образом, теория Цермело остаются полностью корректной избегая всех известных парадоксов связанных с бесконечностью. По сравнению с теорией Кантора, она теряет кардинальную часть, неудачно переопределяя кардинальность через внутренний домен. Можно также указать на проблемы аксиомы бесконечности, которые мы подробно обсуждали в статье ''Несчётность множества континуум, шах и мат''. Данная аксиома не чистая, так как требует особой морфологии бесконечного множества, что приводит к выводимости рефлексивности, которую следует постулировать отдельно. Тем не менее, поскольку на ранних этапах развития теории множеств свойство бесконечности и рефлексивности не различались, предложенная Цермело аксиоматика стала существенным шагом вперёд позволяющим определить ограниченную область объектов, в которой можно безопасно работать. Более того, аксиома выбора, как постулат обеспечивающий универсализацию, то есть возможность вписать любое множество в общее пространство с некоторой дискретной топологией, вообще говоря противоречит безграничной рефлексивности. Благодаря этому предложенная аксиоматика, где теорема Кантора безупречно вписывается, в частности, как экспоненциальный ограничитель рефлексивности, сама создаёт напряжение препятствующее закреплению интерпретации смешивающей позиции референта. Однако, инерция консервативных взглядов о потенциальной бесконечности конечных представлений оказывается существенно сильнее. 

Невероятно любопытным оказывается один из первых учебников по теории множеств ''Основы теории множеств'', написанный Феликсом Хаусдорфом и опубликованный в 1914 году. Хаусдорф остаётся в рамках наивной теории множеств и не адаптирует аксиоматику Цермело, тем не менее, он избегает скомпрометированных доказательств и старается максимально разграничить кардинальную и порядковую алгебру, выстраивая мостики между ними там, где их аргументированность безупречна. При этом теорема о счётном объединёнии счётного множества счётных множеств остаётся исключительно кардинальной теоремой $\aleph_0\aleph_0 = \aleph_0 + \aleph_0 + \dots = \aleph_0$ \cite{hausdorff-cardinal-arithmetic}. Хаусдорф сочетает внутреннее кардинальное представление с внешним представлением упорядоченного универсума. С помощью первого он доказывает, что множество алгебраических чисел счётно, поскольку счётно множество всех конечных комплексов коэффициентов многочленов $(r_1, r_2, \dots, r_n)$ и его мощность $\aleph_0 + \aleph_0^2 + \aleph_0^3 + \dots = \aleph_0$ (что корректно в кардинальной алгебре Кантора). В то время как, во втором случае речь идёт о бесконечных комплексах натуральных чисел $(a_1, a_2, a_3, \dots)$ эквивалентных множеству вещественных чисел, где теорема Кантора приводит к иерархии мощностей $\aleph_0, 2^{\aleph_0}, 2^{2^{\aleph_0}}$ \cite{hausdorff-cardinal-hierarchy}. Это, в частности, позволяет ему использовать ряды для кардинальных чисел невозможные в неупорядоченной кардинальной арифметике Кантора. Он выводит арифметику и иерархию вполне упорядоченных типов, независимо от кардинальной алгебры: $1, 2, \dots, \omega, \omega + 1, \dots, \omega^2, \dots, \omega^{\omega}, \dots, \omega^{\omega^{\omega}}, \dots$ \cite{hausdorff-ordinal-arithmetic}. Однако, Хаусдорф избегает вопроса о мощности $\omega^{\omega} = \omega + \omega^2 + \omega^3 + \dots$, он говорит о счётности $2^{\omega} = \omega$, а также $n^{\omega}$ для любого конечного $n$, поскольку относит их к счётным подмножествам континуума. Более того, он в разделе о связи порядковых типов и мощностей, судя по всему, вполне осознанно, ограничивает последовательность счётных порядковых чисел составляющих счётный класс $Z(\aleph_0)$: $\omega, \omega + 1, \omega + 2, \dots, \omega2, \omega3, \dots, \omega^2, \omega^3$, не включая в неё ни тип $\omega^{\omega}$, ни следующие за ним. И точно также, вводит первый несчётный ординал $\omega_1$, следующий за всеми конечными и счётными ординалами: $0, 1, 2, \dots, \omega, \omega + 1, \dots, \omega2, \dots, \omega^2, \dots$, связанный с первой несчётной мощностью $\aleph_1$ \cite{hausdorff-cardinal-ordinal}. Иными словами, Хаусдорфу удаётся построить две параллельные иерархии кардинальных и порядковых чисел в упорядоченном универсуме, которые не зависят от счётности $\omega^{\omega}$. Конструкция Хаусдорфа допускала как счётность порядкового числа $\omega^{\omega}$, что ведёт к современной устойчиво недостоверной интерпретации, так и его несчётность, означающую, что кардинальная (где сама кардинальность понимается в терминах упорядоченного универсума) формула $\aleph_0\aleph_0 = \aleph_0$ применима лишь в случае, счётности предела последовательности объединяемых счётных множеств (формально если предел несчётен, то данный случай не может быть выражен формулой $\aleph_0\aleph_0 = \aleph_0$). Если углубляться в детали, то Хаусдорф даже показывает, что $\omega^{\omega*}$ является типом комплексов натуральных чисел в лексикографическом порядке, который эквивалентен типу комплексов в натуральном порядке (то есть $\omega^{\omega}$), и оба этих произведения являются типом иррациональных чисел \cite{hausdorff-uncountable-product}. С современных позиций это корректное утверждение, поскольку порядковый тип $\omega*$ как и $\omega$ имеет внешнюю доминанту изолирующую экспоненциальные уровни, в то время как для счётности нужна внутренняя составляющая, которую можно описать как $\omega^{\aleph_0}$, именно такая потенциальность требуется для доказательства счётности алгебраических чисел. Таким образом, парадоксы связанные с теоремой о конечности представления были в полной мере осознаны, равно как и их влияние на результирующую теорию. 

В контексте теоремы о счётности счётного объединения счётных множеств необходимо также упомянуть работу Вацлава Серпински 1918 года, который показал, что данная теорема недоказуема в общем случае без аксиомы выбора \cite{serpinski-choice-axiom}. Серпински указал на данную невозможность вне контекста аксиоматики Цермело, а просто как на факт того, что данная аксиома играет центральную роль во многих доказательствах. То есть, если бы мы аксиоматизировали кардинальную алгебру Кантора, то нам неизбежно потребовалась бы аксиома выбора в той или иной форме, чтобы доказать данную теорему. В это же время, Абрахам Френкель публикует первое издание ''Введения в теорию множеств'', где вслед за Хессенбергом смешивает ранние работы Кантора, где понятие мощности ещё не оформилось в полноценную концепцию, с его зрелыми работами разделившими кардинальную и порядковую алгебры \cite{fraenkel-introduction-1919}. Первое издание написано ещё с позиций наивной теории множеств и не опирается на аксиоматический базис Цермело, здесь воспроизводится теорема кардинальной алгебры о счётность счётного объединения счётных множеств, но в то же время, порядковые типы упоминаются бегло без анализа их мощностей. Френкель честно отмечает недостоверность предложенных в книге рассуждений в силу возникших парадоксов и намечает возможный путь исправления ситуации через строгую аксиоматику Цермело. Любопытно, что на начальных этапах проблемы теории множеств вызванные парадоксами были осознаны и математики использовали два подхода к преодолению кризиса. Первый подход, используемый Хессенбергом и Френкелем, заключается в использовании сомнительных методов доказательств, с честным указанием на ограниченную достоверность вследствие описанных парадоксов. Фактически данный подход фиксирует ключевую задачу: вот то что мы хотим иметь в результате, а вот то что препятствует этому. Второй подход Хаусдорфа, который заключается в исключении недостоверных методов из теории и формулировке того, что остаётся без них. Что в свою очередь ставит вопрос о целесообразности результатов получаемых недостоверными методами. Не ясно по какой причине, точка зрения Хаусдорфа, который явно склонялся к невозможности нарушения теоремы Кантора в упорядоченном универсуме, несмотря на большую достоверность оказалась невостребованной. Стоит отметить также, что в этот период, проблема смешения позиций референта была интуитивно прочувствована и хорошо осознана. В частности, Бертран Рассел сформулировавший знаменитый парадокс предлагает ввести иерархию типов для решения парадоксов \cite{russel-type-theory}. В сущности при таком подходе внутренняя перспектива замкнута в рамках конкретного типа, в то время как внешняя всегда реализуется через типы более высокого уровня. Рассел видит проблему в форме цикличности - порочного круга, когда целостность содержит элементы определённые через эту целостность. Таким образом, он практически вплотную подходит к заключению о том, что невозможно одновременно воспринимать одну и ту же сущность изнутри и снаружи. Однако, концентрируется на радикальном запрете смешения перспектив, вместо того, чтобы изолировать только случаи приводящие к противоречию.

Следует отметить также, что новая аксиоматическая система Цермела подвергалась и убедительной критике, как например, в опередившей своё время работе Туральфа Сколема ''Некоторые ремарки об аксиоматизированной теории множеств'' \cite{skolem-axiomatization}, которая осталась сильно недооцененной. Ключевой тезис Сколема заключается в том, что аксиоматика Цермело - это теория алгебраических структур, определённых набором правил, а не множеств как реальных объектов. Сколем сразу предлагает избежать ошибки рефицикации и считать, что мы оперируем лишь аннотациями множеств, а не самими множествами. Вся теория множеств, представляет собой не более чем описание простого домена - совокупности пар определяющих отношение принадлежности, а значит она никак не может иметь привилегированного положения, например, по сравнению с сюрреальными числами или гросс-объектами. Сколем доказывает теорему, что для удовлетворения счётного непротиворечивого набора аксиом достаточно счётного домена. Таким образом, само понятие бесконечности оказывается относительным, делая аксиоматику полностью внутренней (что по-видимому является свойством самой логики первого порядка). Если Кантор давал нам инструменты смены позиции референта - внутренняя кардинальная алгебра и гибридная порядковая алгебра, а Хессенберг расширил арсенал индексной алгеброй, то для смены позиции референта в аксиоматической теории множеств нам требуется метатеория. Это вызывает существенные проблемы с воспроизводимостью, как показывает Сколем сам домен теории не может быть определён однозначно (что даже не исправляется аксиомой регулярности). Например, пусть задана совокупность пар определяющая универсум нашей теории, во внешней теории воссоздадим этот универсум кумулятивной иерархией индексируя каждый шаг целыми омнифическими числами или гросс-числами. Если у нас получилось $\dot{\omega}$ шагов, мы можем построить подобные универсумы $2\dot{\omega}$ (делая 2 шага вместо одного) и $\dot{\omega}/2$ (используя только чётные шаги) или даже $2^{\dot{\omega}}$ или $\log_2 \dot{\omega}$, и каждый из них будет полноценным доменом теории. Это приводит к тому, что аксиома выбора, которая могла бы обеспечить переход из кардинальной алгебры (например, теории множеств с позитивной логикой) в алгебру сюрреальных чисел, оказывается ослабленной поскольку не может полноценно усилить контекст воспроизводимости, так как определена в рамках теории, которая сама имеет относительную воспроивзодимость. Фактически идеализируя одну аннотацию, мы блокируем естественный механизм переключения перспективы, и именно об этом говорит Сколем: при формулировке аксиоматики мы должны понимать, что кристаллизуем одну из алгебраических структур описания. Кроме того, он обращает внимание на импридикативность в аксиоме бесконечности, что сильно затрудняет её корректную интерпретацию, на чём мы подробно останавливались в статье ''Несчётность множества континуум, шах и мат''. В конце концов, Сколем совершенно справедливо отмечает, что идея всеобщности индукции выводимой из определённых правил аннотации мыслимых объектов абсурдна. С современных позиций мы легко заключаем, что каждой алгебре: кардинальной, порядковой и индексной, соответствует своя индукция, как эмерджентное свойство алгебраической структуры. Сколем делает попытку описать чем опасен релятивизм в теории претендующей на базис математики, вслед за Пуанкаре вскрывая структурную зависимость индукции, интуитивно чувствуя необходимость переключения алгебраических структур при описании действительности, задолго до того, как это становится наглядным и очевидным.  К сожалению, многие проблемы поднимаемые в данной статье были оставлены без внимания математическим сообществом, что в свою очередь способствовало закреплению интерпретации смешавшей две фундаментально различные алгебры.

Долгое время теория множеств описывала множества математических объектов из других областей математики, что мешало полноценному замыканию аксиоматики в изолированную теорию. Эту проблему удалось элегантно решить Джону фон Нейману в статье ''К введению трансфинитных чисел'' \cite{von-neumann-ordinals}. Фон Нейман предлагает зафиксировать канонический вид порядковых чисел как определённых множеств внутри теории. Эти множества строятся следующим образом: $\emptyset, \{ \emptyset \}, \{ \emptyset, \{ \emptyset \} \}, \dots$, где каждый следующий элемент является множеством всех предыдущих. Таким образом, для построения таких чисел нам нужно лишь пустое множество и аксиома объединения, благодаря чему используя аксиому выбора и схему преобразования мы легко определяем нумерацию элементов любого множества такими числами. Фон Нейман показывает, что для любого порядкового типа таким образом существует один и только один подобный ординал (формально, корректно говорить один и только один тип ординала). Фон Нейман чётко проговаривает, что с формальной точки зрения мы всегда должны иметь ввиду, что всякий ординал представляет собой тип множества содержащего все предшествующие ординалы. Эту мысль он предлагает выражать простым языком: всякий ординал представляет собой множество содержащее все предшествующие ординалы. Он подчёркивает, что это не утверждение о ординалах, а определение, которое можно использовать лишь в контексте заданной трансфинитной индукции. Иными словами, в рамках строгой аксиоматики мы сначала должны прийти к трансфинитной индукции работая с типами ординалов, и лишь затем определить для типов ординалов канонических представителей. Этот ключевой предохранитель от рефикации очень быстро затерялся, и первый этап практически никем не воспроизводился. Между тем он играет ключевую роль в понимании того, как конструируется аксиоматическая теория множеств, в частности, для корректного понимания сути аксиомы бесконечности. Мы предполагаем не существование бесконечного множества, мы предполагаем существование типа бесконечных множеств. У нас принципиально отсутствует какая-либо возможность избежать этого ограничения, если только мы хотим претендовать на непротиворечивость теории, поскольку как показал Сколем сам домен теории неоднозначен. Мы способны зафиксировать лишь определённую алгебраическую структуру заданную формулой. Более того, данный тип может не является объектом нашей теории, поскольку совокупность доменов теории вовсе необязательно является множеством. Таким образом, даже когда мы ищем пересечение бесконечных объектов, чтобы получить минимальный из них мы в любом случае получаем минимальность типа, как универсальное свойство нашей алгебраической структуры, в полном соответствии с классическим построением Кантора. Иными словами, наша аксиоматическая конструкция с формальной точки зрения не позволяет нам уйти дальше структурных типов Кантора или классов эквивалентности Рассела. А значит, для любого объекта нашей теории, если только мы предполагаем существование за ним некоторой реальности (например, в платоновском смысле), доступна смена структурной аннотации, переход в другую аксиоматику, позволяющий переключать позиции референта. И эти альтернативные картины, должны дополнять, а не противоречить сформулированной нами, формируя единый сложный ландшафт математики.

Судя по всему, Джон фон Нейман прекрасно понимал это, поскольку предложенная им в следующих работах система аксиом воспроизводила всю мощь исходной программы Кантора, включая кардинальную и порядковую алгебры \cite{von-neumann-system}, в отличии от аксиоматики Цермело, которая концентрировалась исключительно на порядковой части. Фон Нейман вводит два типа объектов: I-объекты (аргументы или элементы) и II-объекты (функции или классы). Две этих совокупности объектов имеют пересечение: I-II-объекты (дуальные), являющиеся одновременно и аргументами и функциями, которые при погружении в множественный синтаксис становятся собственно множествами. Аксиомы сгруппированы в шесть блоков. Вводные аксиомы фиксируют базовый синтаксис системы, определяя экстенсиональное равенство II-объектов и взаимоотношения между I-объектами и II-объектами: функциональные пары и кортежи аргументов. Аксиомы арифметики вводят базовые комбинаторные конструкции имеющие нормальное представление. Блок аксиом логики позволяет выражать равенство, квантификацию и связь явного и неявного представлений через II-объекты. Весьма любопытно, что данный блок позволяет выражать метатеоретические логические конструкции на внутреннем языке самой системы, мостик который не всегда тривиален, особенно в контексте задач кодирования собственного синтаксиса. Следующая группа аксиом требует, чтобы дуальность всегда выражалась на языке системы, то есть была представима функцией, и вводит ключевое ограничение системы - аксиому ограничения размера: функция не обладает дуальностью, тогда и только тогда, когда существует сюръекция из её области определения в совокупность всех аргументов. Уже на данном этапе аксиоматика фон Неймана блокирует все парадоксы самореференции. Поскольку гарантируется только существование конструктивных формул, циклические формулы, вроде $f = \{ \emptyset, \{ g \mid [f, g] \neq A \wedge [g, g] = A \} \}$ получают статус уравнений на существование, для которых метатеория не гарантирует наличие корней в универсуме. Парадокс Рассела приводит к логическому взрыву, но он происходит в объекте за пределами самой системы. В свою очередь аксиома ограничения размера работает в обратную сторону, она не блокирует парадоксы, она постулирует, что все собственно классы обладают общим свойством определяемым через импредикативность, возникающую из необходимости говорить о всех элементах домена. При этом на данном этапе система сохраняет возможность самореферентных конструкций, например, $[f, f] = f$.

Два следующих блока аксиом описывают погружение объектов теории в пространство с заданной морфологией. Аксиомы бесконечности обеспечивают реализацию морфологии бесконечной вполне упорядоченности для I-II-объектов. При этом альтернативные морфологии бесконечности не блокируются, если совокупность исходных аргументов содержит омнифические целые числа или гросс-объекты, то I-II-объект содержащий все такие числа до $\dot{\omega}$ вполне может существовать. Финальная группа аксиом, включающая аксиому фундирования, строит упорядоченный универсум множеств. Здесь I-II-объекты окончательно становятся множествами и единственными возможными аргументами, однако, даже это не делает структуру вполне упорядоченности категоричной морфологией теории. И на данный факт Джон фон Нейман явно указывает: аксиома фундирования запрещающая класс обладающий свойством $\forall n \in \omega: [\alpha, n + 1] \in [\alpha, n]$ привязана исключительно к структуре множества $\omega$, что допускает бесконечные нисходящие цепочки с альтернативными морфологиями. Иными словами, если мы вводим ещё одну аксиому бесконечности, которая обеспечивает индуктивно-дедуктивную морфологию наборов, постулируя наличие множества $\{ 0, 1, 2, \dots, \dot{\omega} - 2, \dot{\omega} - 1 \}$, то в системе будут сосуществовать две альтернативных морфологии бесконечности. С точки зрения примитивов это просто два различных II-объекта, один ($\dot{\omega}$) содержит дедуктивные элементы, в то время как другой ($\omega$) - нет. Дедуктивные аргументы видны исключительно функциям реализующим морфологию индексных чисел, весь вполне упорядоченный универсум является всего лишь индуктивной подсистемой и совершенно не чувствует дедуктивную подсистему. Поэтому аксиома фундирования становится вполне логичным ограничением на взаимоотношение индуктивной и индуктивно-дедуктивной морфологий, запрещающим однозначное сопоставление чисто индуктивных структур более детальным индексным структурам. Более того, биекция в контексте данной аксиоматики является II-объектом, то есть классом, а не множеством, например, биекция собственного класса ординалов на собственный класс всех множеств, является собственным классом. Поэтому понятие кардинальности в систему фон Неймана очень точно соответствует интуиции Кантора. Погружение в общее пространство с вполне упорядоченной морфологией, гарантирует существование множеств-биекций поддерживающих структуру данного порядка. Здесь работает теорема Кантора, блокирующая существование подобного множества в упорядоченной подсистеме для экспоненциально различных бесконечных множеств, в то время, как возможность альтернативных морфологий, допускает существование класса-биекции (множества с альтернативной морфологией) между ними. Связь логики с функциональными отношениями позволяет иметь классическую логику в изолированных фрагментах теории, в точности как в независимых алгебрах Кантора, не совершая ошибку ложной дихотомии. Интересно как в данном подходе раскрывается роль аксиомы выбора, позволяющей переводить объекты из одной морфологии бесконечности в другую, например, можно вполне упорядочить $\dot{\omega}$: $\{ 0, \omega - 1, 1, \omega - 2, 2, \dots \}$, и упорядочить $\omega$ в индексный тип: $\{ 0, 2, 4, \dots, 3, 1 \}$. Таким образом, Джон фон Нейман аксиоматизирует систему максимально близкую к теории Кантора, строго выделяя полностью внутреннюю перспективу, воплощающуюся на уровне классов, и различные способы погружения в общее внешнее пространство на уровне множеств. Невероятно обидно, что настолько гениальное творение осталось невостребованным, и даже было искажено до неузнаваемости в контексте смешения позиций референта.

Так или иначе, с работами фон Неймана теория получила мощнейший аппарат, способный сохранить различение внутренней и внешней перспектив. Однако последующая традиция акцентировала внимание на упрощённой трактовке его первой работы и прочитала этот аппарат уже через смешанную оптику, типы множеств и классы эквивалентности были реифицированны в математические объекты, многие естественные для конструкции Кантора вещи, такие как $\omega - 1 = \omega$ стали недопустимыми. Переопределение пределов через супремумы, по-видимому почерпнутое из работ Хаусдорфа совершенно скрыло напряжение вокруг объектов $2^{\omega}$ и $\omega^{\omega}$, которое было ещё осязаемо в начале столетия. Мы видим, что у обнаруженного нами изъяна современной теории множеств нет единого автора, скорее причина была в постепенной эволюции коллективного понимания математического сообщества в условиях кризиса, вызванного парадоксами самореференции. Кантор обнаруживает конфликт позиций референта в своей первой работе и посвящает себя его исследованию, в конце концов он находит элегантное решение - разделение кардинальной и порядковой алгебры. Оно сформулировано безупречно в его кульминационной работе, но он явно не связывает его со своими предыдущими работами. Хессенберг пытается упорядочить наследие Кантора и не замечает различие алгебр, каждая из которых имеет своё понятие функции, он сливает обе алгебры в одну интегрируя в них диагональный метод, но при этом завершает программу Кантора введя чисто внешнюю алгебру натуральных операций. Движимый желанием спасти теорию от парадоксов Цермело, предлагает аксиоматизацию тории множеств сохраняющую лишь её упорядоченную часть - порядковую алгебру, где безраздельно властвует теорема Кантора. Хаусдорф, видит напряжение, он склоняется к несчётности порядковых типа $\omega^{\omega}$, но оставляет счётность алгебраических чисел теоремой кардинальной алгебры, как потенциально бесконечную сумму конечных комплексов. Френкель, Серпински и другие принимают эстафету исследуя предложенный Цермело фреймворк. Тем временем, теория множеств набирает популярность, однако, воспроизводимая в наивном контексте она закрепляет противоречивые картины в разных частях математического сообщества. Сколем интуитивно чувствующий опасность закрепления одной алгебраической структуры в качестве фундамента математики предупреждает о возможных последствиях наметившегося движения. Одновременно со Сколемом Френкель указывает на то, что кардинальную арифметику невозможно вывести в аксиоматике Цермело, предлагая схему аксиом преобразование. Даннная схема аксиом действительно позволяет говорить о некоторой кардинальности, для которой необходимо доказать эквивалентность понятию Кантора, однако, данный вопрос остаётся без внимания и интерпретация фиксируется случайным образом. Фон Нейман осознанно реифицирует объекты теории, противоречия которой ещё не были разрешены, что с одной стороны упрощает работу, с другой усложняет анализ проблем. В то же время, он предлагают мощнейший формальный аппарат фиксирующий как кардинальную так и порядковую интуицию Кантора, имеющий строгое разделение между сущностями внутренней перспективы - классами (кардинальными объектами) и сущностями внешней перспективы - множествами (упорядоченными объектами). В этой системе он определяет мощность как начальный ординал класса эквивалентности изоморфных множеств, где биекция является классом, а не обязательно множеством, что полностью соответствует подходу Кантора. В конечном итоге, закрепляется наиболее распространившаяся в сообществе интерпретация, которая просто компилирует все теоремы на аксиоматические рельсы без критического осмысления влияние парадоксов связанных с теоремой о конечном представлении на достоверность доказательств некоторых из них в новой аксиоматике, а впоследствии маскируется внутренним определением предела через супремум. Несмотря на то, что ключевые авторы демонстрируют явную неопределённость в атрибуции теоремы Кантора кардинальной и порядковой алгебре, консенсус сливает обе перспективы в одну, сначала в рамках упорядоченного универсума ZFC, а затем, что особенно трагично, проецирует уже сформированную интерпретацию на гениальную эпистемологически выверенную систему фон Неймана. Судя по всему, данный консенсус впервые кристаллизуется в фундаментальном труде ''Начала математики'' \cite{burbaki-set-theory}, написанным группой математиков под псевдонимом Никола Бурбаки. Они впервые принципиально отказались от наивного введения, начав с формального языка и аксиом, и скомпилировали с этим фундаментом практически весь массив доказательств теории множеств. Они поместили теорию множеств в основания математики, возможно, в частности, для того, чтобы никто не прошёл мимо парадоксов с ней связанных, интуитивно всё ещё чувствуя недосказанность современного им формализма. Вряд ли они были авторами закрепившейся интерпретации, скорее они зафиксировали представления, которые уже сформировались в сообществе, не обратив внимания на ключевые неопределённости, связанные с различием кардинальной и порядковой алгебр, а также отношением к ним теоремы Кантора. Это явление вполне объяснимо колоссальной инерцией классических представлений о запрете актуальной бесконечности, которые помогли оправдать счётность $\omega^{\omega}$ с помощью потенциально бесконечного множества функций с конечным носителем. Мы подсознательно хотим, чтобы множество натуральных чисел было замкнутым, не позволяя никаким комбинаторным операциям выйти за его пределы, настолько, что распространяем это стремление на случай внешнего описания, когда множество натуральных чисел становится сформированной целостностью. Такое стремление как минимум нуждается в формализации.

Мы наблюдаем сложнейшую историко-эпистемологическую драму, в которой просматривается удивительная тактическая беспомощность альтернативных направлений математики. Более ста лет подряд сотни ученых, интуитивно чувствуя изъян теории множеств, направляли всю мощь своей критики на её самый укреплённый форпост - диагональный аргумент и теорему Кантора. Пытаясь опровергнуть несчётность континуума, построить прямую нумерацию вещественных чисел или разрушить логику диагонального метода, они неизменно терпели крах, разбиваясь о синтаксический гранит строгого формализма. Подобная тактика обречена на провал, поскольку внутри правил игры упорядоченного универсума теорема Кантора абсолютно корректна и неуязвима. Кардинал в современных аксиоматиках - это не абстракция Кантора от природы элементов и структуры, а начальный ординал, фиксирующий границу экспоненциальной области. Теорема Кантора является честным внутренним стражем иерархии рангов, защищающим теорию от коллапса, и атаковать её на её собственном поле - математическое самоубийство. Ирония ситуации заключается в том, что, пока критики безуспешно штурмовали неприступную цитадель, тылы официальной аксиоматики оставались уязвимы. Настоящая ахиллесова пята теории скрывается там, где теория совершает логическую контрабанду, маскируя её под очевидность - в теореме о счётности счётного объединения счётных множеств. Причина этого кроется в когнитивном искажении - свойстве нашего мышления коллапсировать к внутреннему описанию, предпочитая его внешней сложности, которое, судя по всему, является эволюционной оптимизацией помогающей нам экономить ресурсы. Данный феномен на сегодняшний день имеет бесчисленное множество подтверждений в когнитивистике и хорошо демонстрируется уже упомянутым в настоящей работе экспериментом Питера Уэйсона \cite{wason-246}. Большинство испытуемых не угадывают правило, поскольку экспериментируют с тройками чисел заведомо удовлетворяющих более строгим правилам, например, $(4, 6, 8), (6, 8, 10), (3, 5, 7)$ и так далее. С точки зрения современной когнитивистики, такое поведение объясняется стратегией положительного теста и принципом когнитивной экономии, мозг фиксируется на внутренних объектах имеющегося фокуса, избегая затратного расширения пространства поиска альтернативных гипотез \cite{positive-test, rational-selection}. Подавляющее большинство людей, в том числе и критики современных теорий множеств, воспринимают внутреннюю кардинальную логику этого доказательства естественным образом, не находя никаких оснований для сомнения даже в контексте аксиоматики упорядоченного универсума. Можем ли мы винить математическое сообщество начала прошлого века, за то что оно предпочло объяснить доказательство, которое считалось верным, введением новых сущностей, вместо того, чтобы критически переосмыслить его в контексте новой аксиоматики, если даже в настоящее время, когда накоплено изрядное количество проблем теории множеств, критики пропускают это доказательство, но видят проблему в доказательстве теоремы Кантора, требующей внешней перспективы? Более того, тест Уэйсона играет какую-то невероятно поучительную роль во всей этой истории, ведь именно этот эффект ответственен за твердокаменную веру мейнстрима в то, что кроме индуктивной морфологии бесконечности не существует никаких других. Именно редукция мышления к восприятию свидетельства (известной альтернативы) в качестве правила, а не как одного известного варианта из многих неизвестных, привела к тому, что мы долгое время не могли осознать отсутствие общезначности в современных теориях множеств. В то же время, деконструкция рассуждения о счётности счётного объединения не требуется привлекать сложные внешние концепции; порочный круг обнаруживает себя на уровне базовой индуктивной структуры, доступной для анализа любому математику.

Рассмотрим индуктивную последовательность вложенных ординалов, где каждый элемент представляет собой расширение размерности полинома (базы носителей функций): $\omega \subset \omega^2 \subset \omega^3 \subset \dots \subset \omega^n \subset \dots$. Каждый изолированный слой $\omega^n$ описывает множество функций со строго конечным носителем размера $n$. Объединим эту последовательность по всем конечным степеням: $\mathcal{F} = \bigcup_{n < \omega} \omega^n = \omega^{\omega}$. Поскольку для любого элемента последовательности степень $n$ конечна, мейнстрим утверждает, что предел $\mathcal{F}$ всё ещё представляет собой множество функций исключительно с конечным носителем. Однако комбинаторика построения подмножеств пробегает по всему множеству конечных степеней, которое само по себе бесконечно ($\omega$). Возникает двойной стандарт: с одной стороны, классическая индукция открыто постулирует, что объединение конечных элементов порождает бесконечную целостность ($\bigcup_{n < \omega} n = \omega$), но с другой - директивно запрещает носителю комбинаций в пределе стать бесконечным ($\omega^{\omega} \neq \omega^{\bigcup_{n < \omega} n}$). Но если степень представляет собой не $\omega$, а какой-то иной объект (не конечный, но и не бесконечный в том же смысле, что и $\omega$), который мы условно обозначаем как $\omega$, то для него необходимо вывести ординальную арифметику. Всё стройное здание трансфинитных объектов Кантора зиждется на одном ключевом свойстве - самоподобии, которое он вместе с Дедекиндом позаимствовал у непрерывного. Когда для защиты возникшей интерпретации мы вводим промежуточный объект, у нас просто отсутствует какое-либо свойство отличающее его от $\omega$, что с точки зрения математической логики означает тождественность данного объекта $\omega$. Однако, классика совершает прыжок веры и утверждает, что комбинаторный объём полученного объекта каким-то образом удерживается в рамках полиномиальной области и не допускает смещение мощности в следующую экспоненциальную область $\omega^\omega$. Доказать, что множество функций с конечным носителем отлично от множества всех функций, можно только опираясь на теорему о счётном объединении счётных множеств. Но саму эту теорему в рамках упорядоченного универсума удаётся доказать лишь в том случае, если мы заранее постулируем существование таких промежуточных объектов, мощность которых не совершает экспоненциальный переход. Мы сталкиваемся с чистым порочным кругом. Аксиома выбора здесь оказывается бессильной. Она способна лишь констатировать легитимность одновременного выбора и проиндексировать его ординалами, но она принципиально не может гарантировать, что последовательность этих индексов останется счётной. Современная теория удерживает счётность счётного объединения не за счёт логики аксиом, а путём волевого запрета бесконечному пределу вести себя как бесконечность в точке перехода от полиномов к степенным башням. По всей видимости такие промежуточные объекты возникают из-за того, что мы не осознаём перспективу, с которой рассматриваем объекты. В порядковой алгебре Кантор определяет первое трансфинитное порядковое число, которое делает пределы и бесконечные ряды объектами данной алгебры. Бесконечное суммирование и объединение множеств происходит внутри теории, независимо от того выразимы или нет все объекты этих операций в мета-теории. Но когда мы смотрим на это из мета-теории, мы видим что не можем описать все элементы подобных операций, $\bigcup_{n < \omega} \omega^n$ представляется нам как объект, где все множества конечны и могут быть выражены, но внутри теории это не так, поскольку $\omega$ - это бесконечный объект определяющий бесконечное объединение вариативных позиций, а наше ограничение на выразимость, существует только внешне. Иными словами, из-за отсутствия языка описания переключения перспектив мы атрибутируем ограничение внешней мета-теории внутренней, что и порождает промежуточные объекты.

Алгебраические структуры не существуют в изоляции, а пересекаются и дополняют друг друга. Идеологический коллапс произошедший в теории множеств привёл к блокировке альтернативных подходов, которые могли бы пролить свет на точность формулировок и границы их применимости. Диагональный аргумент предъявляет элемент, которому невозможно сопоставить объект в исходном множестве, вещественное число, для которого нет натурального числа. Если мы упорядочим множество упорядоченных пар натуральных чисел, сравнивая сначала первые элементы, а затем вторые, например, $(1, 1) < (1, 2)$, $(32, 1) > (25, 246)$, то не сможем сопоставить натуральное число паре $(2, 1)$ или $(153, 1)$. Это наблюдение указывает на то, что дело вовсе не в множестве как таковом, а в его структуре. Это наблюдение можно усилить, указав на проблему с сохранением индуктивных инвариантов в стандартной интерпретации теории множеств. Так диагональный метод сохраняет инвариант отличия диагонального элемента от всех предыдущих. Но если мы рассмотрим доказательство счётности пар натуральных чисел и функцию сопоставления: $x = n + (n + m - 1)(n + m - 2)/2$ \cite{cantor-countable}, то для $m = 1$ мы имеем зависимость только от одного множества: $x = n*(n + 1)/2$, значит множество индексирующее пары содержит элемент равный сумме всех предыдущих индексов из этого множества, следовательно множество индексирующее пары содержит и сумму вообще всех индексов данного множества, то есть из $1 + 2 + 3 + \dots = \omega$ имеем $\omega \in \omega$. Это указывает на то, что некоторые построения требуют воспроизводимости более детальной структуры, чем структура ординалов. Нарушение инвариантности связано с потерей информации о морфологии, это сознательное игнорирование части структуры, оно возможно только через явное включение внутренних акцентов в теорию и не может иметь догматического статуса, который приобрело в аксиоматических теориях множеств. За век победного шествия современной интерпретации теории множеств накопилось огромное количество противоречий и нестыковок подобного рода, многие исследователи обращали внимание на противоречивость подходов теории множеств, объём такой критики уже требует для полноценного изложения целой книги и скрупулёзная работа по созданию справочника изъянов теории была проделана Вольфгангом Мюкенхаймом \cite{transfinity}. Многие критики исходят из предположения, что противоречивой оказывается именно концепция актуальной бесконечности, но как мы видели, в действительности, противоречия возникают не из-за актуальной бесконечности как таковой, а из-за попытки совмещать внутреннюю и внешнюю перспективы в рамках одного построения. Возможно самым правильным вариантом, стал бы отказ от современных представлений о бесконечности в попытке пересобрать фундамент математики имея чёткое понимание причин и механизмов кризиса. Однако, такой подход вряд ли представляется возможным, поскольку требует колоссального интеллектуального напряжения от сообщества. В этой связи, наиболее продуктивным авторам видится честная формализация и изоляция найденных дефектов.

В результате предпринятого исследования нам удалось выяснить, что сами по себе аксиоматики ZFC или NBG не содержат некорректных утверждений. Выбранные формы аксиом лишь затрудняют их однозначную интерпретацию и понимание механики систем. К контринтуитивным результатам приводит лишь скрытое допущение о том, что внутреннее описание распространяется на ординальную операцию возведения в степень, что мы лаконично записывали как $\aleph_0 = \omega$. Таким образом, мы можем полностью сохранить обе теории в их современном виде со всеми полученными результатами расширив их аксиомой легализующей данное скрытое допущение. В отличие от кардинальной арифметики в ординальной транзитивное отношение порядка задаёт направление, что позволяет ввести понятие предела. Аксиома выбора в интерпретации требующей воспроизводимости гарантирует, что всякое множество при любом переупорядочении остаётся в пределах класса эквивалентности изоморфных множеств. Таким образом, когда мы говорим об упорядоченном универсуме возникает неоднозначность с объединением элементов всех элементов множества. В самом деле, если мы возмём некоторое вполне упорядоченное множество, например, $\omega$, и построим эквивалентное ему множество множеств, каждое из которых содержит все комбинации соответствующего элемента со всем комбинациями предыдущих, то есть $f(x) = \mathcal{P}(x + 1)$, фактически воссоздавая степенной предел Кантора: $\lim 2^x = 2^{\omega}$, то по определению в комбинациях объединения всех элементов полученного множества участвуют все элементы исходного ($\omega$). Однако, нам нужно сделать так, чтобы индуктивно мы не могли получить степень в множественном смысле, иначе объединение всех элементов не будет эквивалентно исходному множеству. Фактически мы хотим формализовать нашу веру в комбинаторную замкнутость упорядоченного универсума, то есть что предел последовательности элементов эквивалентных самой последовательности всегда эквивалентен этой последовательности. Наиболее элегантной для данного утверждения представляется следующая форма:
\begin{equation}
\begin{gathered}
\forall X, x \; \bigl(x \in X \to x \sim X \wedge (X \cup x) \sim X \bigr) \to \bigcup X \sim X
\label{axiom:limitation-axiom}
\end{gathered}
\end{equation}

Иными словами, если всякий элемент $x$ некоторого множества $X$ эквивалентен этому множеству и объединение $x \cup X$ также эквивалентно ему, то объединение всех элементов этого множества $\bigcup X$ эквивалентно $X$. Преимущество данной формы в том, что она избегает всякого представления о бесконечности и формально просто фиксирует свойство алгебраической структуры: если все элементы множества являются нейтральными элементами объединения в смысле эквивалентности, то их объединение эквивалентно самому множеству. В конечном случае эта аксиома сводится к простому факту, что пустое множество эквивалентно объединению своих элементов. Но когда мы добавляем к системе аксиому бесконечности, кардинальные числа становятся нейтральными элементами объединения для самих себя, что замыкает мультипликативную структуру, а также все уровни воспроизводства аддитивной структуры, в изолированных классах эквивалентности. Данная аксиома не вступает в противоречие с теоремой Кантора, поскольку она всего лишь записывает то, что уже считается истинным - отличие индуктивного построения от множества степени. Именно поэтому она не делает теории сильнее или слабее, а лишь формализует молчаливое предположение ответственное за всю контринтуитивность связанную с ней. Не сложно видеть, что данная аксиома выражает ничто иное как третий принцип порождения Кантора - принцип стеснения или ограничения. Как мы видели Цермело формализовал два первых принципа порождения Кантора, оставив без внимания третий. Данное наблюдение подчёркивает глубину анализа Кантора, осознающего, что сформированная веками интуиция потенциальной бесконечности, какой бы очевидной она ни казалась, нуждается в явной фиксации. Поэтому мы называем предложенный постулат {\it аксиомой стеснения или ограничения (limitation axiom)}, тогда классические интерпретации ZFC и NBG становятся полноценными формальными системами ZFCL и NBGL, где воспроизводимы все результаты современной теории множеств. Стоит также отметить, что аксиома ограничения не полностью фиксирует современную интерпретацию, в частности, она не гарантирует различие мощностей множеств функций с конечным носителем и множеств всех функций. Данная интерпретация остаётся лишь неформальной иллюстрацией причин невозможности индуктивного выхода за пределы фиксированной мощности. 

Теории ZFCL и NBGL безусловно представляют невероятный интерес для понимания механики самовоспроизводства порочного круга заложенного в их основание и масштабов, которых может достигать данное явление. Необходимо подчеркнуть разницу между порочными кругами и прыжками веры, которые мы часто вынуждены делать, так как не имеем возможности установить исчерпывающие начала для всех понятий. Прыжки веры осуществляются не на пустом месте, но являются фиксацией понимания явлений почерпнутого из непосредственного опыта. Это честная фиксация ограниченности языка и понятийного аппарата, кристаллизация формулировки, которая доступна для текущего уровня развития, но которая непременно допускает верификацию. Все аксиомы ZFC или NBG кроме аксиомы бесконечности верифицируемы в реальном опыте. Даже аксиома бесконечности по своей сути является фиксацией образа направленной абстракции, она выражает не произвольную мысль, а описание свойства нашего мышления видеть цель последовательного процесса. Более того, аксиома стеснения сама по себе - это аксиома замкнутости алгебраической структуры, гарантирующая изолированность домена объектов относительно определённого набора операций. И Кантор демонстрирует как корректно использовать данную аксиому: порядковые числа формируют замкнутую область - второй числовой класс, где функции являются внутренними алгебраическими объектами данной области, их невозможно использовать для того чтобы выйти за пределы области, поскольку мы потребовали её замкнутости. Теорема Кантора является внутренней теоремой этой области, как и любая другая теорема сформулированная в терминах таких функций. Именно поэтому для перехода к следующему числовому классу Кантор меняет позицию референта чтобы использовать весь предыдущий класс как целостный объект. Порочный круг возникает тогда, когда мы пытаемся оправдать своё стремление одновременно замкнуть числовой класс и позволить внутренним объектам - функциям применяться к самому классу и выходить за его пределы. Когда мы хотим получить то, чего не существует в опыте - систему, которая одновременно изолированна и способна внутренними средствами туннелировать за свои границы. Когда мы пытаемся сделать внутреннюю сущность - начальный ординал одновременно внешним контуром замыкающим алгебраическую структуру, объектом которой он является. У Кантора кардинал является внешним объектом по отношению к алгебре порядковых чисел, поэтому он корректно замыкает структуру второго числового класса. У фон Неймана функции первичны, что обеспечивает возможность существования биекции как класса не ограниченного порядковой морфологией. Но в современной интерпретации аксиома бесконечности ответственна за существование бесконечных функций, которые нужны для определения класса эквивалентности изоморфных множеств, в результате бесконечное множество одновременно и внешне замыкает этот класс (даёт смысл понятию всех элементов класса) и является его внутренним объектом. Иными словами, мы в буквальном смысле говорим о числе, которое представляет собой все числа. Именно такая вера, опирающаяся на неформализованную интуицию потенциально бесконечного формирует порочный круг, упаковывающий внешний актуально бесконечный объект в привычные рамки внутреннего описания. Порочный круг возникает только тогда, когда мы принимаем утверждение молчаливо, не осознавая этого, и одновременно используем его следствия для доказательства его самого - что и произошло в современной теорией множеств.

Ирония ситуации заключается в следующем. Мы так модифицировали свой логический аппарат, что с формальной точки зрения мы не можем доказать противоречивость принятой системы аксиом, да она содержит порочные круги и основывается на уязвимых постулатах, но строгое доказательство противоречивости данной интерпретации без дополнительных предположений, судя по всему, невозможно. В строгом логическом смысле теорема Гёделя конечно же фиксирует противоречивость теории, её зависимость от внешних утверждений, которые не доказуемы, а значит не являются истинными, следовательно по закону исключённого третьего являются ложными, обеспечивает тривиальность нашей системы аксиом. Однако, мы переворачиваем картину и говорим, что это не фатальная противоречивость, а особый легитимный случай условной противоречивости, что истинность - это имманентное свойство утверждений, которое лишь обнаруживается в доказательствах, но не является неотделимым аттрактором самих доказательств. Следует признать, что здесь мы корректно используем механизм переключение позиции референта вынося понятие истины во внешний контур. Однако, в таком случае, мы создаём проблему куда более сложную, ведь если для любой достаточно сложной теории существуют внешние истины, для неё могут существовать и внешние противоречия (ситуации когда одно построение требует истинности одного утверждения, а друго построение требует истинности его отрицания), но тогда достоверность теории определяется не только внутренней, но и внешней непротиворечивостью. В итоге мы получаем картину, в которой достоверность теории зависит от внешней консистентности, к доказательству которой на сегодняшний день мы даже не знаем как подступиться. По какой-то причине, мы отказываемся видеть довольно очевидную возможность. Когда мы строим теорию мы всегда опираемся на некоторую систему понятий и предположений, которая имеет свою разрешающую способность, то есть позволяет говорить, о производных объектах, но сами фундаментальные сущности системы требуют прыжков веры и их онтология и эпистемология не поддаются описанию. Если в такой системе возникают противоречия, то эти противоречия обусловлены исключительно дизайном самой системы, то есть именно те свойства, которые мы заложили в систему в процессе проектирования проявляются в ней в форме таких противоречий. Когда вместо того, чтобы понять причины коллизий, мы исправляем ситуацию с помощью аксиом цензуры, исключая из нашего описания существенное количество объектов, мы избавляемся от возможности прийти к противоречию внутри системы, то есть мы просто обрываем цепочки рассуждений на некотором фрагменте, поскольку дальнейшая часть ведёт к коллапсу. Но перестают ли от этого такие исключённые фрагменты существовать? Ведь всё, что мы сделали - это запретили себе их видеть. Иными словами, исходные фундаментальные компоненты нашей системы противоречат друг другу, мы ослабляем нашу теорию, чтобы эти противоречия не возникали в ней самой. Но становится ли теория от этого непротиворечивой? Или же мы просто получаем непротиворечивый фрагмент противоречивой теории? Эти ключевые вопросы противоречивости избегаются современным мейнстримом. Для ответа на них, требуются более сильные теории, позволяющие рассмотреть, исключают ли аксиомы цензуры источник противоречий, или же просто закрывают путь к ним, оставляя их частью теории.

Так или иначе, это промежуточное положение становится фундаментальным ограничителем на силу самой теории, любая попытка существенного усиления тут же приводит к коллапсу в форме антиномий самореференции. Мы имеем ряд фундаментальных проблем, которые оказываются независимыми от теории и требуют дополнительных аксиом для своего разрешения. Но всё что нам доступно - это локальное усиление для конкретных задач, без возможности общепринятых прорывов. Путь к преодолению стагнации не может находится внутри самой теории, а должен определятся её взаимодействием с альтернативными системами, такими как, новые основания Куайна, сюрреальные числа или гросс-множества. Математические теории не существуют в вакууме, они сосуществуют друг с другом в экосистеме человеческого мышления, где различные подходы не опровергают друг друга, а дополняют. Уязвимость аксиоматических теорий множеств в контексте разнообразия математических взглядов толкает сообщество к агрессивной цензуре по отношению к другим интерпретациям и тоталитаризму через догматизацию правильных взглядов. Проблема в том, что индексная морфология плохо интегрируется с классическими подходами смешивающими внутреннее и внешнее описание, равенство $2^{\omega} = \omega$ опровергается простым фактом того, что индексное число $\dot{\omega} - 1$ уже бесконечно имеет порядковый тип $\omega$, вследствие чего $2^{\dot{\omega}} = 2^{\dot{\omega} - 1} \times 2 = \omega \times 2$ и равенство $2^{\omega} = \omega$ ломает согласованность структур. Кантор в своих работах многократно подчёркивал важность такой согласованности, поскольку именно она обеспечивает целостность и достоверность наших теорий. Именно структурная когерентность различных способов описания обеспечивает необходимую обоснованность и надёжность тех актов редукции и прыжков веры, которые лежат в основе наших взглядов. Такая транстеоретическая воспроизводимость и отличает науку от жречества, обеспечивая фальсифицируемость и верифицируемость аксиоматического базиса. В то же время, упорство мейнстрима в отстаивании случайно закрепившейся упрощённой интерпретации, отчётливо демонстрирует насколько человеческий вид превосходит все остальные существа на планете в своём упрямстве. Одним из признаков догматического подхода в современной теории множеств является распространенная практика оправдывать аксиоматические системы ZFC и NBG преимущественно их практическими достоинствами: они хорошо себя зарекомендовали, приводят к значительным результатам, удобны в использовании и стали стандартом де-факто в большинстве разделов математики. Такой прагматический аргумент практически не встречается в обосновании фундамента других математических дисциплин, например, математического анализа, алгебры или комбинаторики, где аксиомы обычно мотивируются интуицией, историческим развитием или внутренними логическими соображениями. Подобный стиль оправдания напоминает риторику, характерную для идеологических систем, где легитимность доктрины подкрепляется её успехом в обществе - от религиозных сект до авторитарных политических режимов.

Фиксация скрытых допущения позволяет нам сформировать кластер альтернативных теорий, которые оказываются невероятно согласованными друг с другом и с оригинальными идеями Кантора. Классическая идея потенциальной бесконечности выделенная Кантором в кардинальную алгебру связана с такими аксиоматическими системами как новые основания Куайна, теория множеств с позитивной логикой, а также другими паранепротиворечивыми логиками. Внутренняя морфология кардинальных объектов замкнута и не позволяет выхода наружу, так внешне счётный универсум новых оснований Куайна может быть замкнутой внутренней областью несчётного универсума. В большинстве современных подходов замкнутость достигается за счёт ограничения выразительной силы формул, так позитивная логика не позволяет сформулировать отрицания избегая самой возможности противоречий, а новые основания требуют стратификации формул, делая ряд утверждений, например, о самопринадлежности, бессмысленными в теории. Для таких система характерно невозможность глобальной структуры допускающей внешнее сравнение, в частности, аксиомы выбора здесь возможна лишь во внутренней невоспроизводимой интерпретации. С чисто внешним взглядом проблем также не возникает, его реализуют гросс-множества, класс сюрреальных чисел, подкласс которого - наборы, мы можем рассматривать как множество индексных чисел, и другие системы. Теорема Кантора в таких системах всегда справедлива. Кроме того, выполняется аксиома Евклида - целое всегда больше своих частей и в стандартном арсенале обычно отсутствуют внутренние операции. Такие операции нужны только для демонстрации рефлексивности - внутреннего самоподобия, чтобы показать, что речь идёт об одной и той же бесконечности. Таким образом, для замыкания картины нам нужны только классические порядковые числа в интерпретации Кантора. Обе классические теории ZFC и NBG с интерпретацией доминирования внешней составляющей в ординальном возведении в степень прекрасно справляются с этой задачей. Возможно, в современном контексте данная интерпретация остаётся не очевидной и требует явной фиксации такого выбора в постулате, для чего можно воспользоваться следующей формой:
\begin{equation}
\begin{gathered}
\forall X, I, f \; \bigl(\emptyset \in I \wedge \forall i \; (i \in I \to (i \cup \{ i \} \in I \wedge f(i) = X^i)) \bigr) \to \bigcup_{i \in I} f(i) \sim X^I
\label{axiom:secure-induction}
\end{gathered}
\end{equation}

Иными словами, для всякой пары множеств $X, I$ и функции таких, что всякий элемент индексирующего множества $I$ отображается функцией в степень $X^i$ (множество функций из $i$ в $X$) и множество $I$ транзитивно замкнуто, имеет место эквивалентность объединения индуктивных значений функции со степенью $X^I$. По-видимому, для данной цели мы не можем избежать концепции бесконечности в формулировке аксиомы, поскольку нам необходимо связать аксиоматические объекты теории с индуктивными конструкциями, которые бесконечны по своей природе. В частности, для $\omega$ имеем $\bigcup_{x \in \omega} \omega^x \sim \aleph_0^{\aleph_0}$. Фактически мы явно постулируем, что множество функций с конечным носителем эквивалентно множеству всех функций. Мы можем попытаться сохранить счётность порядковых чисел вида $n^{\omega}$, где $n$ - конечное число, как это сделал Хаусдорф, однако, данный путь ведёт к довольно сложным механизмам переключения перспектив. Данная аксиома называется {\it аксиомой безопасной индукции (secure induction axiom)}, соответственно аксиоматики включающие такую аксиому обозначаются ZFCS и NBGS. 

Как было показано в статье ''Несчётность множества континуум, шах и мат'' невозможность биекции между $\omega$ и $2^{\omega}$ практически никак не сказывается на порядковой арифметике, поскольку $2^{\omega} = \omega^{\omega}$, что сохраняет структуру нормальной формы ординалов. Неравенство $\omega < 2^{\omega}$ легко обосновать, как в случае Канторовского опрделения пределов через бесконечные ряды, так и в случае современного подхода через супремумы. В самом деле, как супремум так и предел определяются в первую очередь множеством, где задана последовательность: $\lim_{2^i \in \omega} 2^i = \omega$ и $\lim_{2^i \in 2^{\omega}} 2^i = 2^{\omega}$. Если множества различны, то различны и пределы. На первый взгляд, может показаться странным, что одно и то же множество чисел $\{2^i | i \in \omega \}$, имеет разный предел в зависимости от того, какому именно множеству мы его атрибутируем $\omega$ или $2^{\omega} = \omega^{\omega}$, ведь порядковых чисел $2^i$ с конечной степенью $i$ превосходящих $\omega$ действительно не существует. Однако, в порядковой алгебре мы сознательно отказываемся от части структурной информации об области непосредственно предшествующей предельным числам для того, чтобы иметь свойство самоподобия маркирующее бесконечность. Это область импредикативной сингулярности, не доступная нашему описанию и с формальной точки зрения мы не можем делать никаких предположений относительно её свойств, а должны обнаруживать эти свойства в результате изучения объектов теории. В частности, из теоремы Кантора следует $\omega < 2^{\omega}$, а значит подпоследовательность $\{2^i | i \in \omega \}$ внутри $\omega$ имеет один предел, в то время как внутри $2^{\omega}$ другой, в силу того, что область сингулярности данной последовательности различна в двух этих контекстах. Мы уже неоднократно сталкивались с данным эффектом, связанным с установленным Кантором простым фактом того, что области предикативности вполне упорядочения связаны с аддитивной структурой, в то время как такие области в структуре континуума требуют более высоких мультипликативных уровней воспроизводства аддитивной структуры, начинающихся с $2^{\omega}$. С современных позиции мы всегда можем сменить позицию референта и перейти к индексному описанию, в котором данный эффект становится невероятно наглядным. Действительно, все числа порядковой сингулярности $2^{\omega - 1}, 2^{\omega - 2}, \dots$ не попадают в области регулярности порядковых чисел между $\omega$ и $2^{\omega}$. Мы не найдём их ни в окрестности $2\omega$, ни в окрестности $\omega^3$ и вообще ни в какой окрестности полинома $\omega$. Таким образом, альтернативные морфологии бесконечности помогают нам понять, что у нас действительно нет никаких основания ожидать, что сингулярность мультипликативной структуры ($x^1, x^{1 + 1}, x^{1 + 1 + 1}, \dots$) - первого уровня воспроизводства аддитивной структуры, будет себя как то проявлять в выстроенной нами аддитивных областях регулярности.

Несложно видеть, что классическое определение кардинальных чисел через начальные ординалы класса изоморфных множеств в такой теории просто классифицируют экспоненциальные области ординалов. Большого смысла выделять такие объекты нет, поскольку они не будут подчинятся кардинальной арифметике открытой Кантором. Тем не менее, несмотря на то, что в упорядоченном универсуме нет объектов способных описать соответствие между $\omega$ и $\omega^{\omega}$ кардинальные числа в смысле Кантора допускают определение и здесь. В самом деле, мы можем воспользоваться принципом стеснения Кантора и сказать, что алгебраически замкнутая область порядковых чисел называется {\it кардинальной областью}, где всем объектам соответствует определённое {\it кардинальное число}. То есть, всё что мы можем получить из бесконечного порядкового числа $\omega$ через сложение, умножение, возведение в степень, бесконечные последовательности, ряды и произведения, имеющие порядковый тип полученный таким образом, образует область первого трансфинитного кардинального числа $\aleph_0$. Из такого определения непосредственно следует вся кардинальная арифметика: $\aleph_0 + 1 = 1 + \aleph_0 = \aleph_0 + \aleph_0 = \aleph_0$, $2\aleph_0 = \aleph_0 2 = \aleph_0$, $\aleph_0^2 = 2^{\aleph_0} = \aleph_0$. Чтобы получить следующее кардинальное число, нам нужно сменить позицию референта и рассмотреть эту замкнутую алгебраическую структуру как целостный объект, заметить, что такой объект вполне упорядочен, а значит является порядковым числом, которое по определению не является внутренним для данного объекта (классическое доказательство несчётности чисел второго числового класса Кантора). Этот объект удовлетворяет аксиоме бесконечности, а значит существует как объект нашей теории (например, по аксиоме выделения). Таким образом, кардинальная алгебра воспроизводится в упорядоченном универсуме в полном объёме, однако, мы не можем говорить о существовании биекции между объектами одной мощности, например, множество вещественных чисел определённых через сечения Дедекинда оказывается счётным, но не при этом не изоморфным множеству натуральных чисел. Данная теория сохраняет классическую интуицию сравнения бесконечных количеств через внутреннюю биекцию алгебраической структуры - принцип Юма, что обеспечивает их Дедекиндову бесконечность. При этом такая интерпретация выводима из индексных теорий, например, теории наборов или сюрреальных чисел, и следовательно всегда допускает восстановление отброшенной информации о морфологии и переходу от порядковой алгебры к индексной. Иными словами, она служит связующим звеном объединяющим различные структурные аннотации бесконечных объектов в единую картину.

Чтобы работать с внешними операциями в принципе можно использовать один из современных подходов нестандартного анализа. В то же время, вдохновляясь концепцией фон Неймана можно существенно упростить аппарат таких теорий. В самом деле, аксиома бесконечности в её классической форме делает гораздо более скромное утверждение, чем принято считать. Она не постулирует существование объекта, который абсолютно рефлексивен. Она утверждает лишь следующее: в рамках данной системы существует подсистема, реализующая рефлексивную бесконечную морфология, а все отображения, которые мы обычно рассматриваем, определены строго внутри данной морфологии. Как мы видели исходная система аксиом фон Неймана идеально подходила для расширения альтернативными структурами бесконечности, что позволяло легко ввести полностью внешнюю индексную морфологию простым добавлением альтернативной аксиомы бесконечности. В современном контексте, когда в NBG как и ZFC доминирует отношение принадлежности $\in$ наиболее естественным представляется введение концепции расслоения данного отношения. Формально мы пытаемся эмулировать функциональную природу данного отношения наследуемую из оригинальной аксиоматики фон Неймана. Пусть все аксиомы NBG и ZFC, кроме аксиомы бесконечности справедливы для любого отношения принадлежности, то есть для характеристических функций с произвольной морфологией. Теперь выделим {\it индуктивное} отношение принадлежности $\epsilon$ и сформулируем для него аксиому бесконечности:
\begin{equation}
\begin{gathered}
\exists I \; (\emptyset \; \epsilon \; I \wedge \forall x \; (x \; \epsilon \; I \to x \cup \{ x \} \; \epsilon \; I))
\label{axiom:ordinal-infinity}
\end{gathered}
\end{equation}

Фактически мы воспроизвели погружение теории в пространство с общей вполне упорядоченной морфологией бесконечности. Аксиомы выбора теперь способно подхватить данную морфологию и индуцировать её на любом объекте теории. На данном этапе системы остаются полностью эквивалентными современным версиям. Теперь выделим альтернативное {\it индуктивно-дедуктивной} отношение принадлежности $\varepsilon$, связанное с $\epsilon$ аксиомой:
\begin{equation}
\begin{gathered}
\forall X, x \; (x \; \epsilon \; X \to x \; \varepsilon \; X)
\label{axiom:induction-deduction}
\end{gathered}
\end{equation}

То есть, индуктивно-дедуктивное отношение принадлежности $\varepsilon$ имеет индуктивную составляющую полностью эквивалентную $\epsilon$. В свою очередь, индуктивное отношение $\epsilon$ просто не видит дедуктивных множеств, словно они не существуют для него. В контексте оригинальной версии аксиоматики фон Неймана это выглядит естественно, поскольку нет выделенного объекта, описывающего отношение принадлежности, морфология же конкретного отношения появляется как класс функций учитывающих те или иные совокупности объектов. Будем называть {\it индексалами} множественные конструкции полностью эквивалентные ординалам, но описываемые в контексте индуктивно-дедуктивного отношения принадлежности $\varepsilon$. Определим для них операцию вычитания: $x - 1 \equiv \bigcup x$. Теперь ничто не препятствует погружению наших теорий в общее пространство с индуктивно-дедуктивной бесконечной морфологией:
\begin{equation}
\begin{gathered}
\exists I \; \bigl(\emptyset \; \epsilon \; I \wedge \forall x, y \; (x \; \epsilon \; I \wedge y = x \cup \{ x \} \to y \; \epsilon \; I \wedge (I - y) \; \varepsilon \; I)\bigr)
\label{axiom:indexal-infinity}
\end{gathered}
\end{equation}

Таким образом, $\dot{\omega}$ становится полноценным объектом нашей теории, интегрирующим в неё альтернативные подходы, такие как сюрреальные числа и гросс-объекты. Аналогично индуктивному случаю, аксиома выбора подхватывает индуктивно-дедуктивную морфологию, позволяя индуцировать её на любых объектах теории. Что позволяет полноценно оперировать как внешними так и внутренними функциями. Предложенное здесь расширение выглядит существенно проще существующих способов введения нестандартного анализа и может быть использовано как с аксиоматиками ZFCL и NBGL, так и с аксиоматиками ZFCS и NBGS. Обратим внимание, что само по себе введение альтернативной морфологии бесконечных объектов не защищает нас от желания распространить кардинальную индуктивность на объект $\omega$, поскольку оно происходит на границе конечной области регулярности и области сингулярности, которая остаётся незатронутой. Только связав данные морфологии друг с другом новыми аксиомами, мы сможем исключить ту или иную возможность. Заметим также, что аксиома фундирования не конфликтует с новой морфологией бесконечности, как можно было бы ожидать, даже в оригинальной формулировке, которая в данном контексте более сильна, поскольку она явно запрещает нисходящие цепочки индексированные ординалами. Сам фон Нейман видит данное ограничение и акцентирует на нём внимание: поскольку у нас нет гарантии, что все модели имеют одну морфологию, нисходящие цепочки возможны в каком-то ином отличном от ординального смысле. В контексте современных теорий всё значительно упрощается, поскольку постулируемое бесконечное множество содержит пустое множество, которое и обеспечивает фундирование. Дело в том, что доказательство невозможности нисходящих цепочек в современных теориях опираются на ординальную морфологию, а в случае индуктивно-дедуктивной морфологии подобное утверждение просто неверно, так как $\dot{\omega}$ является контрпримером. Более того, аксиома фундирования является необходимой для индексной алгебры, она гарантирует, что для всякого множества в индуктивно-дедуктивной морфологии $\bigcup x \neq x$. Действительно, если бы какой-то индексал $\alpha$ не обладал таким свойством, то он бы не могу существовать по аксиоме фундирования. Но аксиома выбора обеспечивает индексальное ранжирование всех множеств, точно также как и ординальное ранжирование, а значит данный вывод справедлив и для всех множеств. Мы наблюдаем потрясающее явление, которое мы отказывались видеть долгое время, поведение всех функций алгебры зависит от морфологии отношения. Кардинальные, ординальные и индексальные функции - это различные, не сводимые друг к другу объекты. Необходимо также сказать о том, что сама по себе индуктивно-дедуктивная морфология не обеспечивает Евклидовости описания, то есть того, что целое больше своих частей. Данная морфология является такой же структурной аннотацией как и индуктивная морфология вполне упорядочения, в частности, множество чётных индексных чисел в индуктивно-дедуктивной морфологии подобно множеству всех индексных чисел, что делает эти множество бесконечными по Дедекинду. В то время, как свойство Евклида необходимо постулировать отдельно.

Несложно также видеть, что такое расслоение отношения принадлежности даёт консервативное расширение теории, поскольку все классические формулы с индуктивным отношением принадлежности истинны в новой теории (можно рассматривать тождество $=$ как сокращение для экстенсиональности определённого отношения принадлежности, фактически отношение эквивалентности для соответствующей морфологии). Стоит также отметить, что символьное разделение различных морфологий - это всего лишь небольшая коррекция выразительного аппарата теории, которая необходима исключительно нам для осознания отсутствия противоречий между различными морфологиями, теория остаётся непротиворечивой даже без такого расслоения. Дело в том, что сама теория ZFC имеет существенный архитектурный недостаток: её аппарат выразимости не является общезначным для всех моделей. Этот изъян особенно ярко проявляется когда за отношением принадлежности подразумевается некоторое более фундаментальная сущность, например, класс или функция как в NBG, что позволяет нам сформулировать альтернативные аксиомы бесконечности без расслоения выразительного аппарата. Его аналогом становятся простые определения морфологических фрагментов внутри теории. Так в силу стандартной аксиомы бесконечности существуют функции (классы), которые связывают исключительно объекты индуктивной морфологии, что даёт классический фрагмент теории. Все теоремы связанное с ними мы называем индуктивными, что придаёт смысл индуктивному использованию отношения принадлежности. Но помимо этого существуют альтернативные функции, которые видят дедуктивные объекты, связанное с ними использование отношения принадлежности мы называем индуктивно-дедуктивным. Отношение принадлежности в таком случае является синтаксическим сахаром, выражающим некоторый функциональный паттерн как для одного так и для другого типа функций. Данное наблюдение позволяет нам определить дедуктивную компоненту отношения принадлежности внутри классической аксиоматики ZFC и сформулировать альтернативную аксиому бесконечности для компонент единого отношения. Например, мы можем воспользоваться тем фактом, что не во всех моделях $\bigcup \omega = \omega$ и сформулировать аксиому бесконечности с альтернативной морфологией непосредственно в языке теорий ZFC и NBG:
\begin{equation}
\begin{gathered}
\exists I \; \bigl[\emptyset \in I \wedge \forall x, y, z, u \; \bigl( x \in I \wedge y = x \cup \{ x \} \wedge z = \bigcup\nolimits^y I \wedge y + u = z \to\\ 
y \in I \wedge z \in I \wedge (\emptyset \in u \wedge \forall v \; (v \in u \to v \cup \{ v \} \in u)) \bigr)\bigr]
\label{axiom:indexal-infinity-no-delamination}
\end{gathered}
\end{equation}

Где $\bigcup\nolimits^y I$ - $y$-кратное взятие объединения всех элементов $\underbrace{\bigcup \bigcup \dots \bigcup}_{y} I$, а $y + u$ - $u$-кратное объединение с собственным синглтоном, оба выражения обеспечиваются схемой аксиом преобразования. Иными словами, существует множество замкнутое индуктивно-дедуктивно так, что разность равноотстоящих от краёв элементов имеет описание в индуктивной морфологии. Данная аксиома полностью сформулирована в языке теорий ZFC и NBG и не вступает в противоречие с какими-либо утверждениями связанными с индуктивной аксиомой бесконечности, поскольку они задают параллельные морфологии преобразуемые друг в друга посредством аксиомы выбора. Благодаря рамке референтных позиций, данное утверждение воспринимается как элементарное, в то время без неё оно с высокой долей вероятности подверглось бы идеологической критике. Таким образом, результаты на которых основаны современные убеждения, теряют общезначность в теории и становятся применимыми исключительно к индуктивной компоненте отношения принадлежности. Формально в данном случае мы говорим, что такое расширение теории {\it консервативно относительно} индуктивного семантического фрагмента теории, то есть в этом фрагменте данное расширение не позволяет доказать новых теорем, зато в других фрагментах старый язык приобретает новый смысл. Более строго: расширение называется консервативным относительно семантического фрагмента теории когда существует расслоение языка изолирующее данный фрагмент и допускающее формулировку изоморфного консервативного расширения. Такое отсутствие общезначности является следствием выбранного нами способа введения бесконечной морфологии, мы просто постулируем, что она существует, но не запрещаем существование альтернативных морфологий, как предлагал сделать Френкель. В сущности данное наблюдение весьма тривиально: из аксиомы существования красных множеств, мы не можем сделать вывод, что все множества красные. Мы имеем дело с классической логической ошибкой поспешного обобщения или в терминах когнитивной науки - сверхобобщения, когда на основании одного свидетельства - аксиомы бесконечности, мы предполагаем общее правило - морфология всех множеств индуктивна и далее применяя FOL совершаем формальную ошибку ложной дихотомии, что полностью разделяет синтаксис и семантику теории. Более того, мы уже де-факто принимаем интерпретацию расслоения ZFC и NBG когда имеем дело с математическим анализом, который определяет континуум - морфологию бесконечного отличную от морфологии вполне упорядочения. Две этих морфологии фундаментально не сводимы друг к другу, а аксиома выбора позволяет лишь преобразовать объект из одной морфологии в другую получив совершенно иной объект. Несложно представить порядок на множестве континуум как морфологию отношения принадлежности: любое вещественное число содержит все предыдущие, и соответствующим образом изменить аксиому непрерывности. Тогда мы будем говорить не о расслоении, а о фрагментации теории с сохранением общего языка. 

Использованный нами подход построения теории множеств от наиболее детализированного описания индексной алгебры к наименее детализированной кардинальной, полностью противоположен подходу Кантора. Его преимущество состоит в том, что он позволяет продемонстрировать факт, который долгое время оставался в тени: функции - это внутренние сущности алгебры, объекты конкретной аннотации. Давайте рассмотрим как некоторые множества пар $f$ и $g$ могут быть аннотированы различными алгебрами. Когда мы работаем в индексной алгебре мы используем внешние функции $f_{ind}$ и $g_{ind}$ с самой богатой морфологией (индуктивной и дедуктивной). Для индексного множества $\dot{\omega}$ функция $f$ реализующая отношение следования переводит данное множество во множество $(\dot{\omega} + 1) \setminus \{ \emptyset \}$: $f_{ind}(x) = x + 1, f_{ind} : \dot{\omega} \to (\dot{\omega} + 1) \setminus \{ \emptyset \}$, в свою очередь само множество $\dot{\omega}$ является областью значений функции $g_{ind}(x) = x + 1, g_{ind} : \dot{\omega} - 1 \to \dot{\omega}$. Обе функции представляют собой объект индексной алгебры - множества пар чисел индексной аннотации множеств: $f_{ind} = \{ (0, 1), (1, 2), \dots, (\dot{\omega} - 2, \dot{\omega} - 1), (\dot{\omega} - 1, \dot{\omega}) \}$, $g_{ind} = \{ (0, 1), (1, 2), \dots, (\dot{\omega} - 3, \dot{\omega} - 2), (\dot{\omega} - 2, \dot{\omega} - 1) \}$. Даже в индексной алгебре наша аннотация не позволяет увидеть реальность скрытую за сингулярностью $\dots$, расположенное между областями регулярности. Когда мы переходим в порядковую алгебру, мы исключаем из описания дедуктивную часть морфологии множеств, но сами множества как объекты остаются не низменными. Если мы обратимся к аннотации наших множеств пар $f$ и $g$ в порядковой алгебре, то увидим следующую картину: $f_{ordinal} = \{ (0, 1), (1, 2), \dots, \{ \{ \dot{\omega} \} \}\}$, $g_{ord} = \{ (0, 1), (1, 2), \dots \}$. В ZFC наша схема аксиом преобразования гарантирует нам существование объекта $g_{ord} : \omega \to \omega \setminus \{ \emptyset \}$, который в этой алгебре не отличим, например, от аннотации объектов индексной алгебры вида $h_{ind} : \dot{\omega} - n \to (\dot{\omega} - n + 1) \setminus \{ \emptyset \}$, где $n$ - конечное число. При этом объект $f_{ord} : \omega \to (\omega + 1) \setminus \{ \emptyset \}$ является абсолютно корректным объектом нашей порядковой аннотации (существование которого следует из схемы аксиом преобразования и аксиомы объединения). Но в классических рассуждениях мы из существования объекта $g_{ord}$ делаем вывод о рефлексивности множества $\omega$, несмотря на существование объекта $f_{ord}$ (этот вывод особенно шокирует в контексте NBG, где отсутствует какое-либо указание на выделенность $g_{ord}$). Иными словами, наш собственный формализм оказывается умнее нас. Даже в рамках ZFCL у нас есть объект $g_{ord}$, характеризующий отношение между объектами $\omega$ и $\omega \setminus \{ \emptyset \}$ с внутренней перспективы, и есть объект $f_{ord}$, характеризующий отношение между объектами $\omega$ и $(\omega + 1) \setminus \{ \emptyset \}$ с внешней перспективы. Именно это Кантор реализовал с помощью левых и правых алгебраических операций, осознание чего мы потеряли в борьбе с парадоксами самореференции. Таким образом, о рефлексивности в современных теориях множеств уместно говорить только по отношению к классу определённых объектов, описывающих взгляд с внутренней позиции референта. Когда мы переходим к кардинальной алгебре мы исключаем из описания индуктивную часть морфологии бесконечных множеств. Здесь оба объекта $f_{card}$ и $g_{card}$ неотличимы:  $f_{card}, g_{card} : \aleph_0 \to \aleph_0 \setminus \{ \emptyset \}$. Мы фиксируем внутреннее описание, где пары $(\dot{\omega} - 2, \dot{\omega} - 1)$ и $(\dot{\omega} - 1, \dot{\omega})$ приобретают вид $(\aleph_0, \aleph_0)$, но объект $\aleph_0$ не является внутренним объектом счётного множества, $aleph_0$ не может ни превосходить сам себя $aleph_0 < aleph_0$, ни быть частью самого себя $aleph_0 \in aleph_0$. Именно так возникает классическая интуиция потенциальной бесконечности, внутри кардиналов мы не можем видеть бесконечной морфологии, мы видим лишь индуктивную область регулярности сформированную исключительно конечными значениями. Если мы теперь аннотируем во внутренней кардинальной алгебре биекцию $\varphi_{ord} : 2^{\omega} \to \mathcal{P}(\omega)$ между множеством $2^{\omega}$ и множеством подмножеств $\mathcal{P}(\omega)$, то получим классический результат Кантора: $\varphi_{card} : \aleph_0 \to 2^\aleph_0 = \aleph_0$, множество алгебраических чисел соответствующих многочленам с коэффициентами $0$ и $1$ счётно. Такие пары как $(2^{\dot{\omega}} - 1, \dot{\omega})$ или $(2^{\dot{\omega} - 1}, \{ 0, 1, 2, \dots, \dot{\omega} - 1 \})$ во внутреннем кардинальном представлении множества $2^{\omega}$ имеют вид: $(\aleph_0, \dot{\omega})$, $(\aleph_0, \{ 0, 1, 2, \dots, \dot{\omega} - 1 \})$, точно также как и другие пары с бесконечными комбинациями элементов $\omega$ в значении, например, $(\aleph_0, \{ 0, 2, 4, \dots \})$ или $(\aleph_0, \{ 1, 3, 5, \dots \})$. Поскольку объект $\aleph_0$ не является внутренним для первой трансфинитной мощности, эти пары в кардинальной алгебре аннотируются как синглтоны, чего уже достаточно, чтобы исключить их из внутреннего функциональное описания $\varphi_{card}$, оставляя от трансфинитной нумерации лишь конечную часть, не говоря уже о том, что значения таких пар, также не могут быть подобными внутренними объектами. Именно здесь конечность носителя возникает естественным образом, вследствие отсутствия морфологии бесконечных объектов, которая могла бы быть вписана во внешнее пространство для сравнения, сама топология такого кардинального пространства не позволяет сравнивать бесконечные объекты. Поэтому у Кантора порядковое тип $\omega^{\omega}$ оказывается счётным, биекция между $\omega$ и $\omega^{\omega}$ упирается в $\aleph_0$ и не видит никакой трансфинитной морфологии, это в чистом виде классическая потенциальная бесконечность, обеспечивающая счётность алгебраических чисел. 

Как мы неоднократно отмечали NBG находится в более выгодном положении, поскольку уже постулирует необходимую смену позиций референта. Классы являются чисто внутренним описанием, наиболее близким к кардинальной алгебре Кантора, в то время, как множества представляют собой объекты вписанные в общее унифицированное пространство с заданными бесконечными морфологиями. Мы видим, что теорема Кантора не доказуема для классов, как это и должно быть в чисто внутреннем представлении. Но даже здесь мы натыкаемся на фундаментальные ограничения логики. Нам необходимо узаконить исключительно иерархическое построение функций, сделав статус циклических формул неопределённым. Например, уравнение $A = \{ \emptyset, \{ x | x \in A \wedge x \notin x \} \}$ не может иметь решением объект внутри нашей теории. Иными словами, такой объект принадлежит пространству, которое наша теория не в состоянии описывать. Сам фон Нейман обращает внимание на то, что аксиома ограничения размера - это всего лишь один из возможных критериев универсальной аннотации валидных с точки зрения теории объектов. Однако, как мы видим на примере цикличности система разделения на функции и аргументы сама справляется с парадоксами, делая вывод о возможности или невозможности тех или иных конструкций. Что открывает путь к описанию гораздо более богатых и разнообразных структур, чем классические множества. Проблема с аксиомами цензуры состоит в том, что они с течением времени догматизируются и становятся границами, выход за которые воспринимается как преступление. Это следствие языковой игры, если для древних греков аксиомы были утверждениями допускающими непосредственную верификацию в опыте, то в современном математическом языке они становятся истинами не требующими доказательств, что кардинально меняет их статус. Мы не способны верифицировать аксиому ограничения размера в непосредственном опыте, поэтому корректней было бы называть подобные утверждения гипотезами. Это нисколько не отменяет того, что гипотеза ограничения размера делает выводимыми схемы аксиом выделения, преобразования и даже аксиому выбора. Которые в свою очередь более подходят на роль аксиом, поскольку допускают непосредственную верификацию на объектах, с которыми мы сталкиваемся в реальности. На взгляд авторов, именно идеологизация мейнстрима препятствует тому, что мы имеем настолько мощный инструмент, но не в состоянии им полноценно пользоваться.

В частности, в NBG мы можем ввести кардинальность как одну из морфологий бесконечности, даже не прибегая к концепции классов эквивалентности. В самом деле, мы обнаружили, что кардинальные функции - это такие функции, которые просто не видят внешней морфологии ординалов, то есть в контексте кардинальной структуры, бесконечных объектов просто не существует, что обеспечивает абсолютную индуктивную замкнутость. Таким образом, аксиома стеснения Кантора даёт достаточный фундамент для погружения классов в такую морфологию бесконечности.
\begin{equation}
\begin{gathered}
\exists I \neq \emptyset \; \forall X \; \bigl(X \sim I \wedge \forall x \; (x \in X \to \exists i \subseteq I \; (x \sim i)) \bigr) \to \bigcup X \sim I
\label{axiom:cardinal-infinity}
\end{gathered}
\end{equation}

Фактически, мы просто постулируем, что существует не пустое комбинаторно замкнутое множество, которое в силу своих свойств не может быть конечным. При этом по аксиоме сводимости оно не является собственно классом, а значит представляет собой некоторую морфологию бесконечности. Мы снова явно не связываем данную морфологию бесконечности с ординальной и индексальной. На исходном уровне описания Джона фон Неймана мы определяем лишь тип функций, которые видят исключительно внутренне замкнутые уровни аргументов, не описывая при этом аргументов, которые связаны с внешней ординальной или индексальной морфологией. Данная формулировка в контексте современной интерпретации без представления о позициях референта кажется провокационной, но если мы прибегаем к расслоению отношения принадлежности и вводим новый символ для кардинальной компоненты, то всё становится на свои места - мы получаем консервативное расширение, которое реализуется в моделях теории. 

Устройство этой морфологии помогает раскрыть аксиома выбора, позволяющая вполне упорядочить любое кардинальное множество в некоторый порядковый тип. Как показал Кантор, как только появляется понятие вполне упорядоченности арсенал алгебраических инструментов обогащается пределами и бесконечными рядами, что позволяет ввести унифицированное координатное пространство с упорядоченной морфологией. Введя различие левых (внутренних) и правых (внешних) операций, он продемонстрировал гениальный способ разделить уровни самовоспроизводства аддитивной структуры, за счёт того, что пределы и ряды оказываются связанными исключительно с внешней морфологией правых операций. Его подход эквивалентен соответствующей интерпретации аксиомы выбора, подразумевающей воспроизводимость, а значит способной индуцировать морфологию вполне упорядоченности заданную классической аксиомой бесконечности на любое множество. В таком пространстве по теореме Кантора $\omega < \omega^{\omega}$, но если мы рассмотрим данное соотношение в кардинальной морфологии, то слева мы видим счётно множество, а справа счётное множество счётных множеств, то есть такое же счётное множество: $\aleph_0 = \aleph_0^{\aleph_0}$. Иными словами, индуктивный обход матрицы, который в ординальной морфологии давал строго больший порядковый тип, в кардинальной морфологии оказывается эквивалентным исходному. Поэтому теорема Кантора не может быть доказана для кардинальной арифметики: когда мы пытаемся сформулировать диагональный аргумент, нам необходимо зафиксировать количество столбцов, мы можем это сделать в ординальной алгебре, взяв например $\omega$, однако в кардинальной алгебре столбцы $\omega + 1, \omega + 2, \dots, \omega^{\omega}, \dots$ - счётны, у нас принципиально отсутствует морфологический критерий ограничивающий возможность продолжения диагонального построения, вследствие внешней импредикативности объектов. Как и в конечной области, оно продолжается неограниченно не давая никакого результата, из конечной области оно переходит в счётно бесконечную, затем в континуально бесконечную и сквозь всю иерархию кардиналов устремляется в пространство собственно классов. 

Мы видим лишь как данное построение упирается в сингулярность кардинальных уровней и снова появляется на следующем ступеньке лестницы. Но в кардинальной алгебре у нас нет никаких морфологических ограничений, таких как ранги, которые могли бы зафиксировать бесконечный объект с внутренней рефлексивностью не распространяющейся на всю область, у нас отсутствует линия радела между кардинальными уровнями, которая позволяла бы идентифицировать целостное построение. В индексной алгебре мы отчётливо видим, что для каждого фиксированного множества теорема Кантора выполняется, но кардинальная алгебра не способна уловить эту фиксацию, она нужна для другого, она обеспечивает замкнутость индуктивной структурны расположенной между кардинальными областями сингулярности. С формальной точки зрения, поскольку у нас есть аксиома регулярности, никакое внешнее диагональное построение невозможно в кардинальной морфологии. В самом деле, если мы попытаемся построить элемент отличный от всех с помощью внешнего сложения $\cdot + 1$ и объединить его со всеми исходными элементами, то мы получим $\aleph_0 + 1 = \aleph_0$, откуда должно следовать $\aleph_0 \in \aleph_0$, что противоречит аксиоме регулярности. Таким образом, аксиома регулярности работает как стратификация не позволяя сформулировать идею самопринадлежности внутри фрагмента с кардинальной морфологией. Внешние функциональные отношения оказываются заблокированными в кардинальной алгебре, но не по нашей прихоти, а по правилу стратификации. В действительности, нам не требуется аксиома регулярности, она помогает лишь увидеть семантику решения синтаксического уравнения, но как только мы её видим, мы можем отказаться от данной аксиомы цензуры. В таком случае, отсутствие внешней компоненты в кардинальной алгебре следует непосредственно из парадокса Рассела, предположив, что внешняя импредикативность допускает предикативную интерпретацию отношения принадлежности мы приходим к противоречию. Теорема Кантора оказывается действительно не доказуема в ZFC и NBG без предположения единственности внутреннего функционального описания (отсутствия кардинальных функций) и единственности индуктивной морфологии (отсутствие кардинальных структур). При этом данная теорема остаётся абсолютно корректной если мы оперируем исключительно индуктивными функциями. 

В аксиоматике ZFC для того, чтобы эмулировать естественную функциональную природу отношения принадлежности, как и прежде, можно воспользоваться расслоением, введя отношение кардинальной принадлежности. Здесь аксиома регулярности, снова помогает избежать циклов вроде $\aleph_0 \in \aleph_0$, одно кардинальное множество $\aleph_0$ представляет всю счётную область ординалов и индексалов, с их внешней морфологией включающую множества подмножеств. Интересно, что когда мы явно вводим кардинальную морфологию в теорию множеств, мы естественным образом обнаруживаем связь отсутствия неподвижных точек в кардинальной, ординальной и индексальной структурах, обеспеченную аксиомой регулярности, с ограничением стратификации в новых основаниях Куайна. Таким образом, стратификация получает морфологическое обоснование, она не позволяет коллапсировать рефлексивным замкнутым алгебраическим структурам за счёт исключения таких отношений как $\aleph_0 \in \aleph_0^{\aleph_0} = \aleph_0$. Поэтому в новых основания Куайна кардинальные объекты, для которых недоказуема теорема Кантора и которые позволяют замкнуть универсум, непротиворечиво сосуществуют с ординалами и индексалами. Концептуально задумка Кантора проста и элегантна, всё конечное можно развернуть последовательно, связав все такие объекты в единый класс. Из конечного мы получаем бесконечное, выделив например числа $0, 1, 2, \dots$ или формулы арифметики $0 < 1, 2 + 2 = 4, \dots$ или их подмножества и комбинаторные объекты. Используя сами такие бесконечные цепочки можно порождать следующее поколение бесконечных объектов и так далее. Всё бесконечное, что мы можем получить таким способом, включая производные любых уровней, таких как возведение в степень $x, x^2, x^3, \dots$, тетрация $x, x \uparrow \uparrow 2,  \uparrow \uparrow 3, \dots$, пентация и прочее. Фактически Кантор предлагает замкнуть область по отношению к понятию предела как следующий класс всего конечного и бесконечного. Выделив в этом классе совокупность объектов расширяющую концепцию натурального ряда в области бесконечности, Кантор показывает, что такой класс не может содержать объекта описывающего количество его собственных элементов, что требует бесконечности нового уровня и предельного аппарата нового уровня, что в свою очередь может быть продолжено на новые и новые уровни, которые и представляют собой кардинальные числа. Фактически аддитивная структура воспроизводит себя даже на кардинальном уровне. При этом не возникает морфологии бесконечного, благодаря чему внутренние отношения обеспечивают тотальную замкнутость алгебраической структуры: подмножества натуральных чисел $\{\}, \{0\}, \{1\}, \{0, 1\}, \dots$ естественным образом разворачиваются в последовательность, равно как и все алгебраические числа и все комбинаторные комбинации алгебраических операций и пределов первого уровня. Несложно видеть, что в ZFCL и NBGL аксиома кардинальной бесконечности оказывается излишней, поскольку ординалы уже обладают данным свойством, в то время как в ZFCS и NBGS мы с её помощью восстанавливаем всю полноту структурного описания бесконечных объектов и обеспечиваем полноценную возможность переключения перспектив. 

Концепция импредикативности или рефлексивности сама по себе не является источником противоречий или контринтуитивных результатов. Единственный по-настоящему рефлексивный объект - это ничто, пустота. Только этот объект обладает свойством тождественности любой своей части. Что математически можно выразить как $0 - 0 = 0$ или $\emptyset \setminus \emptyset = \emptyset$. Откуда можно сделать вывод, что рефлексивность - это воплощение особенности нашего мышления работать в условиях недостатка информации. Когда мы рассуждаем о непрерывности, мы выражаем такой недостаток информации обо всех объектах задающих отношения, в которых непрерывное делимо, в свойстве подобия непрерывного своей части. Подобие, таким образом, - это контролируемая утрата части контекстной информации. В геометрии - это утрата части метрической информации, связывающей внутреннюю морфологию с контекстом внешнего пространства. Поэтому когда Кантор и Дедекинд заимствуют свойство подобия у непрерывного, это неизбежно связано с потерей контекстной информации. В случае вполне упорядоченности, внутренняя морфология характеризуется лишь одним типом точек, имеющим выделенный смысл - это начальные точки, из которых воссоздаётся вся структура. Однако, речь идёт о контролируемой утрате информации, допускающей её восстановление. Именно поэтому Кантор и Дедекинд используют понятия типов - классов эквивалентности, фиксирующих структуру объекта, а не его сущность. Когда мы восстанавливаем тип до множества, вписываем его в общее пространства с другими объектами, оно становится целостным объектов, для которого работает аксиома Евклида - целое больше части. Реификация типов до ординалов приводит к тому, что мы перестаём различать карту и территорию, что особенно наглядно демонстрирует парадокс Банаха - Тарского. Когда утрата информации контролируется подобие становится мощным инструментом, с помощью которого мы можем показать, что $0.99(9) = 1$, или решать такие уравнения как $x^{x^{x^{\dots}}} = n$. Но как только мы теряем контроль, отвлекаемся от критериев стабилизации, рефлексивность может приводить к противоречиям, как в случае, когда мы позволили внутренней рефлексивной составляющей ординалов доминировать над структурным ограничением теоремы Кантора. Именно поэтому в аксиоматиках теории множеств стоит явно формулировать, что рефлексивность возникает из объективизации внутреннего индуктивного описания. Частичную рефлексивность связанную с внутренней морфологией объектов мы называем {\it подобием}, возвращаясь к исходной интуиции Кантора и Дедекинда. Таким образом, если мы хотим полноценную рефлексивность в смысле современной интерпретации, то её необходимо постулировать отдельной аксиомой, поскольку как было неоднократно показано она не выводима из современных аксиоматик. Более того, полноценная рефлексивность онтологически является свойством пустоты, именно поэтому с рефлексивностью множеств связаны парадоксы вида $1 = 2$, это проявление природы онтологической категории ничто. Когда мы наделяем бесконечное полноценной рефлексивностью мы в буквальном смысле обращаем её в ничто. В то же время, подобие как контролируемая рефлексивность выводимо и нуждается лишь в корректном определении класса объектов теории, которые мы считаем выражением отношений внутренней морфологии. Данный факт хорошо иллюстрирует доказанное Кантором подобие дискретного описания одномерного континуума и многомерного, здесь внутреннее отношение $1 = 2$ свидетельствует о частичной утрате информации при переходе к дискретному описанию непрерывности, в то время как двумерное пространство структурно богаче одномерного и проявляет свойства невоспроизводимые в одномерном.

В работе ''К изучению о многообразиях'' \cite{cantor-set-study} Кантор строит взаимно однозначное соответствие между многомерным пространством и непрерывной числовой прямой и, в частности, упоминает теорему, полученную при помощи некоторого метода индукции, утверждающую, что число классов линейных точечных многообразий является конечным и даже равным двум, первый класс можно привести к виду функций от натурального аргумента, второй класс - к виду функций от вещественного аргумента, этим двум классам соответствуют лишь две мощности линейных многообразий. Этой теореме впрочем не суждено было обрести доказательство, но предстояло возглавить список кардинальных проблем математики, сформулированный Давидом Гильбертом на рубеже веков, где она получила название континуум-гипотезы \cite{hilbert-problems}. Стоит отметить, что утверждение Кантора было существенно упрощено адаптацией мейнстрима. На взгляд авторов, в данном контексте Кантор имеет ввиду нечто более глубокое - принципиальную несводимость дискретного и непрерывного друг к другу, а также отсутствие иных категорий множественности в нашем восприятии. Дело в том, что Кантор мыслил континуум как совершенное множество, которое тождественно всем своим производным. Он начинал с множества рациональных чисел и строил предельные производные сквозь всю иерархию трансфинитных чисел второго числового класса. Таким образом, совершенное множество получалось как $S = P^(\Omega)$ - производная минимального числа третьего числового класса (первой несчётной мощности) \cite{cantor-perfect-set}. Данное наблюдение о природе континуума уже демонстрирует принципиальное различие кардинальности для Кантора и тем, что в последствии закрепилось в мейнстриме. Континуум Кантора не ограничен мощностью множества степени как в упрощении стандартной интерпретации, диагональным аргументом Кантор лишь демонстрирует, что конфигурация континуума как всюду плотного непрерывного множества связана с предикативностью как минимум экспоненциального уровня, в то время как аддитивная структура остаётся импредикативной. Но его строгие формулировки остаются не востребованы в математике на протяжении десятилетий, как и его фундаментальное открытие ключевой роли конфигурации предикативности и импредикативности в анализе, непосредственно связанное с позициями референта. Однако, таким образом определённое совершенство является совершенством только по отношению к пределам первого порядка, то есть определяемых счётными последовательностями. Фактически он предлагает принципиально иную концепцию континуума, континуум - это замыкание предельного аппарата некоторой алгебры, что гораздо точнее описывает фрактальность, в частности, свойство аддитивной структуры воспроизводить себя на различных уровнях.  Кантор вплотную подошёл к сюрреальным числам и совершенным множествам любой мощности, что делало структуру непрерывности гораздо более богатой. Так или иначе, с современной точки зрения континуум-гипотеза формулируется как отсутствие промежуточных мощностей между $\aleph_0$ и $2^{\aleph_0}$. В контексте ZFCL и NBGL континуум-гипотеза остаётся доказано независимой от аксиоматик. Даже если мы легализуем объекты описывающие внешние функциональные отношения, воссоединив их с натуральными операциями, и введём полноценную операцию разности, мы получим дедукцию симметричную индукции, которая также не сможет сформировать аксиоматических множеств таких как степень в контексте данных теорий. В то же время, если мы вводим внутренние функции кардиналов, и хотим, чтобы теорема Кантора была применима к кардинальной алгебре, то наличие промежуточных мощностей становится необходимым, поскольку если бы в кардинальной морфологии мощность континуума представлялась как $2^{\aleph_0} = \{0, 1, 2, \dots, \aleph_0\}$, то она допускала бы соответствие с $\aleph_0$. Однако, здесь она не представляет существенного интереса, поскольку является проблемой абстрактной синтетической алгебраической структуры, с относительным понятием мощности. В контексте ZFCS и NBGS континуум-гипотеза не имеет смысла, поскольку континуум в понимании мейнстрима, как замыкание множества рациональных чисел, здесь счётен. Любопытно, что самая первая формулировка оказывается самой меткой, допускающей всего два фундаментальных типа описания - дискретный и непрерывный, и возвращает нас к античному делению на числа и величины, на которой, в частности, строит своё изложение Аристотель. В такой формулировке данная проблема представляется весьма интересной и интригующей, возможно ли математически доказать или опровергнуть то, что эти два способа описания реальности необходимы и достаточны, и что таким образом наше мышление достаточно универсально для описания законов вселенной. Цензурированные теории очень любопытное направление для исследований, поскольку их развитие позволит нам выработать формальные критерии необоснованного исключения объектов из описания, что даст нам возможность по-новому взглянуть на существующие математические теории и проанализировать не содержат ли они признаков цензуры.

В частности, здесь следует упомянут парадокс Скулема, ясно, что он остаётся в равной степени применим как ZFCL и NBGL, так и к ZFCS и NBGS. Действительно, смысл теоремы Лёвенгейма-Скулема можно сформулировать следующим образом, если множество формул выполнимо в какой-либо области, то оно непротиворечиво, а любое счётное множество непротиворечивых формул языка первого порядка выполнимо в счётной области, что должно было противоречить любой теории множеств, в которой существуют несчётные множества \cite{fraenkel-skolem}. Данный парадокс иллюстрирует нечто фундаментальное, применимое к любым нашим описаниям действительности. Всякая конечная система правил определяет некоторую замкнутую песочницу со своим внутренним миром и своими внутренними структурами. Однако, наше мышление таково, что мы всегда можем выйти из этого мира во внешнее пространство и посмотреть на всю систему со стороны. Если система допускает неограниченное индуктивное построение, то внешне мы неизбежно сталкиваемся с определённым произволом восприятия её целостности. У нас нет единого для всех людей ментального пространство, которое мы могли бы использовать для того, чтобы гарантировать воспроизводимость внешнего взгляда на данную систему. А значит бесконечность, вероятно как и любая абстракция, является относительным понятием, фиксирующим структуру, а не конкретную сущность. Таким образом, теорема Лёвенгейма-Скулема - это фундаментальная теорема мышления строго разделяющая карту и территорию, наши теории всегда оперируют картами, а не территориями. Несмотря, на то что данная теорема не доказывается в логике второго порядка, из-за иерархии мета-теорий Тарского она становится универсальной, вследствие которой у нас отсутствует универсальная математическая теория и всякую совокупность теорий можно вписать в контекст более сильной первопорядковой мета-теории. Поэтому абсолютная категоричность невозможна и даже арифметика второго порядка может иметь модели разной мощности, за счёт некатегоричности мета-теории, хотя иногда некатегоричность первопорядковой арифметики связывает с теоремой Лёвенгейма-Скулема, что не совсем корректно \cite{fraenkel-categoricity}. Иными словами, бесконечность следует рассматривать как структурное свойство, а не как экзистенциальную характеристику. А следовательно, не может существовать и универсальной морфологии бесконечного, любая фрактальная структура воспроизводящая себя на различных уровнях, описывается данным свойством. И наоборот, всякий объект обладающий свойством бесконечного может описываться различными морфологиями. Подобие или изоморфизм - это внутренние понятия структуры, которые не могут определять внешние свойства описываемых объектов. Именно теорема Лёвенгейма-Скулема помогает избежать противоречий в аксиоматиках ZFCL и NBGL, смешение перспектив не приводит к коллапсу теории поскольку в рамках аксиоматики остаётся формальная возможность разделить индуктивные множества и аксиоматические множества. Вследствие этого индуктивные множества могут достигать степени на одном уровне бесконечности $2^{\omega}$, в то время, как аксиоматические способны сразу перепрыгивать через эту индуктивную иерархию и формировать $2^{\aleph_0}$. Подобие всех уровней бесконечности не позволяет нам прийти к коллизии внутри теории вследствие данного факта, пока теория остаётся достаточно слабой. Но в то же самое время, такая интерпретация очень тяжело согласуется с альтернативными морфологиями, и практически несовместима с целостностью по Евклиду, что способствует её идеологизации и маскировке единственной уязвимости - устойчивой условной непротиворечивости. В то же самое время, совокупность данных фактов препятствует существенному усилению теорий использующих такую интерпретацию.

Описанные в данной статье подходы невероятно обогащают математическую логику, где диагональный метод получил, пожалуй, наибольшее распространение и используется, например, в таких важных утверждениях как неразрешимость проблемы остановки \cite{computable-halting}, теорема Клини о неподвижной точке \cite{computable-kleene}, теорема Гёделя о неполноте, теорема Тарского о невыразимости истины \cite{computable-tarski-godel} и прочих. Однако, диагонализация здесь принципиально отличается от диагонализации Кантора. Судя по всему, первым на это фундаментальное различие обращает внимание Фрэнк Рамсей разделивший парадоксы на логические (синтаксические) и эпистемологические (семантические) \cite{ramsey-paradoxes}. К первой группе относятся парадоксы Рассела \cite{russel-paradox}, Кантора \cite{cantor-paradox}, Бурали-Форти \cite{burali-forti}, которые имеют непосредственное отношение к диагональному методу Кантора, в то время как, во вторую группу входят парадоксы Лжеца, Ришара \cite{richard-paradox}, Греллинга-Нельсона \cite{grelling-nelson} и Берри \cite{berry-paradox}, которые можно связать с диагонализацией предложенной Гёделем. Первая группа была обнаружена для бесконечных объектов, и данная иллюзорная, как мы обнаружили, связь закрепилась и помешала обнаружить истинную причину конфликта. Вторая группа считалась смысловой ошибкой и часто игнорировалась, из-за чего современные теории остались к ней уязвимыми. Как не сложно видеть обе группы имеют непосредственное отношение к позиции референта, и с этой точки зрения разница между ними становится более отчётливой. Действительно, логические парадоксы явно пытаются рассматривать один и тот же объект одновременно с внутренней и внешней позиции. В парадоксе Рассела множество внутренне конструируется через свойство, которым оно внешне должно обладать как целостный объект. В парадоксе Кантора множество всех подмножеств $\mathcal{P}(X)$ некоторого множества $X$ определено внешне по отношению к $X$ (содержит $X$ как элемент), что несовместимо с попыткой вписать его как подмножество $X$ в общее пространство с упорядоченной морфологией. В парадоксе Бурали-Форти мы внутренне рассматриваем класс ординалов как контейнер, но внешне видим в нём обычный ординал. В то время как эпистемологические парадоксы, используют стандартное расслоение на внутреннюю и внешнюю позиции референта в языке. В парадоксе Лжеца высказывание внутренне (синтаксически) утверждает свою ложность, что внешне (семантически) означает противоположное, синтаксически оно пытается занять внешнюю позицию по отношению к себе, но семантически это происходит во внутренней. Парадоксы Ришара и Берри синтаксически принимают построения как данность, но семантически отрицают её, внутренне они формулируют внешнее по отношению к конструкции утверждение, а внешне оказываются её внутренними элементами и посему контрпримерами. В случае парадокса Греллинга-Нельсона мы синтаксически определяем дихотомию свойств, которые семантически ей противоречат. Иными словами, логические парадоксы смешивают позиции референта уже в синтаксисе, а эпистемологические маскируют это смешение в уровнях встроенных в язык. 

Анри Пуанкаре в своей книге ''О науке'' довольно подробно разбирает парадокс Берри \cite{poincare-berry} и приходит к выводу, что классификации можно разделить на предикативные и импредикативные, те которые не нарушаются введением новых элементов и те которые нарушаются. Первые - это актуально существующие списки, такие как классификация принадлежности военнослужащих или биологическая классификация живых организмов, вторые - это потенциально существующие системы, которые описываются общими словами, но не фиксируются конкретными правилами построения. Последние позволяют определять импредикативные объекты, которые в своём определении ссылаются на саму классификацию. Предикативность - это исключительно внутреннее свойство, которое не может быть использовано для выхода за пределы системы. Он связывает импредикативность с неаккуратным введением бесконечных объектов и приходит практически к тому же выводу, который мы сформулировали в своей статье: при работе с бесконечными объектами необходимо строго следить за предикативностью отношений. Это именно то, что произошло в консенсусе математического сообщества, которое превратило начальный ординал - внутренний объект области, во внешний контур замыкающий алгебраическую структуру, проиллюстрировав ту самую импредикативность, о которой предупреждал Пуанкаре. Он справедливо говорит, о том, что все утверждения о бесконечности это обобщённые утверждения о конечных числах, собственно это доказанные математический факт, воплощённый в теореме Лёвенгейма-Сколема (все объекты счётного универсума конечны). Однако, Пуанкаре использует довольно сложный технический язык, импредикативность - это результат, а не причина феномена. Ему приходится говорить о импредикативной классификации, как такой которая допускает порочные определения, то есть фактически денотировать причину через её результат. Он также связывает предикативность с неизменяемостью, то с достаточностью информации для статичной фиксации классификации. На взгляд авторов, Аристотель в своих работах гораздо нагляднее формулирует то различие о котором, хочет сказать Пуанкаре. Он разделяет две категории: бытие в возможности (потенциальное) и бытие в реальности (актуальное). Бытие в реальности - это внутреннее описание фактического существования, в том числе, того что современные конструктивисты называют потенциально конструируемым, это предикативная непротиворечивая система описания реальности. Бытие в возможности - это выход за пределы физической реальности, взгляд на неё со стороны, описание основанное на недостатке информации и фактического материала. Блок мрамора в возможности это и статуя Афродиты и статуя Апполона и статуя Диониса, но как только статуя Афродиты становится реальностью, пребывание её в форме Апполона или Диониса исключается. Альтернативы в возможности в отличии от реальности не противоречат друг другу. Если нам предъявлен напечатанный список всех фраз определяющий целые числа, длина которых не превышает двухста слов, то никакого противоречия с фразой ''наименьшее число списка, которое не определено менее чем ста словами'' не возникает. Наше мышление корректно интерпретирует недостаток информации о актуально существующем списке и его внешнем контексте, где формулируется фраза. Таким образом, сам по себе недостаток информации не является достаточным источником вариабельности и противоречий, он требует ещё одного условия - существования объекта в возможности, а не в действительности. В возможности контур груди Афродиты не противоречит торсу Апполона или канфару Диониса, это части различных альтернатив, точно также список фраз не рассматривающий импредикативную фразу и список учитывающий её - это различные альтернативы, не вступающие в противоречие из-за онтологического статуса данной категории бытия. Когда Аристотель говорит о существовании бесконечности в возможности, он не отрицает её существование как целостной сущности, он лишь говорит, что она существует в мире абстракций, предсказаний нашего мышления, внешнем по отношению к предикативной реальности. Однако, в средневековой схоластике этот язык утрачивается, категория внешнего потенциального бытия и категория внутреннего фактического бытия сливаются, в результате чего Пуанкаре вынужден переизобретать этот язык, что читается довольно тяжело, и вероятно вследствие этого не находит должного отклика. В онтологии Аристотеля пространство предсказаний содержит альтернативные морфологии бесконечного - кардинальную, порядковую, индексную - точно так же, как глыба мрамора содержит в возможности альтернативные формы. Принятие на веру одной морфологии не делает остальные невозможными. Однако в современной теории множеств мы легитимизируем одну форму и объявляем альтернативы исключёнными самой теорией, не утруждая себя каким-либо доказательством этого. С точки зрения Аристотеля это выглядело бы примитивным культом Афродиты, а мы с таким описанием действительности казались бы ему абсолютными дилетантами. В терминах языковых игр Витгенштейна полезно различать: аксиому как утверждение, которое допускает непосредственную верификацию в опыте, то есть утверждение о реальности, и гипотезу как утверждение, которое фиксирует одну из альтернативных форм, таким образом, являясь утверждением о возможности. Когда мы не делаем такого различия, аксиома бесконечности остаётся аксиомой в возможности, которую мы пытаемся интерпретировать как аксиому реальности, но представление формы Афродиты, не меняют реальности глыбы мрамора, которую описывает наша теория.

Гёдель невероятно проницательно замечает, что самореференция которую считали причиной парадоксов возникает естественным образом в любых достаточно сложных системах \cite{incompleteness-original}. Причина данного явления, вероятно, кроется во фрактальности - способности к собственному воспроизводству на различных уровнях. Так аддитивная группа в арифметике воспроизводит себя в мультипликативной, и далее сквозь все уровни степенных башен. Он показывает, что арифметика может кодировать свои формулы благодаря своей основной теореме - уникальности разложения на простые множители: если каждому символу присвоить уникальный номер, то любую формулу $x_1x_2x_3 \dots x_i$ можно закодировать числом $n = p_1^{x_1}p_2^{x_2}p_3^{x_3} \dots p_i^{x_i}$, где $p_i$ - простые числа. Гедёль строит рекурсивную функцию $diag(x, y)$, которая подставляет число $y$ в формулу с одной переменной, имеющей код $x$. Благодаря чему, ему удаётся доказать в мета-теории красивую диагональную лемму: для любого свойства $\varphi(x)$ существует такое утверждение $G$ с кодом $g$, что в арифметике доказуемо утверждение $G \leftrightarrow \varphi(g)$. Иными словами, самореференция возникает в области конечных чисел и позволяет формулам делать утверждения о своих собственных свойствах. Такая самореференция абсолютно безобидна, выражения ''мой код больше 42'' или ''мой код простое число'' абсолютно корректны и не приводят к каким-либо парадоксам арифметики. Однако, проблемы начинаются когда мы пытаемся заставить формулу говорить о внешних метатеоретических конструкциях, например, кванторах логики или понятии истины, то есть смешать уровни, как это делают семантические парадоксы. Рекурсивная функция $proofCheck(x, y)$ верифицирующая факт того, что рассуждение, имеющее код $x$, доказывает утверждение с кодом $y$ позволяет записать предикат доказуемости $Provable(y) = \exists x \; proofCheck(x, y)$, что приводит к неполноте $G \leftrightarrow \neg Provable(g)$. В свою очередь, Альфред Тарский показал, что предикат истинности $True(x)$ не может существовать, то есть арифметика вообще не может оперировать понятием истины внутри себя. На первый взгляд может показаться, что мы имеем дело со смешением внешней перспективы - формул и внутренней перспективы - чисел, как в парадоксе Рассела. Действительно, арифметика второго порядка видит свой универсум $\mathbb{N}$ и может доказать для него теорему Кантора, то есть защититься от возможности закодировать все свои формулы числами. Однако, коварство подхода Гёделя в том, что он не требует кодирования всех формул, какими бы ни были предикаты истины и доказуемости, они представляют собой фиксированные формулы, конечные или бесконечные, благодаря чему нам достаточно лишь возможности закодировать применение предикатов ко всем числам. Поскольку предикаты статичны связанные с ними компоненты кода также статичны, а значит нам достаточно описать вариации аргументов чтобы получить всю механику применимости (семантику), в частности, мы можем выделить все чётные числа для описания самого предиката, а аргументы кодировать нечётными числами. Или просто ввести константу обозначающую внешний оракул имеющий семантику необходимых понятий. Таким образом, аргумент Гёделя схватывает саму способность внешнего семантического обобщения, уклоняясь от теоремы Кантора, что делает его устойчивым к теориям второго порядка, $\omega$-логикам и даже конструктивизму. Если мы считаем, что на внешнем уровне существует какая-либо абстракция или универсализация индукции вытекающая из $P(0) \wedge P(n) \to P(n + 1)$, то подход Гёделя способен её использовать. 

Единственным способом ускользнуть из этой ловушки оказываются альтернативные бесконечные морфологии. В нестандартном анализе, сюрреальных числах Конвея или гросс-объектах аргумент Гёделя всё-таки даёт сбой. Для того чтобы гарантировать, что система видит все возможные применения предикатов внешней семантики к своим объектам, в том числе $\dot{\omega} - 1$, необходимо чтобы область значений кодирования немного выходила за рамки домена теории, что касается, в частности, область значений функции $diag(x, y)$. Таким образом, если мы позволяем системе иметь чисто внешнюю (Евклидову) морфологию, то мы устраняем саму возможность семантических петель самореференции с внешним контекстом. Это именно то о чём предупреждает Пуанкаре, когда мы отказываемся от части морфологической информации, мы открываем дорогу импредикативности: область значений функции $diag(x, y)$ выходит за рамки исходного домена, что и обеспечивает тот самый механизм изменения исходной классификации. Фактически Пуанкаре предвосхищает теорему Гёделя, задолго до того как она будет сформулирована. Благодаря Аристотелю мы знаем, что ключевое условие позволяющее недостатку информации привести к конфликту также выполнено, вследствие того что семантика обобщения существует в возможности, будь то квантор существования или всеобщности, представление натурального ряда как законченной сущности или конструктивного продолжение индукции на сколь угодно большое число. Иными словами, причина уязвимости классических теорий к эпистемологическим парадоксам кроется в том, что мы блокировали в этих теориях естественных механизм смены позиции референта, заставив этого референта созерцать исключительно внутреннюю перспективу - рефлексивные морфологии отбрасывающие часть структурной информации о бесконечных объектах. Но сохранили при этом синтаксическую возможность говорить о внешней семантике благодаря фрактальной воспроизводимости структуры. Мы не оставили никакой возможности этим теориям оперировать понятиями внешнего контура и восстановить предикативность бесконечных объектов с помощью утраченной дедуктивной компоненты морфологии, вследствие чего любая попытка внутренней интерпретации этих понятий ведёт к коллапсу. Теоремы Гёделя, Тарского, Клини и Тьюринга показывают нам две ключевые вещи: то, что самореференция неизбежна и то, что позиция референта - это неотъемлемая часть математики. Истребление самореференции было ошибкой сильно ослабившей наши теории. Самореференция не является источником парадоксов, она прекрасно функционирует в конечном предикативном мире не приводя к каким-либо противоречиям. Настоящий источник парадоксов - смешение позиций референта, попытка находится в двух различных точках одновременно, долгое время находился в тени. Теории лишённые естественного механизма смены позиции референта, стремящиеся лишь к внутреннему описанию, например, в терминах морфологии вполне упорядоченности (вижу цель, не вижу препятствий), обречены на устойчивую неполноту или условную непротиворечивость. Механизм устойчивости оказывается фундаментально неустраним на структурном уровне без апелляции к единой среде, допускающей Евклидову объективизацию любой сущности. Такая среда играет роль реального существования экземпляра классификации, например, в распечатанном виде, как в случае с парадоксом Барри. Нам требуется внешняя перспектива, жёстко фиксирующая превосходство целого над своими частями, что позволяет разглядеть скрытые манипуляции изменяющие домен теории. Индуктивно-дедуктивная морфология сама по себе помогает вскрыть подвох лишь для конкретных видов импредикативных утверждений, однако, поскольку она сама является структурной аннотацией и рефлексивна по своей природе, такая морфология сохраняет уязвимость к импредикативным высказываниям. Полностью избавится от таких импредикативных петель возможно лишь постулировав существование среды - внешней перспективы, позволяющей ввести сравнимый эталон бесконечности, превосходящий любую свою часть. Технически теоремы Гёделя и Тарского выводят противоречие из предположения о замкнутости домена, означающее некорректность предположения. В самом деле, мы предполагаем, что система замкнута и может интерпретировать все свои объекты, оставаясь при этом строго упорядоченной. Это позволяет нам оправдать использование метатеоретических концепций, таких как кванторы всеобщности и существования или даже понятие истины, которые требуют возможности интерпретировать весь домен. Но из этого выводится необходимость иметь возможность интерпретировать объекты (Гёделево утверждение или предикат истины) за пределами исходного домена, что противоречит исходному предположению. Такое толкование остаётся за пределами концептуальной рамки мейнстрима, несмотря на то, что Кантор вплотную приблизился к данному логическому ограничению, что привело его к структурному разделение независимых алгебр. Данное ограничение, по-видимому, обусловлено особенностями мышления разделяющего внутреннюю и внешнюю позицию референта. В действительности, мы уже используем эти перспективы в математике, иногда более осознанно. Например, в концепции мета-теорий, мата-уровень выступает в качестве внешней перспективы, а уровень самой теории - внутренней. А иногда менее осознанно, как в случае классов и множеств или двух мер Лебега и Банаха. Фактически меру Банаха мы можем ввести как меру Лебега просто отказавшись от сигма-аддитивности. А сигма-аддитивность - это как раз и есть условие внешнего взгляда, блокирующее рефлексивность внутренних отношений. Проблема лишь в том, что мы вводим эти концепции на очень высоком уровне, либо не осознаём этих перспектив в определённых понятиях, но без них в математике не обойтись. Так или иначе, данное наблюдение возвращает актуальность второй проблеме Гильберта о (не)противоречивости аксиом арифметики \cite{hilbert-problems}. 

Обсуждаемые в статье концепции могут оказаться полезными и в рамках математического анализа. Проблема разделения дискретного и непрерывного описания действительности представляет собой один из самых древних вопросов математики. Классическая идея континуума как множества точек не выдерживает столкновения с реальностью, это понимал ещё Аристотель \cite{aristotle-pointless-continuum}, который характеризует непрерывное как делимое до бесконечности \cite{aristotle-countinuum-divisibility}, тождественное своей части и упорядоченное \cite{aristotle-countinuum-similarity}. С современных позиций мы можем лишь констатировать, что все наши описания непрерывного - рациональные числа, вещественные числа, сюрреальные числа обладают указанными свойствами. Аннотация непрерывного рациональными числами обеспечивает порядок и позволяет делить непрерывное в рациональных пропорциях сколь угодно большое количество раз, при этом любые части будут подобны друг другу в отношении этих свойств делимости. Аналогичным образом вещественные числа по Кантору обеспечивают порядок и трансфинитную делимость в рамках второго числового класса, сохраняя при этом подобие. Сюрреальные числа проносят эти свойства через всю трансфинитную иерархию. Таким образом, континуум можно определить наиболее общим образом как контейнер упорядоченных отношений, обладающий внутренней рефлексией (подобием). В рамках дискретных теорий мы можем иметь различные аннтоции континуума, обладающие также различной морфологией, среди которых наиболее изученной является несчётная аннотация ZFCL и NBGL. ZFCS и NBGS дают счётные аннотации континуума, здесь различается непрерывность Дедекинда определённая через сечения множества рациональных чисел и непрерывность Вейерштрасса образованная числами с бесконечной позиционной записью. Индуктивно-дедуктивная морфология добавляет слой дедуктивных точек, позволяющих детально рассмотреть приближение к пределу. Сюрреальная аннотация континуума даёт чистую морфологию отношений. Различные аннотации континуума обогащают теорию мер, в частности, мера Лебега получает внешнюю поддержку от дедуктивных точек, что позволяет нам избежать ряда геометрических парадоксов теории множеств, таких, как парадокс Хаусдорфа \cite{hausdorff} и парадокс Банаха-Тарского \cite{banach-tarski}. Кроме того, ZFCS и NBGS делают вновь актуальным определение континуума Кантора, как замыкания предельного аппарата определённого уровня. Благодаря такой концепции мы преодолеваем ограничения рекурсивных подходов и становимся способны описывать континуум во всей полноте воспроизводства аддитивной структуры. Диверсификация моделей аннотации континуума позволяет нам лучше понимать смысл классической аксиомы непрерывности, связав её с конкретной морфологией бесконечности и возможно найти обобщённые формулировки. В частности, аксиома стеснения обеспечивающая замкнутость континуума в концепции Кантора, представляется одним из фундаментальнейших открытий о природе континуума, долгое время остававшимся в тени. Любопытно, что Пуанкаре также предвосхищает данное наблюдение, демонстрируя импредикативность описания точек пространства в классическом понимании континуума, что свидетельствует о связи наших представлений о непрерывном с морфологиями не обладающими полнотой информации. Кроме того, отмеченная диверсификация обеспечивает достоверность воспроизводимых в различных моделях результатов, что критически важно для того, чтобы избежать догматизации в математике.

Обратим внимание также на ещё один обнадёживающий результат: внешняя перспектива в равной степени с теоремами семантической диагонализации, позволяет противостоять теореме Пенроуза о невозможности алгоритмического искусственного интеллекта, в доказательстве которой используется конструктивная теорема Гёделя о неполноте \cite{penrose}. Вообще говоря, с математической точки зрения данная теорема заслуживает доверия, хотя её утверждение немного более скромно: невозможен алгоритм искусственного интеллекта, в том числе, алгоритм использующий нейронные сети. В этом плане Пенроуз абсолютно прав, как в утверждениях теоремы, так и в ответах на критику. Возможность преодоления ограничений теоремы Пенроуза следовало искать не в доказательстве, а в понятии алгоритма, и в том факте, что компьютеры могут выполнять программы не являющиеся алгоритмами, например, несколько машин Тьюринга работающие параллельно с одной лентой не описываются каким-либо алгоритмом, откуда следует, что изменив способ синхронизации вычислений в нейронных сетях мы получаем программу не ограниченную теоремой Пенроуза. Однако, судя по тому, что основательной критики не последовало, мало кто воспринял данную теорему со всей серьёзностью, что странно, учитывая внимание прикованное в настоящее время к проблеме искусственного интеллекта. В контексте возможности переключения позиций референта, формальный запрет на алгоритмический искусственный интеллект перестаёт существовать, а значит к нему могут привести и работы ведущиеся сейчас по всему миру, что оставляет место некоторым фантастическим надеждам. Тем не менее, нужно сказать и то, что в функционировании мозга квантовая составляющая, как отмечает Пенроуз, действительно может оказаться весьма существенной. Но в таком случае, состояние нейронов может оказаться некоторой непрерывной целостностью, точного описания которой не в состоянии дать наши дискретные модели, даже такие мощные как  сюрреальные числа, не говоря уже о конечных моделях реализованных в компьютерах. Тогда единственную надежду на создание искусственного интеллекта дают лишь квантовые компьютеры, в развитии которых мы, к сожалению, едва достигаем стадии квантовых калькуляторов.

История становления и развития теории множеств представляет большой интерес с точки зрения деталей и механизмов эволюции консенсуса научного сообщества. В рамках теории Томаса Куна о структуре научных революций, процесс консерватизации парадигм представляет собой сложный психологический переход от восприятия новой теории как маргинальной и провокационной (обладающей некоторым интуитивным изъяном, вызывающим вызов и сопротивление) к её догматическому воспроизводству как эталона, где изъян трансформируется в фундаментальную истину. Мы наблюдаем как обнаружение парадоксов теории множеств стимулирует взрывной рост интереса к данной области, сопровождаемый случайным закреплением утверждений имеющих неопределённый статус. Этот переход опирается на комплекс когнитивных искажений и социальных динамик, эволюционно укоренённых в человеческой психике, и может быть описан через призму модели стереотипного содержания как сдвиг от квадранта "холодной некомпетентности" (отрицание, критика) к "тёплой компетентности" (лояльность, восхищение), где бывший оппонент становится дороже старого союзника. Среди ключевых психологических механизмов можно назвать следующие. Эффект привыкания: интенсивное взаимодействие с парадоксальной теорией приводит к адаптации и эмоциональной лояльности. Эффект вложения и упрямство: исследователи, инвестируя усилия в разбор коллизий, развивают настойчивость, превращая неудачу в разрешении противоречий в эмоциональную привязанность, что усиливает лояльность через отказ зафиксировать убытки - потраченные ресурсы. Фундаментальная ошибка атрибуции и склонность к подтверждению: не обнаружив ошибки, исследователь атрибутирует её отсутствие истинности теории, игнорируя аномалии и фокусируясь на подтверждающих данных. Генерализация: переоценка универсальности теории на основе ограниченного опыта ускоряет её переход в догму. В итоге, многократное воспроизводство обеспечивает быстрый рост новых результатов, что вызывает срочную потребность в стандартизации положений, в процессе которого закрепляются сомнительные утверждения, закладывая семена будущих кризисов в результате накопленных аномалий. В сообществе формируются авторитеты, происходит переход к нормальной науке и социальное давление нормализует парадигму, подстраивая индивида под консенсус, что эволюционно обеспечивает стабильность групп. История теории множеств фиксирует саму динамику данного процесса. Гений Кантора обнаруживает конфликт позиций референта и находит решение в разделений кардинальной и порядковой алгебры. Уже в первое десятилетие после его работ появляется упрощённая интерпретация, сливающая обе перспективы в единую морфологию, которая соседствует с альтернативными взглядами. Упрощённая версия становится доминирующей и финальная вспышка в форме гениальной работы фон Неймана разделяющей совокупности на внутренние классы и в внешние множества уже не способна повлиять на конвейер научного тиражирования знаний, несмотря на свою яркость.

Период поздней нормальной науки по Куну можно описать как идеологический консерватизм. Структура, которая в начале давала стабильность и мощный рост, к концу превращается в репрессивный аппарат, охраняющий себя от изменений. Ключевые интерпретации теории становятся канонами, абсолютизируются, ключевые фигуры превращаются в идолов. Финальная фаза научной школы неотличима от жречества, она бюрократична и инертна, а энергия импульсов к изменению гасится о процедурную необходимость. Если Абрахам Френкель, один из архитекторов ZFC, называл диагональный аргумент Кантора рискованным и постоянно интересовался достижениями альтернативных концепций (подчёркивая философские уязвимости в основаниях множеств), а совместно с Бар-Хиллелем заключал: ''не существует, да и не предвидится, никакого единого и общепринятого способа перестройки математики, и в этом смысле кризис оснований всё ещё продолжается''. То современная математика эволюционировала в догматическую систему, подавляющую любую критику. Так в статье ''Воспоминания редактора о некоторых безнадежных статьях'' Вильфрид Ходжес упоминает десятки работ (в редакции одного единственного журнала), оспаривающих интерпретацию теоремы Кантора и её применимость к кардинальной алгебре. Но все они идут в разрез с консенсусом математического сообщества и цензурируются. Нечто подобное происходит на математических форумах и в Интернет сообществах, о чём пишет Вольфганг Мюккенхайм в своей книге ''Трансфинитность'': он приводит свыше 100 аргументов о проблемах теории, указывая на признаки цензуры идей, аналогичной идеологическому контролю - от блокировки публикаций до маргинализации критиков на онлайн-платформах. Такое положение дел говорит о социальной некомпетентности (или институциональной дисфункции) математического сообщества. Как социальный институт, оно призвано удовлетворять потребности общества в инструментах для описания наблюдаемой и воображаемой реальности, включая альтернативные математические парадигмы. Подавляющее большинство отвергаемых и маргинализируемых работ связано именно с проблемами теории множеств, практически никто не подвергает сомнению теорему Пифагора или основную теорему арифметики. Это напрямую свидетельствует о общественном запросе на более разнообразный математический аппарат для описания множественных сущностей - от философских интерпретаций бесконечности до прикладных моделей в физике или информатике. Однако данный запрос игнорируется институтом математики, который тем самым не выполняет прямую обязанность перед социумом. Проблема в том, что такое отношение блокирует авторам развивающим альтернативные концепции доступ к институциональному инструменту обмена опытом, блокируя развитие языка и понятийного фундамента. Вследствие чего, бюрократическому аппарату системы легко взаимодействовать с одиночками не способными последовательному и результативному изложению своих идей, продолжая период консервативной стагнации. Математика - это язык описания действительности, язык на котором функционирует наше собственное мышление, и тысячи работ говорят о том, что оно куда более широко и многогранно, чем пытается утверждать современный математический консенсус. Ограничивая эволюцию язык описания реальности, мы ограничиваем возможности собственного развития. 

Мейнстрим обвиняет альтернативных авторов в непонимании, но это всего лишь примитивная языковая игра, в действительности концептуальная основа современной теории множеств чрезвычайно проста, она интегрирует классические представления о потенциальной бесконечности в актуально бесконечный аппарат Кантора. Множество функций с конечным носителем - это именно потенциально бесконечный объект, принятие факта существования которого достаточно для того, чтобы осознать картину современенного математического консенсуса. Логика доказательства счётности счётного объединения счётных множеств вполне понятна, она апеллирует к той самой классической концепции неограниченного продолжения, многократное повторение которого сквозь столетия приводит к непоколебимой вере в корректность данного доказательства, вне зависимости от контекста. Но математика не строится на вере, она строится на строгом определении границ применимости тех или иных принципов. Математика требует обоснования распространения индуктивного аппарата потенциальной бесконечности на актуально бесконечный, сравнимый и упорядоченный объект $\omega$. Но такое обоснование отсутствует в современной интерпретации, то что мейнстрим называет пониманием представляет собой ничто иное как согласие со сложившимися практиками, наследование институциональной традиции, однако, - это понимание культа, а не понимание науки. Суть этой примитивной языковой игры сводится к обозначению отсутствия принадлежности к сообществу, где непонимание приобретает смысл унтерменша (недочеловека). Взаимно однозначное соответствие или биекция не имеет никакого сакрального смысла, это простые функциональные отношения, которые определены в рамках конкретной алгебраической структуры. Эти функции живут внутри алгебры и не позволяют выйти наружу за её пределы, в частности, чтобы описать свой собственный домен как законченный объект. Попытка заставить взаимно однозначное соответствие определить внешние контуры теории - это и есть та самая импредикативность, о которой настойчиво предупреждал Пуанкаре. Функция замкнутая в потенциально бесконечном кардинальном домене не способна увидеть актуально бесконечных объектов домена порядковых чисел, точно также как индуктивная функция не способна увидеть дедуктивных объектов. Кантор строго формулирует данное условие, для применения диагонального аргумента нам нужна система бесконечных координат, единая для всех объектов пространства. Мы сначала постулируем идентифицируемость бесконечности как объекта, и лишь затем применяем функциональную диагонализацию, которая требует, чтобы бесконечность была объектом. Практически никто впоследствии не воспроизводил это доказательство с такой степенью ясности и математической строгости. Поэтому понимание на которое претендует мейнстрим - есть ничто иное как декларация собственного статуса внутри идеологической системы.

Итак, самое сложное позади, мы разобрались с чрезвычайно запутанным клубком противоречий лежавшем в самом сердце теории множеств, из-за которого одна интерпретация стала доминирующей, плохо согласуясь при этом с альтернативными подходами. Причиной данного феномена становится некритически адаптированные к аксиоматикам рассуждения в рамках кардинальной алгебры, не воспроизводимые в рамках порядковой алгебры без дополнительных предположений. Явная формулировка таких предположений и противоположных разделяет теории на две группы, непротиворечиво объединяя практически все имеющиеся подходы в рамках единой концептуальной модели, возвращающей в математику утраченные категории позиции референта. Что позволяет универсальным образом интерпретировать все парадоксы самореференции и анализировать необходимость шагов по их нейтрализации в современных теориях. Сама по себе концепция связанная с позицией референта также носит гипотетический характер, находя лишь косвенные подтверждения в когнитивистике, лингвистике, истории науки и философии. Она имеет огромную объяснительную силу, позволяющую не только объединить различные альтернативные подходы, но и аккуратно позиционировать сильные классические теоремы, такие как теоремы Гёделя и Тарского. Интерпретация мейнстрима не обладает подобным потенциалом, но силы современных теорий не хватает, чтобы явно продемонстрировать её противоречивость, что связано с фундаментальной относительностью понятия бесконечности, которая обеспечивает импредикативную перестройку домена позволяющую избегать явных противоречий. Таким образом, объяснительная сила альтернативной картины, позволяет если не усомнится в интерпретации мейнстрима, то рассматривать альтернативную концептуальную рамку как равноправную интерпретацию, заслуживающую столь же пристального внимания. ZFCL и NBGL представляют собой классическую интерпретацию теории множеств смешивающую перспективы. Здесь $\omega = 2^{\omega}$ не приводит к немедленному противоречию благодаря теореме Лёвенгейма-Скулема, которая позволяет считать, что аксиоматические бесконечные множества существуют на одном кардинальном уровне, в то время как, индуктивные множества формируют иерархию предыдущих. Такое отсутствие связи между индуктивным и аксиоматическим существованием множеств находит своё выражение в континуум-гипотезе. Однако, данный подход блокирует естественный механизм переключения перспектив, в результате чего альтернативные теории оказываются плохо совместимыми с ним, а простые интуитивные модели требуют сложного аппарата нестандартного анализа. В свою очередь теории ZFCS и NBGS обеспечивают переключение контекста внутреннего и внешнего описания, что приводит к согласованности кардинальной, порядковой и индексной морфологий. Здесь индуктивные множества однозначно связываются с аксиоматическими, и множество $2^{\omega}$ становится полноценной степенью множества $\omega$. Более того, в рамках данного подхода все внутренние морфологии могут быть вписаны в единое универсальное Евклидово пространство, обеспечивающее предикативность всех отношений вписанных структур, что мощным внутренним механизмом верификации теории. 

Авторы настоящей серии статей надеются, что их скромный вклад окажется способным принести толику пользы на пути математики к выходу из этого кризиса, сдерживающего её развитие в течении длительного периода. После обнаружения неоднозначности в диагональном методе, перед ними стояла неимоверно сложная задача, найти механизмы реализующие альтернативные перспективы на одни и те же объекты возникающие в нашем сознании, что возлагало на них огромную ответственность, не оставляя им никакого права посеять сомнения, способные привести к тому, что мучительная операция их возвращения в математику затянется. Кроме того, хоть авторы и видят пути к более строгой формализации понятий теории множеств, было бы ошибкой предъявлять их научному сообществу, несмотря на то, что это бы существенно упростило ряд рассуждений. Гораздо увлекательнее и честнее представить перед исследователями открытую проблему, чем предложить готовое решение, ведь таких решений может быть очень много, и упускать это разнообразие означало бы значительно ограничить наше понимание. Хуже того, отсутствие разнообразных описаний обнаруживающих когерентность между собой несёт в себе опасность новой догматизации и идеологизации. Анализ парадигм математики приведённый в данной работе ярко демонстрирует факт того, что различные подходы не противоречат, а дополняют друг друга. Формализм несмотря на всю свою силу, оказывается уязвим к реификации и бюрократизации процессов, которая ведёт к консерватизму и догматизации. Формализм не имеет механизмов выработки понятийного аппарата и языка, его основное предназначение - это инженерное исследование данных зафиксированных систем. Формализм - это огораживание территорий - песочниц, где безопасно вести раскопки. Но это не снимает с нас задачи исследования мира в целом, включая всего того, что находится за пределами данных песочниц. Интуиционизм наоборот концентрируется на процессе концептуализации  мыслеобразов и формировании языка и понятийного аппарата для их описания. Этот процесс фундаментально неформализуем, в нём нет места строгим правилам наделения смыслом внутренних объектов мышления. Логицизм обеспечивает механизмы связи между различными подходами, не случайно ограничения логицизма, которые мы отказывались формулировать в явном виде, воспроизводятся во всех теориях. Однако связи в рамках логицизма оказываются невероятно сложными и громоздкими. Гораздо эффективнее оказывается выделение инвариантных свойств различных композиций в рамках структурализма и перенос их с одних построений на другие. Проблема лишь в том, что механизм переноса не предусматривает верификации утверждений, вследствие чего, например, смешение перспектив преодолевает границы теории множеств и обнаруживает себя, например, в интуиционисткой теории типов Мартина-Лёфа. В то время как, в идеале всё должно происходить с точностью до наоборот, ведь наука - это моделирование действительности, основанное на изоморфизме различных подходов, языков и понятийных систем. Выявление такой когерентности различных описаний и является критерием научного знания, взаимно верифицирующим альтернативные подходы к исследованию действительности. Некритический перенос концепций бесконечности из одной теории в другую, вместо естественного возникновения собственного представления, существенно тормозит развитие. Построенный в рамках интуиционистского подхода конструктивный язык математического описания, является минимальным базисом верифицирующим все остальные направления. Данному описанию порой не достаёт силы, но оно проводит чёткую границу предикативной и импредикативной реальности. Поэтому в современной математике владение всеми этими подходами рассматривается авторами как безусловная необходимость. Руководствуясь этими соображениями авторы приложили все усилия, которые только могли, и просят у читателей прощения, если где-то переборщили с подробностями, а где-то наоборот аргументы кажутся недостаточно обоснованными. Впоследствии мы увидим, как избавление от оков разобранных заблуждений позволяет с невероятной лёгкостью окончательно разрешить практически все известные парадоксы теории множеств, не опасаясь реванша, но что ещё более важно насколько обширные и интересные направления для исследований открываются перед нами, когда к математике возвращается утерянная стройность и непротиворечивость. И пусть авторам настоящей работы, если представится такой случай, доведётся проиллюстрировать лишь ничтожную часть всего многообразия захватывающих творений математической реальности, которые готовит новая эпоха, гораздо важнее то, чтобы выход из тупика, найденный ими, как можно скорее стал достоянием всей математической общественности, поскольку, по их глубокому убеждению, он способен подтолкнуть продвижение не только в самой математике но и в физике, в частности, квантовой механике, провозгласившей элементарными объектами описания вселенной функции, инструментарий в пространстве которых, судя по всему, был существенно ограничен классическими представлениями о мощности множеств. Перестройка математики неизбежна, однако всегда лучше иметь спектр альтернативных концепций предложенных разными авторами, чтобы когда настанет время для фазы синтеза, мы могли надеяться, что новое здание не даст столь прискорбных трещин в фундаменте.

\printbibliography

 \end{document}